\documentclass[a4paper, 11pt, oneside]{book}

\usepackage{fontspec}
\usepackage[margin=3.25cm]{geometry}
\usepackage[italian]{babel}
\usepackage{amsmath}
\usepackage{amssymb}
\usepackage{amsthm}
\usepackage{mathtools}
\usepackage{tikz}
\usepackage{XCharter}
\usepackage[xcharter]{newtxmath}
\usepackage[hidelinks]{hyperref}
\usepackage[most]{tcolorbox}
\usepackage{float}
\usepackage{graphicx}
\usepackage{tabularx}
\usepackage{booktabs}
\usepackage{stmaryrd}
\usepackage{centernot}
\usepackage{needspace}
\usepackage{microtype}
\usepackage{enumitem}
\usepackage{fancyhdr}
\usepackage{titlesec}
\usepackage{titletoc}
\usepackage{caption}

\hypersetup{
    pdftitle={Fondamenti dell'Informatica},
    pdfauthor={Michele Castenedoli}
}

\setcounter{tocdepth}{1}
\setlength{\skip\footins}{15pt plus 4pt minus 2pt}
\usetikzlibrary{positioning, fit, arrows.meta, backgrounds, matrix,
decorations.pathreplacing, automata, decorations.pathmorphing}

% --- IMPOSTAZIONI GLOBALI PER I GRAFI TIKZ ---
% Assicura che tutti i nodi (stati degli automi e punti negli insiemi) abbiano esattamente la stessa dimensione
\tikzset{
    every state/.style={draw, thick, minimum size=2.5em, inner sep=2pt},
    punto/.style={circle, fill, minimum size=3pt, inner sep=0pt},
    stato/.style={circle, fill, minimum size=3pt, inner sep=0pt},
    % Stile comune a tutti gli automi: frecce, spessore e distanza tra gli stati
    automa/.style={->, >={Stealth}, auto, thick, initial text={},
        node distance=2.5cm},
    % Sotto-automa di una costruzione: riquadro tratteggiato attorno ai suoi stati
    sottoautoma/.style={draw, dashed, thin, -, rounded corners=8pt,
        inner xsep=10pt, inner ysep=8pt},
    % Parte interna di un sotto-automa, di cui si mostrano solo gli estremi
    interno/.style={-, thin, decorate, decoration={snake, amplitude=0.4mm,
        segment length=2mm, pre length=1mm, post length=1mm}}
}

\defaultfontfeatures{}
\setmonofont{JetBrainsMono Nerd Font}[
    Scale=MatchLowercase, Contextuals=Alternate
]

% Elenchi compatti, con punto grigio
\setlist[itemize]{nosep, leftmargin=2em, topsep=3pt, partopsep=0pt,
    label=\textcolor{black!50}{\small\textbullet}}
\setlist[enumerate]{nosep, leftmargin=2em, topsep=3pt, partopsep=0pt}

% Evita righe isolate a inizio o fine pagina
\widowpenalty=10000
\clubpenalty=10000

% Colori delle figure: solo grigi, per la stampa in bianco e nero
\colorlet{evidenza}{black!70}
\colorlet{tenue}{black!8}

% Capitoli: numero grande in grigio, barra verticale e titolo in maiuscoletto
\newcommand{\numerocapitolo}{\fontsize{60}{60}\selectfont}
\titleformat{\chapter}[block]
    {\normalfont\LARGE}
    {{\color{black!45}\numerocapitolo\thechapter}\hspace{12pt}\vline}
    {12pt}
    {\scshape}
\titleformat{name=\chapter, numberless}[block]
    {\normalfont\LARGE\scshape}{}{0pt}{}
\titlespacing*{\chapter}{0pt}{50pt}{45pt}
\titleformat{\section}
    {\normalfont\Large\scshape\raggedright}{\thesection}{1em}{}
\titleformat{\subsection}
    {\normalfont\large\itshape}{\thesubsection}{1em}{}

% Numerazione sempre preceduta dal capitolo, anche nel Capitolo 0
\renewcommand{\thefigure}{\thechapter.\arabic{figure}}
\renewcommand{\thetable}{\thechapter.\arabic{table}}

% Indice: capitoli in maiuscoletto con il numero di pagina in grassetto
\titlecontents{chapter}[1.5em]{\addvspace{1em}}
    {\contentslabel{1.5em}\scshape}{\hspace*{-1.5em}\scshape}
    {\hfill\bfseries\contentspage}

% Intestazione senza riga: sezione corrente a sinistra, pagina a destra
\renewcommand{\chaptermark}[1]{\markboth{\thechapter\enspace #1}{\thechapter\enspace #1}}
\renewcommand{\sectionmark}[1]{\markright{\thesection\enspace #1}}
\pagestyle{fancy}
\fancyhf{}
\fancyhead[L]{\scshape\rightmark}
\fancyhead[R]{\thepage}
\renewcommand{\headrulewidth}{0pt}
\setlength{\headheight}{14pt}
\fancypagestyle{plain}{\fancyhf{}\renewcommand{\headrulewidth}{0pt}}

% Didascalie: etichetta in maiuscoletto, testo in corsivo
\captionsetup{font=small, labelfont=sc, textfont=it, labelsep=quad, skip=8pt}

% Enunciati: intestazione in maiuscoletto e testo in corsivo, senza riquadri
\newtheoremstyle{formale}
    {\topsep}% spazio sopra
    {\topsep}% spazio sotto
    {\itshape}% corpo
    {}% rientro
    {\scshape}% intestazione
    {.}% punto dopo l'intestazione
    {0.5em}% spazio dopo l'intestazione
    {\thmname{#1}\thmnumber{ #2}\thmnote{ {\normalfont(#3)}}}
\theoremstyle{formale}
\newtheorem{theorem}{Teorema}
\newtheorem{lemma}{Lemma}
\newtheorem{corollary}{Corollario}
\newtheorem*{definition}{Definizione}

% Chiusura rispetto alle epsilon-transizioni: trattino e testo dritti, non un meno
\DeclareMathOperator{\eclosure}{\varepsilon\text{-}closure}

% Esempi: intestazione in corsivo, chiusi da un rombo
\newenvironment{esempio}
    {\begin{proof}[\textit{Esempio}]\renewcommand{\qedsymbol}{$\diamond$}}
    {\end{proof}}

% Va a pagina nuova solo se nella pagina corrente resta meno spazio di #1. A differenza
% di \needspace non aggiunge penalit\`a, che anticiperebbero il cambio pagina prima di
% ogni enunciato anche quando c'\`e spazio, lasciando pagine mezze vuote. Subito dopo un
% titolo non fa nulla, per non separare il titolo da ci\`o che segue
\makeatletter
\newcommand{\spaziominimo}[1]{%
    \par
    \if@nobreak\else
        \ifdim\pagegoal<\maxdimen
            \ifdim\dimexpr\pagegoal-\pagetotal\relax<#1\relax
                \newpage
            \fi
        \fi
    \fi}
\makeatother

% Evita che l'intestazione di un enunciato o di una sezione resti isolata in fondo
% alla pagina
\AddToHook{env/definition/before}{\spaziominimo{5\baselineskip}}
\AddToHook{env/theorem/before}{\spaziominimo{5\baselineskip}}
\AddToHook{env/esempio/before}{\spaziominimo{4\baselineskip}}
\AddToHook{cmd/section/before}{\spaziominimo{10\baselineskip}}
\AddToHook{cmd/subsection/before}{\spaziominimo{7\baselineskip}}

\definecolor{codeborder}{HTML}{707070}
\definecolor{codeback}{HTML}{F0F0F0}
\definecolor{codekeyword}{HTML}{1F3FA8}
\definecolor{codecomment}{HTML}{6A737D}
\definecolor{codestring}{HTML}{A31515}

\lstdefinelanguage{pseudo}{
    morekeywords={
        function, end, returns, return, if, then, else,
        while, do, for, to, each, in, loop, repeat, until,
        persistent, static, local, and, or, not,
        true, false, null, failure
    },
    sensitive=true,
    alsoletter={-},
    morecomment=[l]{//},
    morestring=[b]"
}

\newcounter{escount}[section]
\newcommand{\esercizio}{
    \spaziominimo{8\baselineskip}% Evita il titolo isolato a fondo pagina
    \stepcounter{escount}% Incrementa il contatore di 1
    \subsection*{Esercizio \arabic{escount}}% Genera il titolo non numerato nativamente
}

% Testo dell'esercizio, evidenziato rispetto alla soluzione
\newtcolorbox{enunciato}{
    enhanced,
    breakable,
    sharp corners,
    boxrule=0pt,
    leftrule=2pt,
    colframe=codekeyword,
    colback=codekeyword!5,
    before skip=6pt,
    after skip=10pt,
    left=8pt,
    right=8pt
}


% Titolo di una (parte di) dimostrazione, su una riga a s\'e
\newcommand{\dimostrazione}[1]{
    \par\needspace{4\baselineskip}\medskip
    \noindent\textbf{Dimostrazione #1}\par
}


\newtcblisting{codice}[1][Codice]{
    breakable,
    enhanced,
    sharp corners,
    boxrule=0.6pt,
    colframe=codeborder,
    colback=codeback,
    title=#1,
    fonttitle=\bfseries,
    coltitle=codeborder,
    attach title to upper={\par\smallskip},
    before skip=12pt,
    after skip=12pt,
    boxsep=4pt,
    left=4pt,
    right=4pt,
    grow to left by=1cm,
    grow to right by=1cm,
    listing only,
    listing options={
        language=pseudo,
        basicstyle=\ttfamily\footnotesize,
        keywordstyle=\color{codekeyword}\bfseries,
        commentstyle=\color{codecomment}\itshape,
        stringstyle=\color{codestring},
        breaklines=true,
        breakatwhitespace=true,
        breakautoindent=true,
        breakindent=2em,
        postbreak=\mbox{\textcolor{codecomment}{$\hookrightarrow$}\space},
        tabsize=4,
        columns=fullflexible,
        keepspaces=true,
        literate={<-}{{←}}1 {->}{{→}}1 {<=}{{≤}}1 {>=}{{≥}}1 {!=}{{≠}}1
    }
}

\title{Fondamenti dell'Informatica}
\author{Michele Castenedoli}
\date{13 ottobre 2026}

\begin{document}
    \begin{titlepage}
        \centering
        \vspace*{0.3\textheight}
        {\LARGE\scshape Fondamenti dell'Informatica\par}
        \vspace{1em}
        {\large Appunti di Michele Castenedoli\par}
        \vfill
        {13 ottobre 2026\par}
    \end{titlepage}

    \frontmatter
    \tableofcontents

    \mainmatter
    \setcounter{chapter}{-1}% Il primo capitolo \`e il Capitolo 0
    \chapter{Prerequisiti matematici}
    \section{Basi di logica}
    \subsection{Linguaggio del primo ordine}
    Un \emph{linguaggio del primo ordine} \`e formato da quattro tipi di simboli:
    \begin{enumerate}
        \item \textsc{Simboli di relazione}, indicati con le metavariabili $p, q, r$;
        \item \textsc{Simboli di funzione}, indicati con le metavariabili $f, g, h$;
        \item \textsc{Simboli di costante}, indicati con le metavariabili $c, d$;
        \item \textsc{Simboli logici}:
            \begin{itemize}
                \item le parentesi $($ e $)$ e la virgola;
                \item un insieme numerabile di \emph{variabili}, indicate con le
                    metavariabili $x, y, v$;
                \item i \emph{connettivi logici} e i \emph{quantificatori}
                    $(\land, \lor, \lnot, \implies, \iff, \forall, \exists)$;
                \item il simbolo $\bot$, che indica il falso.
            \end{itemize}
    \end{enumerate}
    I simboli di relazione e di funzione hanno ciascuno la propria \emph{arit\`a}, ovvero
    il numero di argomenti.

    A partire da questi simboli si costruiscono, in ordine, i termini, le formule atomiche
    e le formule.

    \begin{definition}[Termine]
        I \emph{termini} sono definiti per induzione:
        \begin{enumerate}
            \item ogni variabile \`e un termine;
            \item ogni costante \`e un termine;
            \item se $f$ \`e un simbolo di funzione $n$-ario e $t_1, \dots, t_n$ sono
                termini, allora $f(t_1, \dots, t_n)$ \`e un termine.
        \end{enumerate}
    \end{definition}

    \begin{definition}[Formula atomica]
        Le \emph{formule atomiche} sono:
        \begin{enumerate}
            \item $p(t_1, \dots, t_n)$, se $p$ \`e un simbolo di relazione $n$-ario e
                $t_1, \dots, t_n$ sono termini;
            \item il simbolo $\bot$.
        \end{enumerate}
    \end{definition}

    \begin{definition}[Formula]
        Le \emph{formule} sono definite per induzione:
        \begin{enumerate}
            \item ogni formula atomica \`e una formula;
            \item se $\phi$ \`e una formula, allora $\neg \phi$ \`e una formula;
            \item se $\phi$ e $\psi$ sono formule, allora $\phi \circ \psi$ \`e una formula,
                con $\circ \in \{\land, \lor, \implies, \iff\}$;
            \item se $\phi$ \`e una formula e $x$ \`e una variabile, allora
                $\forall x \, (\phi)$ e $\exists x \, (\phi)$ sono formule.
        \end{enumerate}
    \end{definition}

    \subsection{Equivalenze}
    Si scrive $\phi \equiv \psi$ quando le formule $\phi$ e $\psi$ sono
    \emph{equivalenti}, ovvero assumono sempre lo stesso valore di verit\`a.

    Due negazioni consecutive si annullano (\emph{doppia negazione}):
    \begin{equation*}
        \neg (\neg \phi) \equiv \phi.
    \end{equation*}
    Le \emph{leggi di De Morgan} descrivono la negazione di una congiunzione e di una
    disgiunzione:
    \begin{equation*}
        \neg (\phi \land \psi) \equiv (\neg \phi) \lor (\neg \psi),
        \qquad\qquad
        \neg (\phi \lor \psi) \equiv (\neg \phi) \land (\neg \psi).
    \end{equation*}
    Un'\emph{implicazione} pu\`o essere riscritta come disgiunzione, oppure come la sua
    \emph{contronominale} (l'ultima forma):
    \begin{equation*}
        (\phi \implies \psi)
        \equiv (\neg \phi) \lor \psi
        \equiv \bigl((\neg \psi) \implies (\neg \phi)\bigr).
    \end{equation*}
    Infine, la \emph{negazione di un quantificatore} si ottiene scambiandolo con l'altro e
    negando la formula:
    \begin{equation*}
        \neg \bigl(\exists x \, (\phi)\bigr) \equiv \forall x \, (\neg \phi),
        \qquad\qquad
        \neg \bigl(\forall x \, (\phi)\bigr) \equiv \exists x \, (\neg \phi).
    \end{equation*}

    \subsection{Abbreviazioni}
    Per i quantificatori ristretti agli elementi di un insieme $S$ si adottano le seguenti
    abbreviazioni:
    \begin{align*}
        (\exists x \in S.\; \phi) &\equiv \exists x \, (x \in S \land \phi), \\
        (\forall x \in S.\; \phi) &\equiv \forall x \, (x \in S \implies \phi),
    \end{align*}
    dove il punto si legge \emph{tale che}.

    \section{Insiemistica}
    \subsection{Definizione e rappresentazione}
    Un \emph{insieme} \`e una collezione di elementi senza ripetizioni. Esso pu\`o essere
    rappresentato in due modi. La rappresentazione \emph{intensionale} descrive gli
    elementi attraverso una loro caratteristica: ad esempio
    \begin{equation*}
        \{x \mid p(x)\}
    \end{equation*}
    \`e l'insieme di tutti gli elementi $x$ che rendono vero il predicato $p$. La
    rappresentazione \emph{estensionale}, invece, elenca tutti gli elementi
    dell'insieme, come in $\{1, 2, 3\}$.

    Un'altra definizione fondamentale per il resto del corso \`e quella di sottoinsieme.

    \begin{definition}[Sottoinsieme]
        $A \subseteq B$, ovvero $A$ \`e un \emph{sottoinsieme} di $B$, se ogni elemento di
        $A$ appartiene anche a $B$. $A$ \`e invece un \emph{sottoinsieme proprio} di $B$,
        $A \subset B$, se in pi\`u $B$ contiene almeno un elemento che non appartiene ad $A$:
        \begin{gather*}
            A \subseteq B \iff \forall x \, (x \in A \implies x \in B), \\
            A \subset B \iff A \subseteq B \land \exists x \in B.\; x \notin A.
        \end{gather*}
    \end{definition}

    \subsection{Operazioni tra insiemi}
    \label{sec:operazioni-insiemi}
    Le operazioni classiche tra insiemi sono:
    \begin{itemize}
        \item \textsc{Unione}: $A \cup B = \{x \mid x \in A \lor x \in B\}$;
        \item \textsc{Intersezione}: $A \cap B = \{x \mid x \in A \land x \in B\}$;
        \item \textsc{Differenza}: $A \setminus B = \{x \mid x \in A \land x \notin B\}$;
        \item \textsc{Complemento}: $\overline{A} = \{x \in U \mid x \notin A\}$, dove $U$
            \`e l'\emph{insieme universo}, che contiene tutti gli elementi considerati.
    \end{itemize}
    Per ogni $x \in U$ vale quindi $x \in A \iff x \notin \overline{A}$ e, analogamente,
    $x \in \overline{A} \iff x \notin A$.

    \subsection{Relazioni}
    Per definire le relazioni \`e necessario il \emph{prodotto cartesiano}. Dati gli insiemi
    $A_1, A_2, \dots, A_n$, il loro prodotto cartesiano \`e definito come mostrato nella
    Figura~\ref{fig:prodotto}.

    \begin{figure}[H]
        \centering
        \begin{tikzpicture}[
            every node/.style={inner sep=0pt, outer sep=0pt},
            freccia/.style={-{Stealth}, evidenza, thick, shorten <=5pt, shorten >=5pt}
        ]
            \node (A1) {$A_1$};
            \node[base right=2.5pt of A1] (p1) {$\times$};
            \node[base right=2.5pt of p1] (A2) {$A_2$};
            \node[base right=2.5pt of A2] (p2) {$\times \dots \times$};
            \node[base right=2.5pt of p2] (An) {$A_n$};
            \node[base right=3pt of An] (uguale) {$= \{\, \langle$};
            \node[base right=0.5pt of uguale] (x1) {$x_1$};
            \node[base right=0pt of x1] (v1) {$,$};
            \node[base right=2pt of v1] (x2) {$x_2$};
            \node[base right=0pt of x2] (v2) {$, \dots,$};
            \node[base right=2pt of v2] (xn) {$x_n$};
            \node[base right=0.5pt of xn] {$\rangle \mid \forall i.\; x_i \in A_i \,\}$};

            \draw[freccia] (A1.south) to[bend right=35] (x1.south);
            \draw[freccia] (A2.north) to[bend left=30] (x2.north);
            \draw[freccia] (An.south) to[bend right=25] (xn.south);
        \end{tikzpicture}
        \caption{Prodotto cartesiano: la componente $x_i$ di ogni tupla appartiene
            all'insieme $A_i$.}
        \label{fig:prodotto}
    \end{figure}

    Gli elementi del prodotto cartesiano sono le \emph{tuple}
    $\langle x_1, x_2, \dots, x_n \rangle$, in cui la componente $i$-esima $x_i$ appartiene
    all'insieme $A_i$. A differenza degli insiemi, nelle tuple l'ordine \`e rilevante e gli
    elementi possono ripetersi: $\langle 1, 2 \rangle \neq \langle 2, 1 \rangle$, mentre
    $\{1, 2\} = \{2, 1\}$.

    \begin{definition}[Relazione binaria]
        Dati due insiemi $A$ e $B$, una \emph{relazione binaria} tra $A$ e $B$ \`e un
        sottoinsieme del loro prodotto cartesiano:
        \begin{equation*}
            R \subseteq A \times B.
        \end{equation*}
        Se $x_1 \in A$ e $x_2 \in B$ sono in relazione si scrive
        \begin{equation*}
            x_1 \mathrel{R} x_2 \qquad \text{oppure} \qquad (x_1, x_2) \in R.
        \end{equation*}
    \end{definition}
    Le coppie, ovvero le tuple di due elementi, si indicano anche con le parentesi tonde.

    \subsection{Propriet\`a delle relazioni}
    Le propriet\`a seguenti riguardano le relazioni su un insieme $S$, ovvero
    $R \subseteq S \times S$.

    \begin{definition}[Relazione riflessiva]
        Una relazione $R$ \`e \emph{riflessiva} se ogni elemento \`e in relazione con se
        stesso:
        \begin{equation*}
            \forall x \in S.\; x \mathrel{R} x.
        \end{equation*}
    \end{definition}

    \begin{esempio}
        La relazione $\subseteq$ tra insiemi \`e riflessiva, poich\'e ogni insieme \`e
        sottoinsieme di se stesso:
        \begin{equation*}
            \{1, 2, 5, 8\} \subseteq \{1, 2, 5, 8\}.
        \end{equation*}
        La riflessivit\`a deriva dal fatto che il simbolo $\subseteq$ ammette l'uguaglianza;
        il sottoinsieme proprio $\subset$, invece, non \`e riflessivo, poich\'e nessun insieme
        \`e sottoinsieme proprio di se stesso.
    \end{esempio}

    \begin{definition}[Relazione simmetrica]
        Una relazione $R$ \`e \emph{simmetrica} se la posizione degli elementi nella
        relazione non \`e rilevante:
        \begin{equation*}
            \forall x, y \in S.\; x \mathrel{R} y \implies y \mathrel{R} x.
        \end{equation*}
    \end{definition}

    \begin{esempio}
        L'uguaglianza tra insiemi \`e simmetrica: se $A = B$, allora anche $B = A$.
        \begin{equation*}
            \underbrace{\{2, 5, 10\}}_{A} = \underbrace{\{2, 5, 10\}}_{B}
        \end{equation*}
        La relazione $\subseteq$, invece, non \`e simmetrica: $\{1\} \subseteq \{1, 2\}$, ma
        $\{1, 2\} \not\subseteq \{1\}$.
    \end{esempio}

    \begin{definition}[Relazione transitiva]
        Una relazione $R$ \`e \emph{transitiva} se
        \begin{equation*}
            \forall x, y, z \in S.\; x \mathrel{R} y \land y \mathrel{R} z \implies
            x \mathrel{R} z.
        \end{equation*}
    \end{definition}
    In altri termini, se $x$ \`e in relazione con $y$ e $y$ \`e in relazione con $z$, allora
    anche $x$ \`e in relazione con $z$, come illustrato nella Figura~\ref{fig:transitiva}.

    \begin{figure}[H]
        \centering
        \begin{tikzpicture}[
            freccia/.style={-{Stealth}, thick, shorten <=2pt, shorten >=2pt}
        ]
            \node[punto, label=below:$x$] (x) at (0, 0) {};
            \node[punto, label=below:$y$] (y) at (2, 0) {};
            \node[punto, label=below:$z$] (z) at (4, 0) {};
            \draw[freccia] (x) -- node[below] {$R$} (y);
            \draw[freccia] (y) -- node[below] {$R$} (z);
            \draw[freccia, evidenza, dashed] (x) to[bend left=35]
                node[above] {$x \mathrel{R} z$} (z);
        \end{tikzpicture}
        \caption{Relazione transitiva: da $x \mathrel{R} y$ e $y \mathrel{R} z$ segue
            $x \mathrel{R} z$.}
        \label{fig:transitiva}
    \end{figure}

    \begin{definition}[Chiusura transitiva]
        La \emph{chiusura transitiva} di una relazione $R \subseteq S \times S$ \`e la
        relazione $R^*$ definita come il pi\`u piccolo insieme $R^* \supseteq R$ tale che
        $R^*$ sia transitiva: essa si ottiene aggiungendo a $R$ le sole coppie necessarie a
        renderla transitiva.
    \end{definition}

    \begin{esempio}
        Se $R = \{(1, 2), (2, 3)\}$, allora
        \begin{equation*}
            R^* = \{(1, 2), (2, 3), (1, 3)\},
        \end{equation*}
        dove la coppia $(1, 3)$ \`e necessaria poich\'e $1 \mathrel{R} 2$ e
        $2 \mathrel{R} 3$.
    \end{esempio}

    \begin{definition}[Relazione antisimmetrica]
        Una relazione $R$ \`e \emph{antisimmetrica} se
        \begin{equation*}
            \forall x, y \in S.\; x \mathrel{R} y \land y \mathrel{R} x \implies x = y.
        \end{equation*}
    \end{definition}
    In altri termini, elementi diversi non possono essere scambiati nella relazione: se
    $x \neq y$ e $x \mathrel{R} y$, allora non pu\`o valere $y \mathrel{R} x$.

    \begin{esempio}
        La relazione $\subseteq$ \`e antisimmetrica, poich\'e $A \subseteq B$ e
        $B \subseteq A$ implicano insieme $A = B$.
    \end{esempio}

    \subsection{Tipologie di relazione}
    \begin{definition}[Tipi di relazione]
        I tipi classici di relazione $R \subseteq S \times S$ sono:
        \begin{itemize}
            \item \textsc{Equivalenza}: $R$ \`e riflessiva, simmetrica e transitiva.
            \item \textsc{Pre-ordine}: $R$ \`e riflessiva e transitiva. Si tratta di un
                ordinamento molto debole, poich\'e manca l'antisimmetria: elementi diversi
                possono essere in relazione in entrambi i versi.
            \item \textsc{Ordine parziale}: $R$ \`e un pre-ordine antisimmetrico.
                \emph{Parziale} significa che non tutti gli elementi sono in relazione tra
                loro.
            \item \textsc{Ordine totale}: $R$ \`e un ordine parziale in cui due elementi
                qualsiasi sono sempre in relazione:
                \begin{equation*}
                    \forall x, y \in S.\; x \mathrel{R} y \lor y \mathrel{R} x.
                \end{equation*}
        \end{itemize}
    \end{definition}

    \begin{esempio}
        La relazione $=$ \`e di equivalenza, $\subseteq$ \`e un ordine parziale ($\{1\}$ e
        $\{2\}$ non sono in relazione in nessuno dei due versi) e $\leq$ sui numeri naturali
        \`e un ordine totale.
    \end{esempio}

    \begin{definition}[Partizione]
        Una \emph{partizione} di un insieme $S$ \`e una famiglia di sottoinsiemi
        $\{S_1, S_2, \dots, S_m\}$ di $S$, a due a due disgiunti, la cui unione \`e $S$:
        \begin{equation*}
            \bigcup_{i=1}^{m} S_i = S,
            \qquad\qquad
            \forall i \neq j.\; S_i \cap S_j = \emptyset.
        \end{equation*}
    \end{definition}
    Ogni elemento di $S$ appartiene quindi a uno e un solo sottoinsieme della partizione,
    come mostra la Figura~\ref{fig:partizione}.

    \begin{figure}[H]
        \centering
        \begin{tikzpicture}[thick]
            % Punti sul bordo dell'ellisse, dati come angoli polari
            \coordinate (A) at (115:3cm and 1.5cm);
            \coordinate (B) at (235:3cm and 1.5cm);
            \coordinate (C) at (310:3cm and 1.5cm);
            \coordinate (E) at (265:3cm and 1.5cm);
            \coordinate (G) at (15:3cm and 1.5cm);

            % Contorno dell'insieme S
            \draw (0, 0) ellipse (3cm and 1.5cm);
            \node at (3.5, -1) {$S$};

            % Punti interni sul segmento AC, da cui partono i tagli secondari
            \path (A) -- (C) coordinate[pos=0.35] (F) coordinate[pos=0.7] (D);

            % Tagli che separano i sottoinsiemi
            \draw (A) -- (B);
            \draw (A) -- (C);
            \draw (D) -- (E);
            \draw (F) -- (G);

            \node at (-2.0, 0) {$S_1$};
            \node at (0.2, 0.9) {$S_2$};
            \node at (-0.8, -0.6) {$S_3$};
            \node at (0.8, -1.0) {$S_4$};
            \node at (1.8, 0.0) {$S_5$};
        \end{tikzpicture}
        \caption{Una partizione dell'insieme $S$ in cinque sottoinsiemi disgiunti.}
        \label{fig:partizione}
    \end{figure}

    Partizioni e relazioni di equivalenza sono strettamente legate: una relazione di
    equivalenza $R$ su $S$ e una partizione $\{S_1, \dots, S_m\}$ di $S$ si
    corrispondono quando
    \begin{itemize}
        \item due elementi dello stesso sottoinsieme sono sempre in relazione: se
            $a, b \in S_i$, allora $a \mathrel{R} b$;
        \item due elementi di sottoinsiemi diversi non sono mai in relazione: se
            $a \in S_i$, $b \in S_j$ e $i \neq j$, allora $\lnot (a \mathrel{R} b)$.
    \end{itemize}
    Ogni partizione $\{S_1, \dots, S_m\}$ definisce quindi una relazione di equivalenza,
    in cui due elementi sono in relazione se appartengono allo stesso sottoinsieme:
    \begin{equation*}
        x \mathrel{R} y \iff \exists i.\; x, y \in S_i.
    \end{equation*}
    Viceversa, ogni relazione di equivalenza $R$ su $S$ definisce una partizione di $S$,
    i cui sottoinsiemi raccolgono gli elementi in relazione tra loro.

    \begin{definition}[Classe di equivalenza]
        Sia $R$ una relazione di equivalenza su $S$ e sia $y \in S$. La \emph{classe di
        equivalenza} di $y$ \`e l'insieme degli elementi in relazione con $y$:
        \begin{equation*}
            [y]_R = \{x \in S \mid x \mathrel{R} y\}.
        \end{equation*}
    \end{definition}
    I sottoinsiemi della partizione definita da $R$ sono esattamente le classi di
    equivalenza: ogni $S_i$ \`e la classe $[y]_R$ di un qualsiasi suo elemento $y$, e
    \begin{equation*}
        \{S_1, \dots, S_m\} = \{[y]_R \mid y \in S\}.
    \end{equation*}

    \subsection{Elementi minimali e massimali}
    Nel seguito una relazione d'ordine \`e indicata con $\leq$. La negazione di una relazione
    si indica barrandone il simbolo: $y \not\leq x$ significa che $y \leq x$ non vale.

    \begin{definition}[Elementi minimali e massimali]
        Un elemento $x \in S$ \`e \emph{minimale} rispetto a $\leq$ in $S$ se nessun altro
        elemento di $S$ \`e minore o uguale a $x$:
        \begin{equation*}
            \forall y \in S \setminus \{x\}.\; y \not\leq x.
        \end{equation*}
        Analogamente, $x \in S$ \`e \emph{massimale} rispetto a $\leq$ in $S$ se nessun
        altro elemento di $S$ \`e maggiore o uguale a $x$:
        \begin{equation*}
            \forall y \in S \setminus \{x\}.\; x \not\leq y.
        \end{equation*}
    \end{definition}
    In entrambe le definizioni \`e necessario escludere $y = x$: poich\'e $\leq$ \`e
    riflessiva, vale sempre $x \leq x$, e la condizione $\forall y \in S.\; y \not\leq x$
    non sarebbe mai soddisfatta.

    Se l'elemento minimale \`e unico si chiama \emph{minimo}; se l'elemento massimale \`e
    unico si chiama \emph{massimo}.\footnote{Questo vale per gli insiemi finiti. In
    generale il minimo \`e l'elemento $x$ tale che $\forall y \in S.\; x \leq y$ e il massimo
    \`e l'elemento $x$ tale che $\forall y \in S.\; y \leq x$; in un insieme infinito pu\`o
    esistere un unico elemento massimale che non \`e il massimo.}

    \begin{esempio}
        In $S = \{\{1\}, \{2\}, \{1, 2\}\}$ con la relazione $\subseteq$, gli elementi
        minimali sono $\{1\}$ e $\{2\}$, dunque non esiste un minimo, mentre l'unico elemento
        massimale $\{1, 2\}$ \`e il massimo.
    \end{esempio}

    \subsection{Funzioni}
    \begin{definition}[Funzione]
        Una funzione \`e un caso particolare di relazione. Dati due insiemi $D$, il
        \emph{dominio}, e $C$, il \emph{codominio}, una relazione
        \begin{equation*}
            f \subseteq D \times C
        \end{equation*}
        \`e una \emph{funzione} da $D$ a $C$, e si scrive $f \colon D \to C$, se
        \begin{equation*}
            \forall x \in D.\; \exists!\, y \in C \text{ tale che } f(x) = y.
        \end{equation*}
    \end{definition}
    Qui $f(x) = y$ significa $(x, y) \in f$ e il simbolo $\exists!$ si legge \emph{esiste
    un unico}. Equivalentemente, a ogni elemento del dominio corrisponde uno e un solo
    elemento del codominio.

    \begin{definition}[Immagine]
        L'\emph{immagine} di $f$, $\operatorname{range}(f) \subseteq C$, \`e l'insieme dei
        valori del codominio effettivamente assunti da $f$:
        \begin{equation*}
            \operatorname{range}(f)
            = \{\, c \in C \mid \exists d \in D.\; f(d) = c \,\}
            = \{\, f(d) \in C \mid d \in D \,\}.
        \end{equation*}
    \end{definition}

    \begin{figure}[H]
        \centering
        \begin{tikzpicture}[
            freccia/.style={-{Stealth}, thick}
        ]
            \draw[thick] (0, 0) ellipse (1.1 and 1.5);
            \node at (0, 1.85) {$D$};
            \draw[thick] (5, 0) ellipse (1.3 and 1.5);
            \node at (5, 1.85) {$C$};
            \draw[thick, evidenza, fill=tenue] (4.85, -0.1) ellipse (0.65 and 0.8);

            \node[punto] (d1) at (-0.2, 0.7) {};
            \node[punto] (d2) at (0.3, -0.1) {};
            \node[punto] (d3) at (-0.1, -0.8) {};
            \node[punto] (c1) at (4.7, 0.3) {};
            \node[punto] (c2) at (5.0, -0.5) {};
            \node[punto] at (5.6, 1.0) {};

            \draw[freccia] (d1) to[bend left=15] (c1);
            \draw[freccia] (d2) to[bend left=10] (c1);
            \draw[freccia] (d3) to[bend right=10] (c2);
            \node at (2.3, 1.15) {$f$};

            \node[evidenza, align=left, anchor=west] (range) at (6.7, -0.6)
                {$\operatorname{range}(f)$\\immagine di $f$};
            \draw[evidenza] (range.west) -- (5.45, -0.35);
        \end{tikzpicture}
        \caption{Una funzione $f \colon D \to C$ e la sua immagine
            $\operatorname{range}(f)$.}
        \label{fig:immagine}
    \end{figure}

    Nella Figura~\ref{fig:immagine} ogni elemento di $D$ \`e origine di esattamente una
    freccia, come richiesto dalla definizione di funzione, mentre un elemento di $C$ pu\`o
    esserne destinazione pi\`u volte oppure mai; in quest'ultimo caso non appartiene
    all'immagine.

    \begin{definition}[Funzione iniettiva]
        Una funzione $f \colon D \to C$ \`e \emph{iniettiva} se elementi diversi del
        dominio hanno sempre immagini diverse:
        \begin{equation*}
            \forall x, y \in D.\; f(x) = f(y) \implies x = y.
        \end{equation*}
    \end{definition}
    Se $f$ \`e iniettiva, dunque, due elementi distinti del dominio non possono avere la
    stessa immagine.

    \begin{definition}[Funzione suriettiva]
        Una funzione $f \colon D \to C$ \`e \emph{suriettiva} se ogni elemento del
        codominio \`e raggiunto da almeno un elemento del dominio:
        \begin{equation*}
            \forall y \in C.\; \exists x \in D.\; f(x) = y.
        \end{equation*}
    \end{definition}

    \begin{definition}[Funzione biettiva]
        Una funzione \`e \emph{biettiva}, o \emph{biezione}, se \`e sia iniettiva sia
        suriettiva: ogni elemento del codominio \`e raggiunto da \emph{esattamente un}
        elemento del dominio.
    \end{definition}
    Due insiemi in biezione hanno la stessa cardinalit\`a.

    \begin{definition}[Immagine inversa]
        Data $f \colon D \to C$ e un sottoinsieme $Y \subseteq C$, l'\emph{immagine
        inversa} di $Y$ \`e l'insieme degli elementi del dominio che $f$ manda in $Y$:
        \begin{equation*}
            f^{-1}(Y) = \{\, x \in D \mid f(x) \in Y \,\}.
        \end{equation*}
    \end{definition}

    \begin{figure}[H]
        \centering
        \begin{tikzpicture}[
            freccia/.style={-{Stealth}, thick}
        ]
            \draw[thick] (0, 0) ellipse (1.3 and 1.5);
            \node at (0, 1.85) {$D$};
            \draw[thick] (5, 0) ellipse (1.3 and 1.5);
            \node at (5, 1.85) {$C$};
            \draw[thick, evidenza, fill=tenue] (0.2, 0.0) ellipse (0.6 and 0.72);
            \node[evidenza] at (0.2, 1.05) {$f^{-1}(Y)$};
            \draw[thick, evidenza, fill=tenue] (4.8, 0.1) ellipse (0.5 and 0.6);
            \node[evidenza] at (4.8, 0.95) {$Y$};

            \node[punto] (x1) at (0.1, 0.5) {};
            \node[punto] (x2) at (0.4, 0.0) {};
            \node[punto] (x3) at (0.1, -0.4) {};
            \node[punto] (y1) at (4.7, 0.35) {};
            \node[punto] (y2) at (4.9, -0.2) {};

            \draw[freccia] (x1) to[bend left=10] (y1);
            \draw[freccia] (x2) to (y1);
            \draw[freccia] (x3) to[bend right=5] (y2);
        \end{tikzpicture}
        \caption{Immagine inversa $f^{-1}(Y)$ di un sottoinsieme $Y \subseteq C$.}
        \label{fig:inversa}
    \end{figure}

    \chapter{Introduzione}
    \section{Il problema dell'informatica}
    Si consideri una generica funzione matematica $f \colon \mathbb{N} \to \mathbb{N}$. Il
    problema di fondo dell'informatica consiste nello stabilire \emph{quali funzioni siano
    calcolabili}.

    \begin{definition}[Funzione calcolabile]
        \begin{equation*}
            f \text{ \`e \emph{calcolabile}} \iff \text{esiste un algoritmo che la calcola},
        \end{equation*}
        dove per \emph{algoritmo} si intende una sequenza finita di passi di calcolo.
    \end{definition}
    Una stessa funzione calcolabile pu\`o essere calcolata da pi\`u algoritmi diversi, come
    illustrato nella Figura~\ref{fig:calcolabili}.

    \begin{figure}[H]
        \centering
        \begin{tikzpicture}[
            freccia/.style={-{Stealth}, thick, evidenza}
        ]
            % Tutte le funzioni da N in N
            \draw[thick] (0, 0) rectangle (6, 3.6);
            \node[anchor=north west] at (0, 3.6) {$f \colon \mathbb{N} \to \mathbb{N}$};

            % Funzioni calcolabili
            \draw[thick, evidenza, fill=tenue] (4.2, 0.3) rectangle (5.7, 1.7);
            \node[evidenza, font=\footnotesize, align=center, anchor=south]
                at (4.95, 1.7) {funzioni\\calcolabili};
            \node[punto] (f1) at (4.6, 0.8) {};
            \node[punto] (f2) at (5.3, 1.2) {};

            % Algoritmi
            \draw[thick, evidenza] (7.8, 0) rectangle (10.2, 2.6);
            \node[evidenza, anchor=south] at (9, 2.6) {Algoritmi};
            \node[punto] (a1) at (8.4, 0.6) {};
            \node[punto] (a2) at (8.8, 1.9) {};
            \node[punto] (a3) at (9.6, 1.1) {};

            \draw[freccia] (f1) to[bend right=25] (a1);
            \draw[freccia] (f1) to[bend left=10] (a3);
            \draw[freccia] (f2) to[bend left=20] (a2);
        \end{tikzpicture}
        \caption{Le funzioni calcolabili sono quelle per cui esiste almeno un algoritmo che
            le calcola.}
        \label{fig:calcolabili}
    \end{figure}

    Per stabilire se tutte le funzioni siano calcolabili, si confronta la cardinalit\`a
    dell'insieme delle funzioni con quella dell'insieme degli algoritmi.

    \section{Quante funzioni esistono?}
    Una funzione $f \colon \mathbb{N} \to \mathbb{N}$ pu\`o essere vista come un \emph{insieme
    di coppie} in $\mathbb{N} \times \mathbb{N}$:
    \begin{equation*}
        f(2) = 3 \;\mapsto\; (2, 3) \qquad
        f(3) = 5 \;\mapsto\; (3, 5) \qquad
        f(7) = 9 \;\mapsto\; (7, 9) \qquad
    \end{equation*}

    \begin{esempio}
        Una funzione \`e un insieme infinito di coppie, come
        \begin{equation*}
            \{(0, 3), (1, 5), (2, 2), (3, 100), (4, 1000), (5, 20), (6, 2), \dots\},
        \end{equation*}
        in cui ogni numero naturale compare come primo elemento di \emph{esattamente una}
        coppia: non \`e possibile aggiungere un'altra coppia con primo elemento $2$,
        poich\'e $f(2)$ \`e gi\`a fissato.
    \end{esempio}

    Pertanto
    \begin{equation*}
        f \subseteq \mathbb{N} \times \mathbb{N}
        \iff
        f \in \mathcal{P}(\mathbb{N} \times \mathbb{N}),
    \end{equation*}
    dove $\mathcal{P}(X)$ \`e l'\emph{insieme delle parti} di $X$, ovvero l'insieme di tutti
    i suoi sottoinsiemi. Le funzioni da $\mathbb{N}$ in $\mathbb{N}$ sono quindi elementi di
    $\mathcal{P}(\mathbb{N} \times \mathbb{N})$.

    \subsection{\texorpdfstring{$\mathbb{N} \times \mathbb{N}$}{N x N} \`e numerabile}
    Vale $|\mathbb{N} \times \mathbb{N}| = |\mathbb{N}|$. Le seguenti affermazioni sono
    equivalenti:
    \begin{itemize}
        \item la cardinalit\`a di $\mathbb{N} \times \mathbb{N}$ \`e uguale a quella di
            $\mathbb{N}$;
        \item l'insieme $\mathbb{N} \times \mathbb{N}$ \`e \emph{numerabile};
        \item esiste una \emph{biezione} (mappatura univoca) tra
            $\mathbb{N} \times \mathbb{N}$ e $\mathbb{N}$.
    \end{itemize}
    Infatti, le coppie possono essere ordinate e numerate percorrendole per
    diagonali\footnote{Non \`e invece possibile numerare le coppie riga per riga,
    $(0, 0) \mapsto 0$, $(0, 1) \mapsto 1$, $(0, 2) \mapsto 2$, \dots: la riga delle coppie
    che iniziano con $0$ \`e infinita, quindi la coppia $(1, 0)$ non riceverebbe mai un
    numero.} (Figura~\ref{fig:diagonali}):
    \begin{equation*}
        (0, 0) \mapsto 0 \quad
        (0, 1) \mapsto 1 \quad
        (1, 0) \mapsto 2 \quad
        (0, 2) \mapsto 3 \quad
        (1, 1) \mapsto 4 \quad
        (2, 0) \mapsto 5 \quad
        \dots
    \end{equation*}

    \begin{figure}[H]
        \centering
        \begin{tikzpicture}[x=1.6cm, y=1.1cm]
            % Coppia (x, y) in colonna y e riga x, col suo numero d'ordine
            \foreach \x in {0,...,3} {
                \foreach \y in {0,...,3} {
                    \ifnum\numexpr\x+\y\relax<4
                        \pgfmathtruncatemacro{\n}{(\x+\y)*(\x+\y+1)/2+\x}
                        \node[fill=white, inner sep=2pt] (p\x\y) at (\y, -\x) {$(\x, \y)$};
                        \node[evidenza, font=\scriptsize, anchor=south east, inner sep=1pt]
                            at (p\x\y.north west) {\n};
                    \fi
                }
            }
            \node at (4, 0) {$\cdots$};
            \node at (0, -4) {$\vdots$};

            \begin{scope}[on background layer]
                \foreach \d in {1, 2, 3}
                    \draw[-{Stealth}, black!55] (p0\d) -- (p\d0);
            \end{scope}
        \end{tikzpicture}
        \caption{Enumerazione delle coppie di $\mathbb{N} \times \mathbb{N}$ per diagonali.}
        \label{fig:diagonali}
    \end{figure}

    Poich\'e $\mathbb{N} \times \mathbb{N}$ e $\mathbb{N}$ sono in biezione, lo sono anche i
    loro insiemi delle parti. Quindi
    \begin{equation*}
        |\{f \mid f \colon \mathbb{N} \to \mathbb{N}\}|
        = |\mathcal{P}(\mathbb{N} \times \mathbb{N})|
        = |\mathcal{P}(\mathbb{N})|,
    \end{equation*}
    ovvero esistono \emph{tante funzioni quanti sono i sottoinsiemi di $\mathbb{N}$}.

    \section{Quanti algoritmi esistono?}
    \label{sec:algoritmi}
    Sia $\Sigma$ un \emph{alfabeto}, ovvero un insieme finito e non vuoto di simboli: quelli
    utilizzati per scrivere programmi e algoritmi. Un \emph{programma} \`e una
    \emph{stringa} (sequenza di simboli) su $\Sigma$, quindi
    \begin{equation*}
        |\text{Programmi}| = |\text{stringhe su } \Sigma| = |\Sigma^*|,
    \end{equation*}
    dove l'operatore $^*$ \`e la \emph{chiusura di Kleene}: $\Sigma^*$ \`e l'insieme di
    tutte le stringhe ottenute concatenando zero o pi\`u simboli di $\Sigma$. Ogni stringa
    di $\Sigma^*$ \`e quindi \emph{finita}, anche se l'insieme \`e infinito:
    \begin{equation*}
        \underbrace{\Sigma = \{a, b\}}_{\text{finito}}
        \qquad\qquad
        \underbrace{\Sigma^* = \{\varepsilon, a, b, aa, ab, ba, bb, aaa, \dots\}}_
        {\text{infinito}}
    \end{equation*}
    dove $\varepsilon$ \`e la \emph{stringa vuota}.

    $\Sigma^*$ pu\`o essere messo in biezione con $\mathbb{N}$ poich\'e i suoi elementi sono
    ordinabili (prima per lunghezza, poi in ordine alfabetico, come nell'esempio): \`e
    sufficiente numerarli secondo tale ordine. Quindi $\Sigma^*$ \`e numerabile,
    $|\Sigma^*| = |\mathbb{N}|$, e
    \begin{equation*}
        |\text{Programmi}| = |\Sigma^*| = |\mathbb{N}|.
    \end{equation*}

    \section{Il teorema di Cantor}
    Resta da confrontare $|\mathbb{N}|$, la cardinalit\`a degli algoritmi, con
    $|\mathcal{P}(\mathbb{N})|$, la cardinalit\`a delle funzioni.

    \begin{theorem}[di Cantor]
        \label{thm:cantor}
        Per ogni insieme $S$ si ha che
        \begin{equation*}
            |S| < |\mathcal{P}(S)|.
        \end{equation*}
    \end{theorem}

    Per gli insiemi finiti il teorema \`e verificabile direttamente, come mostrato nella
    Tabella~\ref{tab:parti}.

    \begin{table}[H]
        \centering
        \begin{tabular}{llcc}
            \toprule
            \textsc{Insieme} $S$ & \textsc{Insieme delle parti} $\mathcal{P}(S)$ & $|S|$
                & $|\mathcal{P}(S)|$ \\
            \midrule
            $\emptyset$ & $\{\emptyset\}$ & $0$ & $1$ \\
            $\{a, b\}$ & $\{\emptyset, \{a\}, \{b\}, \{a, b\}\}$ & $2$ & $4$ \\
            $\{a, b, c\}$ & $\{\emptyset, \{a\}, \{b\}, \{c\}, \{a, b\}, \{b, c\}, \{a, c\},
                \{a, b, c\}\}$ & $3$ & $8$ \\
            \bottomrule
        \end{tabular}
        \caption{Cardinalit\`a di alcuni insiemi finiti e dei loro insiemi delle parti.}
        \label{tab:parti}
    \end{table}

    \begin{proof}[Dimostrazione del Teorema~\ref{thm:cantor}]
        Si suppone per assurdo che $|S| = |\mathcal{P}(S)|$. Allora esiste una funzione
        biettiva $f \colon S \to \mathcal{P}(S)$, che associa a ogni elemento di $S$ un
        sottoinsieme di $S$ stesso.

        Si consideri l'insieme $A = \{x \in S \mid x \notin f(x)\} \subseteq S$. Poich\'e
        $\mathcal{P}(S)$ contiene ogni sottoinsieme di $S$, si ha $A \in \mathcal{P}(S)$ e,
        dato che $f$ \`e suriettiva, esiste $a \in S$ tale che $f(a) = A$. Si hanno quindi due
        possibilit\`a:
        \begin{itemize}
            \item $a \in A$: allora, per definizione di $A$, $a \notin f(a) = A$, assurdo
                ($\lightning$);
            \item $a \notin A$: allora, per definizione di $A$, $a \in f(a) = A$, assurdo
                ($\lightning$).
        \end{itemize}
        In entrambi i casi si giunge a una contraddizione, quindi non esiste $a \in S$ tale
        che $f(a) = A$: l'insieme $A$ non ha controimmagine e $f$ non pu\`o essere
        suriettiva. Di conseguenza $|S| \neq |\mathcal{P}(S)|$ e, poich\'e
        $x \mapsto \{x\}$ \`e una funzione iniettiva da $S$ in $\mathcal{P}(S)$,
        $|S| < |\mathcal{P}(S)|$.
    \end{proof}

    In generale
    \begin{equation*}
        |\mathcal{P}(S)| = 2^{|S|}.
    \end{equation*}

    Per $\mathbb{N}$ vale $|\mathcal{P}(\mathbb{N})| = |\mathbb{R}|$: i sottoinsiemi di
    $\mathbb{N}$ sono tanti quanti i numeri reali, che \emph{non sono numerabili}.
    L'\emph{ipotesi del continuo} afferma inoltre che non esiste alcun insieme con
    cardinalit\`a strettamente compresa tra $|\mathbb{N}|$ e $|\mathbb{R}|$.

    Riassumendo:
    \begin{equation*}
        |\text{Programmi}| = |\mathbb{N}|
        < |\mathcal{P}(\mathbb{N})|
        = |\{f \mid f \colon \mathbb{N} \to \mathbb{N}\}|.
    \end{equation*}
    Poich\'e la cardinalit\`a dei programmi \`e strettamente inferiore a quella delle
    funzioni, esistono \emph{funzioni non calcolabili}.

    \chapter{Induzione}
    \section{Relazioni ben fondate}
    \begin{definition}[Relazione ben fondata]
        Sia $A$ un insieme e $\prec \; \subseteq A \times A$ una relazione su $A$,
        antisimmetrica e transitiva. La relazione $\prec$ \`e \emph{ben fondata} se non
        esistono \emph{catene discendenti infinite}, ovvero sequenze infinite
        $a_0 \succ a_1 \succ a_2 \succ \dots$ di elementi di $A$.
    \end{definition}

    Ne segue che, se $A$ non \`e vuoto, esistono \emph{elementi minimali}, ovvero elementi
    di $A$ privi di elementi pi\`u piccoli. L'insieme degli elementi minimali di $A$ si
    indica con $\operatorname{Base}(A)$:
    \begin{equation*}
        \operatorname{Base}(A) = \{\, a \in A \mid \nexists b \in A.\; b \prec a \,\}.
    \end{equation*}
    La situazione \`e illustrata nella Figura~\ref{fig:base}.

    \begin{figure}[H]
        \centering
        \begin{tikzpicture}[
            minimale/.style={circle, fill=evidenza, inner sep=1.8pt}
        ]
            \draw[thick] (0, 0) circle (2);
            \node at (1.7, 1.7) {$A$};

            % Catene discendenti: ogni freccia va verso un elemento pi\`u piccolo
            \foreach \x/\n [count=\i] in {-1/3, 0/4, 1/2} {
                \node[minimale] (m\i) at (\x, -1.3) {};
                \foreach \k in {1,...,\n}
                    \node[punto] (p\i\k) at (\x, -1.3 + 0.65*\k) {};
                \draw[-{Stealth}] (p\i1) -- (m\i);
                \foreach \k [evaluate=\k as \j using int(\k-1)] in {2,...,\n}
                    \draw[-{Stealth}] (p\i\k) -- (p\i\j);
            }

            \node[draw=evidenza, dashed, rounded corners, fit=(m1)(m3), inner sep=5pt]
                (base) {};
            \draw[evidenza] (base.east) -- (2.6, -1.3)
                node[right, align=left] {$\operatorname{Base}(A)$\\elementi minimali};
        \end{tikzpicture}
        \caption{Catene discendenti in un insieme ben fondato e suoi elementi minimali.}
        \label{fig:base}
    \end{figure}

    \section{Principio di induzione}
    \begin{theorem}[Principio di induzione]
        \label{thm:induzione}
        Sia $(A, \prec)$ ben fondato e sia $\Pi$ una propriet\`a sugli elementi di $A$.
        Allora $\forall a \in A.\; \Pi(a)$ se e solo se valgono:
        \begin{enumerate}[label=\textup{(\alph*)}]
            \item \textup{\textsc{Base dell'induzione}}: $\forall a \in
                \operatorname{Base}(A).\; \Pi(a)$;
            \item \textup{\textsc{Passo induttivo}}: per ogni
                $a \in A \setminus \operatorname{Base}(A)$
                \begin{equation*}
                    \underbrace{\bigl(\forall b \prec a.\; \Pi(b)\bigr)}_{\text{ipotesi
                    induttiva}}
                    \implies
                    \underbrace{\Pi(a)}_{\text{tesi}}
                \end{equation*}
        \end{enumerate}
    \end{theorem}

    \begin{esempio}
        Si vuole dimostrare che
        \begin{equation*}
            \forall n \geq 1. \quad \sum_{i=1}^{n} i = \frac{n(n+1)}{2}.
        \end{equation*}
        Sia $\Pi(n)$ tale uguaglianza. In questo caso $A = \mathbb{N} \setminus \{0\}$ con
        l'ordine usuale $<$, che \`e ben fondato, e $\operatorname{Base}(A) = \{1\}$.

        \textsc{Base}. Per $n = 1$, $\Pi(1)$ diventa:
        \begin{equation*}
            \sum_{i=1}^{1} i = 1
            \qquad\qquad
            \frac{1 \cdot (1 + 1)}{2} = \frac{1 \cdot 2}{2} = 1.
        \end{equation*}
        I due membri coincidono, dunque $\Pi(1)$ \`e verificata.

        \textsc{Passo induttivo}. Sia $n \in \mathbb{N} \setminus \{0, 1\}$ e si assuma
        l'\emph{ipotesi induttiva}, ovvero che $\Pi(m)$ valga per ogni $m < n$. In
        particolare vale $\Pi(n - 1)$, utilizzata nel secondo passaggio:
        \begin{align*}
            \sum_{i=1}^{n} i
                &= \sum_{i=1}^{n-1} i + n \\
                &= \frac{(n-1)(n-1+1)}{2} + n && \text{(ipotesi induttiva)} \\
                &= \frac{(n-1)n + 2n}{2} \\
                &= \frac{n(n-1+2)}{2} \\
                &= \frac{n(n+1)}{2},
        \end{align*}
        che \`e esattamente $\Pi(n)$.
    \end{esempio}

    \chapter{Linguaggi formali}
    \section{Alfabeti e stringhe}
    \begin{definition}[Alfabeto]
        Un \emph{alfabeto} $\Sigma$ \`e un insieme finito e non vuoto di simboli, ovvero di
        entit\`a primitive come lettere e caratteri.
    \end{definition}

    \begin{definition}[Stringa]
        Una \emph{stringa}, o \emph{parola}, \`e una sequenza finita di simboli su un
        alfabeto $\Sigma$.
    \end{definition}

    Come visto nella Sezione~\ref{sec:algoritmi}, l'insieme di tutte le stringhe su
    $\Sigma$ si indica con $\Sigma^*$. Esso contiene sempre la \emph{stringa vuota}
    $\varepsilon$, priva di simboli, ed \`e infinito numerabile.

    \begin{definition}[Lunghezza di una stringa]
        La \emph{lunghezza} $|x|$ di una stringa $x$ \`e il numero di \emph{occorrenze} di
        simboli al suo interno.
    \end{definition}

    Si contano le occorrenze e non i simboli: ad esempio $|abba| = 4$, anche se $abba$
    contiene solo due simboli distinti. La stringa vuota ha lunghezza $|\varepsilon| = 0$.

    \begin{definition}[Prefisso, suffisso e sottostringa]
        Sia $x = a_1 a_2 \cdots a_n$ una stringa su $\Sigma$, con $a_1, \dots, a_n \in \Sigma$,
        e siano $1 \leq i \leq j \leq n$. Allora:
        \begin{itemize}
            \item $a_1 \cdots a_i$, la parte iniziale di $x$, \`e un \emph{prefisso} di $x$;
            \item $a_j \cdots a_n$, la parte finale di $x$, \`e un \emph{suffisso} di $x$;
            \item $a_i \cdots a_j$, una qualsiasi parte contigua di $x$, \`e una
                \emph{sottostringa} di $x$.
        \end{itemize}
    \end{definition}

    Anche la stringa vuota $\varepsilon$ \`e prefisso, suffisso e sottostringa di ogni
    stringa. Come mostra la Figura~\ref{fig:sottostringhe}, ogni prefisso e ogni suffisso
    di $x$ sono anche sue sottostringhe.

    \begin{figure}[H]
        \centering
        \begin{tikzpicture}[
            simbolo/.style={inner sep=1pt, text height=1.6ex, text depth=0.5ex},
            graffa/.style={decorate, thick,
                decoration={brace, amplitude=5pt, raise=3pt}},
            etichetta/.style={font=\small\itshape}
        ]
            % Simboli di x, affiancati sulla stessa linea di base
            \node[simbolo] (s1) {$a_1$};
            \foreach \s [count=\k from 2, evaluate=\k as \j using int(\k-1)]
                in {a_2, \cdots, a_i, \cdots, a_j, \cdots, a_{n-1}, a_n}
                \node[simbolo, anchor=base west] (s\k)
                    at ([xshift=6pt] s\j.base east) {$\s$};

            % Sottostringa sopra, prefisso e suffisso sotto
            \draw[graffa] (s4.north west) -- (s6.north east)
                node[midway, above=9pt, etichetta] {sottostringa};
            \draw[graffa] (s4.south east) -- (s1.south west)
                node[midway, below=9pt, etichetta] {prefisso};
            \draw[graffa] (s9.south east) -- (s6.south west)
                node[midway, below=9pt, etichetta] {suffisso};
        \end{tikzpicture}
        \caption{Prefisso $a_1 \cdots a_i$, sottostringa $a_i \cdots a_j$ e suffisso
            $a_j \cdots a_n$ della stringa $x = a_1 a_2 \cdots a_n$.}
        \label{fig:sottostringhe}
    \end{figure}

    \section{Linguaggi}
    \begin{definition}[Linguaggio]
        Un \emph{linguaggio} $\mathcal{L}$ su un alfabeto $\Sigma$ \`e un sottoinsieme delle
        stringhe di $\Sigma^*$, ovvero $\mathcal{L} \subseteq \Sigma^*$.
    \end{definition}

    Alcuni esempi di linguaggi sono $\emptyset$, dato che \`e sottoinsieme di qualsiasi
    insieme, e $\{\varepsilon\}$, il linguaggio che contiene solo la stringa vuota. I due
    non vanno confusi: $\emptyset$ non contiene alcuna stringa, mentre $\{\varepsilon\}$
    ne contiene esattamente una, di lunghezza $0$. Entrambi sono linguaggi su qualsiasi
    alfabeto, mentre $\Sigma^*$ \`e il linguaggio di tutte le stringhe su $\Sigma$.

    Poich\'e un linguaggio \`e, di fatto, un insieme, su di esso si possono eseguire le
    operazioni insiemistiche della Sezione~\ref{sec:operazioni-insiemi}, con $\Sigma^*$
    come insieme universo:
    \begin{itemize}
        \item \textsc{Complemento} di $\mathcal{L} \subseteq \Sigma^*$:
            $\overline{\mathcal{L}} = \neg \mathcal{L} = \Sigma^* \setminus \mathcal{L}$;
        \item \textsc{Unione} di $\mathcal{L}_1, \mathcal{L}_2 \subseteq \Sigma^*$:
            $\mathcal{L}_1 \cup \mathcal{L}_2 = \{x \in \Sigma^* \mid x \in \mathcal{L}_1
            \lor x \in \mathcal{L}_2\}$;
        \item \textsc{Intersezione} di $\mathcal{L}_1, \mathcal{L}_2 \subseteq \Sigma^*$:
            $\mathcal{L}_1 \cap \mathcal{L}_2 = \{x \in \Sigma^* \mid x \in \mathcal{L}_1
            \land x \in \mathcal{L}_2\}$.
    \end{itemize}
    Il complemento \`e quindi sempre relativo a $\Sigma^*$. Ad esempio, il complemento del
    linguaggio vuoto \`e il linguaggio di tutte le stringhe, mentre quello di
    $\{\varepsilon\}$ \`e il linguaggio delle stringhe non vuote:
    \begin{equation*}
        \overline{\emptyset} = \Sigma^* \setminus \emptyset = \Sigma^*,
        \qquad\qquad
        \overline{\{\varepsilon\}} = \{x \in \Sigma^* \mid x \neq \varepsilon\}
        = \Sigma^* \setminus \{\varepsilon\}.
    \end{equation*}

    Altre operazioni, come la concatenazione e la chiusura di Kleene, sono invece
    specifiche dei linguaggi e sono descritte nella Sezione~\ref{sec:operazioni-linguaggi}.

    \section{Funzioni come linguaggi}
    Una funzione $f \colon \mathbb{N} \to \mathbb{N}$ pu\`o essere vista come un linguaggio
    $\mathcal{L}_f$ su un alfabeto $\Sigma$ composto da un singolo simbolo.

    \begin{esempio}
        Si consideri l'alfabeto $\Sigma = \{1\}$ e una funzione
        $f \colon \mathbb{N} \to \mathbb{N}$. Allora
        \begin{equation*}
            f \colon \mathbb{N} \to \mathbb{N}
            \quad\longleftrightarrow\quad
            \mathcal{L}_f = \{1^{f(x)} \mid x \in \mathbb{N}\},
        \end{equation*}
        dove $1^n$ indica la stringa formata da $n$ occorrenze del simbolo $1$, con
        $1^0 = \varepsilon$.
    \end{esempio}

    Questo linguaggio permette di porre una domanda con risposta s\`i o no: data in input
    una stringa $\sigma \in \Sigma^*$, stabilire se essa appartenga o meno al linguaggio,
    \begin{equation*}
        \sigma \in \Sigma^*
        \;\longmapsto\;
        \begin{cases}
            \sigma \in \mathcal{L}_f & \text{se $\sigma$ appartiene al linguaggio,} \\
            \sigma \notin \mathcal{L}_f & \text{se $\sigma$ non appartiene al linguaggio.}
        \end{cases}
    \end{equation*}
    Si ragiona cos\`i su insiemi invece che su funzioni, dove gli elementi dell'insieme
    dipendono dal calcolo della funzione.

    \begin{esempio}
        Per $f(x) = 2x$ si ottiene
        \begin{equation*}
            \mathcal{L}_f = \{1^{2x} \mid x \in \mathbb{N}\}
            = \{\varepsilon, 11, 1111, 111111, \dots\},
        \end{equation*}
        il linguaggio delle stringhe di lunghezza pari su $\{1\}$: ad esempio
        $1111 \in \mathcal{L}_f$, poich\'e $f(2) = 4$, mentre $111 \notin \mathcal{L}_f$,
        poich\'e nessun $x$ ha $f(x) = 3$.
    \end{esempio}

    \section{Automi a stati finiti deterministici}
    I \emph{linguaggi regolari} sono i linguaggi per cui esiste un automa a stati finiti
    deterministico che li riconosce.

    Un \emph{automa a stati finiti} \`e un modello matematico di un sistema con input, ed
    eventualmente output, a valori discreti. Il sistema si trova sempre in uno stato, scelto
    in un insieme finito di stati possibili, e l'essere in uno stato gli permette di tenere
    traccia della storia precedente.

    Un automa a stati finiti si pu\`o rappresentare come una testina che scorre, sempre
    nella stessa direzione, un nastro di lunghezza illimitata contenente dei simboli. La
    testina si trova in un certo \emph{stato} $q$ e, a seconda del simbolo $s_i$ che legge,
    passa in un altro stato oppure resta nello stesso. La situazione \`e illustrata nella
    Figura~\ref{fig:nastro}.

    \begin{figure}[H]
        \centering
        \begin{tikzpicture}[
            cella/.style={draw, thick, minimum width=1cm, minimum height=0.8cm}
        ]
            % Nastro illimitato a destra, con i simboli da leggere
            \foreach \s [count=\k] in {s_1, s_2, s_3}
                \node[cella] (c\k) at (\k, 0) {$\s$};
            \node[minimum width=1cm, minimum height=0.8cm] (c4) at (4, 0) {$\cdots$};
            \draw[thick] (c4.north west) -- (c4.north east)
                (c4.south west) -- (c4.south east);

            % Testina sul primo simbolo, nello stato corrente q
            \node[draw, thick, minimum size=0.7cm] (q) at (1, -1.9) {$q$};
            \draw[evidenza, double, double distance=1.6pt, -{Implies}]
                ([yshift=3pt] q.north) -- ([yshift=-4pt] c1.south);
        \end{tikzpicture}
        \caption{Automa a stati finiti: la testina, nello stato $q$, legge il primo simbolo
            $s_1$ del nastro.}
        \label{fig:nastro}
    \end{figure}

    Quando la lettura dei simboli termina, a seconda dello stato raggiunto dalla testina,
    l'automa fornisce un risultato di accettazione o di rifiuto della stringa letta.

    Il comportamento dell'automa si definisce in maniera univoca mediante una tabella
    detta \emph{matrice di transizione}.

    \begin{esempio}
        La matrice di transizione della Tabella~\ref{tab:matrice} descrive un automa con
        stati $q_0$, $q_1$, $q_2$ sull'alfabeto $\{a, b\}$. Ogni riga corrisponde a uno
        stato e ogni colonna a un simbolo: la casella nella riga $q$ e nella colonna $a$
        contiene lo stato in cui l'automa passa leggendo $a$ dallo stato $q$. Ad esempio,
        dallo stato $q_0$, leggendo $b$, l'automa passa in $q_2$.

        \begin{table}[H]
            \centering
            \begin{tabular}{ccc}
                \toprule
                $\delta$ & $a$ & $b$ \\
                \midrule
                $q_0$ & $q_1$ & $q_2$ \\
                $q_1$ & $q_1$ & $q_0$ \\
                $q_2$ & $q_1$ & $q_0$ \\
                \bottomrule
            \end{tabular}
            \caption{Matrice di transizione di un automa con tre stati.}
            \label{tab:matrice}
        \end{table}
    \end{esempio}

    In alternativa, l'automa si rappresenta con un grafo: gli stati sono i nodi e ogni
    transizione \`e un arco etichettato con il simbolo letto, come negli esempi seguenti.

    \begin{definition}[Automa deterministico]
        Un \emph{automa a stati finiti deterministico} (DFA) \`e una quintupla
        $\langle \mathcal{Q}, \Sigma, \delta, q_0, \mathcal{F} \rangle$, dove:
        \begin{itemize}
            \item $\mathcal{Q}$ \`e un insieme finito di \emph{stati};
            \item $\Sigma$ \`e un alfabeto;
            \item $\delta \colon \mathcal{Q} \times \Sigma \to \mathcal{Q}$ \`e la
                \emph{funzione di transizione};
            \item $q_0 \in \mathcal{Q}$ \`e lo stato \emph{iniziale};
            \item $\mathcal{F} \subseteq \mathcal{Q}$ \`e l'insieme degli stati
                \emph{finali}.
        \end{itemize}
    \end{definition}

    La funzione di transizione di un DFA \`e deterministica: per ogni stato e per ogni
    simbolo esiste sempre una e una sola transizione,
    \begin{equation*}
        \forall q \in \mathcal{Q}.\; \forall a \in \Sigma.\;
        \exists!\, p \in \mathcal{Q}.\; \delta(q, a) = p.
    \end{equation*}
    L'unica forma di memoria dell'automa \`e lo stato in cui si trova.

    \textsc{Notazione}. Si useranno $p, q, r$ per denotare gli stati, $P, Q, R, S$ per gli
    insiemi di stati, $a, b$ per i simboli di $\Sigma$ e $x, y, z, u, v, w$ per le
    stringhe, tutti con o senza pedici.

    Poich\'e la funzione $\delta$ riceve in ingresso un solo simbolo, \`e necessario
    estenderla per poter leggere stringhe. A questo scopo si definisce la funzione
    $\hat{\delta}$, che si ottiene in modo univoco da $\delta$:
    \begin{equation*}
        \hat{\delta} \colon \mathcal{Q} \times \Sigma^* \to \mathcal{Q},
        \qquad
        \begin{cases}
            \hat{\delta}(q, \varepsilon) = q, \\
            \hat{\delta}(q, wa) = \delta(\hat{\delta}(q, w), a).
        \end{cases}
    \end{equation*}

    Una stringa $x$ \`e \emph{accettata}, o \emph{riconosciuta}, da un DFA
    $\mathcal{M} = \langle \mathcal{Q}, \Sigma, \delta, q_0, \mathcal{F} \rangle$ se
    $\hat{\delta}(q_0, x) \in \mathcal{F}$. Il linguaggio riconosciuto (o accettato) da
    $\mathcal{M}$, indicato con $\mathcal{L}(\mathcal{M})$, \`e l'insieme delle stringhe
    accettate:
    \begin{equation*}
        \mathcal{L}(\mathcal{M})
        = \{x \in \Sigma^* \mid \hat{\delta}(q_0, x) \in \mathcal{F}\}.
    \end{equation*}

    \begin{definition}[Linguaggio regolare]
        Un linguaggio $\mathcal{L}$ \`e detto \emph{regolare} se \`e accettato da qualche
        DFA, ovvero se esiste un DFA $\mathcal{M}$ tale che
        \begin{equation*}
            \mathcal{L} = \mathcal{L}(\mathcal{M}).
        \end{equation*}
    \end{definition}

    \begin{esempio}
        Il linguaggio vuoto $\mathcal{L} = \emptyset$ \`e regolare. Lo riconosce il DFA
        $\mathcal{M}$ della Figura~\ref{fig:dfa-vuoto}, con $\mathcal{Q} = \{q_0\}$,
        $\Sigma = \{a, b\}$, stato iniziale $q_0$ e $\mathcal{F} = \emptyset$.

        \begin{figure}[H]
            \centering
            \begin{tikzpicture}[automa]
                \node[state, initial] (q0) {$q_0$};
                \path (q0) edge[loop above] node {$a, b$} ();

                % Matrice di transizione, a destra del diagramma
                \node[right=2.5cm of q0] {%
                    \begin{tabular}{ccc}
                        \toprule
                        $\delta$ & $a$ & $b$ \\
                        \midrule
                        $q_0$ & $q_0$ & $q_0$ \\
                        \bottomrule
                    \end{tabular}};
            \end{tikzpicture}
            \caption{DFA che riconosce il linguaggio vuoto e la sua matrice di
                transizione.}
            \label{fig:dfa-vuoto}
        \end{figure}

        Qualunque simbolo legga, l'automa resta in $q_0$, quindi
        $\hat{\delta}(q_0, x) = q_0$ per ogni $x \in \Sigma^*$. Poich\'e
        $\mathcal{F} = \emptyset$, per\`o, la condizione di accettazione
        $\hat{\delta}(q_0, x) \in \mathcal{F}$ \`e sempre falsa: nessuno stato appartiene
        all'insieme vuoto. Nessuna stringa viene quindi accettata e
        \begin{equation*}
            \mathcal{L}(\mathcal{M})
            = \{x \in \Sigma^* \mid \underbrace{\hat{\delta}(q_0, x) \in \mathcal{F}}_{\text{
            sempre falsa, poich\'e } \mathcal{F} \;=\; \emptyset}\}
            = \emptyset.
        \end{equation*}
        Riconoscere $\emptyset$ non significa che l'automa non esista, ma che non accetta
        alcuna stringa: in generale, ogni DFA con $\mathcal{F} = \emptyset$ riconosce il
        linguaggio vuoto.
    \end{esempio}

    \begin{esempio}
        Il linguaggio $\mathcal{L} = \{\varepsilon\}$ \`e regolare. Lo riconosce il DFA
        $\mathcal{M}$ della Figura~\ref{fig:dfa-epsilon}, con
        $\mathcal{Q} = \{q_0, q_1\}$, $\Sigma = \{a, b\}$, stato iniziale $q_0$ e
        $\mathcal{F} = \{q_0\}$.

        \begin{figure}[H]
            \centering
            \begin{tikzpicture}[automa]
                \node[state, initial, accepting] (q0) {$q_0$};
                \node[state, right=of q0] (q1) {$q_1$};
                \path (q0) edge node {$a, b$} (q1)
                    (q1) edge[loop above] node {$a, b$} ();

                % Matrice di transizione, a destra del diagramma
                \node[right=2cm of q1] {%
                    \begin{tabular}{ccc}
                        \toprule
                        $\delta$ & $a$ & $b$ \\
                        \midrule
                        $q_0$ & $q_1$ & $q_1$ \\
                        $q_1$ & $q_1$ & $q_1$ \\
                        \bottomrule
                    \end{tabular}};
            \end{tikzpicture}
            \caption{DFA che riconosce il linguaggio $\{\varepsilon\}$ e la sua matrice di
                transizione: lo stato finale $q_0$ \`e indicato dal doppio cerchio.}
            \label{fig:dfa-epsilon}
        \end{figure}

        Lo stato iniziale \`e anche finale, quindi la stringa vuota, che non fa compiere
        alcuna transizione, viene accettata. Il primo simbolo letto porta invece in $q_1$,
        da cui l'automa non esce pi\`u:
        \begin{equation*}
            \hat{\delta}(q_0, \varepsilon) = q_0 \in \mathcal{F},
            \qquad\qquad
            \hat{\delta}(q_0, x) = q_1 \notin \mathcal{F}
            \quad \text{per ogni } x \neq \varepsilon.
        \end{equation*}
        Quindi $\mathcal{L}(\mathcal{M}) = \{\varepsilon\}$.
    \end{esempio}

    \subsection{Come dimostrare che un linguaggio \`e regolare}
    \label{sec:dimostrare-regolare}
    Un modo per dimostrare che un linguaggio \`e regolare \`e costruire un DFA che lo
    riconosca. L'uguaglianza $\mathcal{L} = \mathcal{L}(\mathcal{M})$ verr\`a dimostrata
    mediante la duplice inclusione $\mathcal{L} \subseteq \mathcal{L}(\mathcal{M})$ e
    $\mathcal{L}(\mathcal{M}) \subseteq \mathcal{L}$. In particolare:
    \begin{align*}
        & \mathcal{L}(\mathcal{M}) \subseteq \mathcal{L} \\
        \iff {} & \forall x \in \Sigma^*.\;
            \bigl(x \in \mathcal{L}(\mathcal{M}) \implies x \in \mathcal{L}\bigr) \\
        \iff {} & \hat{\delta}(q_0, x) \in \mathcal{F} \implies x \in \mathcal{L}
            && \text{(definizione di $\mathcal{L}(\mathcal{M})$, con $x$ arbitrario)} \\
        \iff {} & x \notin \mathcal{L} \implies \hat{\delta}(q_0, x) \notin \mathcal{F}
            && \text{(per contrapposizione)}
    \end{align*}
    e
    \begin{align*}
        & \mathcal{L} \subseteq \mathcal{L}(\mathcal{M}) \\
        \iff {} & \forall x \in \Sigma^*.\;
            \bigl(x \in \mathcal{L} \implies x \in \mathcal{L}(\mathcal{M})\bigr) \\
        \iff {} & x \in \mathcal{L} \implies \hat{\delta}(q_0, x) \in \mathcal{F}
            && \text{(definizione di $\mathcal{L}(\mathcal{M})$, con $x$ arbitrario)}
    \end{align*}
    Riassumendo:
    \begin{equation*}
        \begin{cases}
            x \in \mathcal{L} \implies \hat{\delta}(q_0, x) \in \mathcal{F}, \\
            x \notin \mathcal{L} \implies \hat{\delta}(q_0, x) \notin \mathcal{F}.
        \end{cases}
    \end{equation*}

    Dal momento che le condizioni da dimostrare sono analoghe, verranno portate avanti
    insieme ed \`e quindi necessario che la base dell'induzione sia la stessa per entrambe.
    In particolare, \`e richiesto di trovare una lunghezza tale per cui esistano stringhe in
    $\mathcal{L}$ e stringhe non in $\mathcal{L}$. I passi sono dunque:
    \begin{itemize}
        \item \textsc{Base}: si considera la lunghezza minima di $x$ e, per tutte le
            lunghezze inferiori o uguali a quella, si dimostra che valgono entrambe le
            condizioni;
        \item \textsc{Passo induttivo}: si suppone che entrambe le condizioni valgano per
            ogni $x$ tale che $|x| \leq n$ e si dimostra che valgono anche per
            $|x| = n + 1$.
    \end{itemize}

    \begin{esempio}
        Si consideri il linguaggio
        $\mathcal{L} = \{x \in \Sigma^* \mid \text{``$x$ contiene almeno due $1$''}\}$
        definito sull'alfabeto binario $\Sigma = \{0, 1\}$. L'automa a stati finiti che
        riconosce $\mathcal{L}$ \`e quello della Figura~\ref{fig:dfa-due-uno}, con stato
        iniziale $q_0$ e $\mathcal{F} = \{q_2\}$.

        \begin{figure}[H]
            \centering
            \begin{tikzpicture}[automa]
                \node[state, initial] (q0) {$q_0$};
                \node[state, right=of q0] (q1) {$q_1$};
                \node[state, accepting, right=of q1] (q2) {$q_2$};
                \path (q0) edge node {$1$} (q1)
                    (q1) edge node {$1$} (q2)
                    (q0) edge[loop above] node {$0$} ()
                    (q1) edge[loop above] node {$0$} ()
                    (q2) edge[loop above] node {$0, 1$} ();

                % Matrice di transizione, a destra del diagramma
                \node[right=1.8cm of q2] {%
                    \begin{tabular}{ccc}
                        \toprule
                        $\delta$ & $0$ & $1$ \\
                        \midrule
                        $q_0$ & $q_0$ & $q_1$ \\
                        $q_1$ & $q_1$ & $q_2$ \\
                        $q_2$ & $q_2$ & $q_2$ \\
                        \bottomrule
                    \end{tabular}};
            \end{tikzpicture}
            \caption{DFA che riconosce le stringhe binarie con almeno due $1$ e la sua
                matrice di transizione.}
            \label{fig:dfa-due-uno}
        \end{figure}

        \`E comunque necessario dimostrare che vengono riconosciute tutte e sole le stringhe
        di $\mathcal{L}$. In particolare, bisogna dimostrare che
        \begin{equation*}
            \mathcal{L} = \mathcal{L}(\mathcal{M})
            = \{x \in \Sigma^* \mid \hat{\delta}(q_0, x) \in \{q_2\}\}.
        \end{equation*}
        La dimostrazione avviene per induzione. La base non pu\`o riguardare \emph{solo}
        stringhe con meno di due simboli: servono almeno due simboli per comporre una stringa
        $z \in \mathcal{L}$.\footnote{In realt\`a, non \`e mai possibile
        considerare $|x| = 0$ come unico caso base, perch\'e l'unica stringa di tale lunghezza,
        $\varepsilon$, o appartiene a un linguaggio o non vi appartiene. Servono, invece,
        lunghezze tali per cui esistono stringhe sia appartenenti che non.} Quindi:
        \begin{equation*}
            \forall x \in \Sigma^*.\; \bigl(|x| < 2 \implies x \notin \mathcal{L}\bigr)
        \end{equation*}
        Allora:
        \begin{equation*}
            \left.
            \begin{alignedat}{3}
                |x| = 0 \quad & x = \varepsilon \notin \mathcal{L} & \qquad
                    & \hat{\delta}(q_0, \varepsilon) = q_0 \notin \mathcal{F} \\
                |x| = 1 \quad & x = 0 \notin \mathcal{L} &
                    & \hat{\delta}(q_0, 0) = q_0 \notin \mathcal{F} \\
                & x = 1 \notin \mathcal{L} &
                    & \hat{\delta}(q_0, 1) = q_1 \notin \mathcal{F} \\[4pt]
                |x| = 2 \quad & x = 00 \notin \mathcal{L} &
                    & \hat{\delta}(q_0, 00) = q_0 \notin \mathcal{F} \\
                & x = 01 \notin \mathcal{L} &
                    & \hat{\delta}(q_0, 01) = q_1 \notin \mathcal{F} \\
                & x = 10 \notin \mathcal{L} &
                    & \hat{\delta}(q_0, 10) = q_1 \notin \mathcal{F} \\
                & x = 11 \in \mathcal{L} &
                    & \hat{\delta}(q_0, 11) = q_2 \in \mathcal{F}
            \end{alignedat}
            \right\}
            \quad \text{\textsc{Base}}
        \end{equation*}

        \textsc{Passo induttivo}. \emph{Ipotesi induttiva}: si suppone che per tutte le
        stringhe $x$, con $|x| \geq 2$, di lunghezza $|x| \leq n$, per un certo
        $n \in \mathbb{N}$, valga che
        \begin{equation*}
            \begin{cases}
                x \in \mathcal{L} \implies \hat{\delta}(q_0, x) \in \mathcal{F}\\
                x \notin \mathcal{L} \implies \hat{\delta}(q_0, x) \notin \mathcal{F}
            \end{cases}
        \end{equation*}
        dove si \`e usato che $\mathcal{F} = \{q_2\}$. La seconda condizione, per\`o, non
        basta: dice solo che l'automa \`e in $q_0$ oppure in $q_1$, mentre per leggere il
        simbolo successivo serve sapere in quale dei due. La si scorpora quindi in due
        condizioni pi\`u precise, ottenendo l'ipotesi
        \begin{equation*}
            \Psi(x) =
            \begin{cases}
                x \in \mathcal{L} \implies \hat{\delta}(q_0, x) = q_2, \\
                x \notin \mathcal{L} \land \text{``$x$ non contiene $1$''}
                    \implies \hat{\delta}(q_0, x) = q_0, \\
                x \notin \mathcal{L} \land \text{``$x$ contiene un solo $1$''}
                    \implies \hat{\delta}(q_0, x) = q_1.
            \end{cases}
        \end{equation*}
        \`E lecito farlo perch\'e ogni nuova condizione implica quella di partenza, dato che
        $q_0$ e $q_1$ sono diversi da $q_2$, e la base verifica gi\`a gli stati esatti. In
        cambio, nel passo induttivo va dimostrata per $x$ l'intera $\Psi(x)$.

        Si dimostri allora il caso base per $\Psi$:
        %Dimostrazione del caso base con Psi(x)%
        \begin{equation*}
            \left.
            \begin{alignedat}{3}
                |x| = 0 \quad & x = \varepsilon \notin \mathcal{L} & \qquad
                    & \hat{\delta}(q_0, \varepsilon) = q_0
                    &\quad \text{(caso 2 di $\Psi$)} \\
                |x| = 1 \quad & x = 0 \notin \mathcal{L} &
                    & \hat{\delta}(q_0, 0) = q_0
                    &\quad \text{(caso 2 di $\Psi$)} \\
                & x = 1 \notin \mathcal{L} &
                    & \hat{\delta}(q_0, 1) = q_1
                    &\quad \text{(caso 3 di $\Psi$)} \\
                |x| = 2 \quad & x = 00 \notin \mathcal{L} &
                    & \hat{\delta}(q_0, 00) = q_0
                    &\quad \text{(caso 2 di $\Psi$)} \\
                & x = 01 \notin \mathcal{L} &
                    & \hat{\delta}(q_0, 01) = q_1
                    &\quad \text{(caso 3 di $\Psi$)} \\
                & x = 10 \notin \mathcal{L} &
                    & \hat{\delta}(q_0, 10) = q_1
                    &\quad \text{(caso 3 di $\Psi$)} \\
                & x = 11 \in \mathcal{L} &
                    & \hat{\delta}(q_0, 11) = q_2
                    &\quad \text{(caso 1 di $\Psi$)}
            \end{alignedat}
            \right\}
            \quad \text{\textsc{Base}}
        \end{equation*}

        \textsc{Passo induttivo}. Per verificare $\Psi(x)$ su tutte le stringhe di
        $\Sigma^*$ si applica l'ipotesi induttiva
        \begin{equation*}
            \forall x \in \Sigma^*.\; \bigl(|x| \leq n \implies \Psi(x)\bigr)
        \end{equation*}
        Sia $x$ una stringa
        con $|x| = n + 1$: allora $x = x'0$ oppure $x = x'1$, con $|x'| = n$, e a $x'$ si
        pu\`o applicare l'ipotesi induttiva. Si distingue un caso per ciascuna condizione
        di $\Psi(x)$. In ogni caso si ragiona prima sul linguaggio, contando gli $1$ di $x'$;
        poi si applica l'ipotesi induttiva a $x'$ e infine si legge l'ultimo simbolo. Nei
        calcoli, \emph{def} indica la definizione di $\hat{\delta}$ e \emph{Hp} il punto in
        cui si sostituisce lo stato dato dall'ipotesi induttiva.
        \begin{enumerate}
            \item $x \in \mathcal{L}$, quindi $x$ contiene almeno due $1$.
                \begin{enumerate}[label=\arabic{enumi}\alph*.]
                    \item $x = x'0$. Se $x'$ contenesse un solo $1$ (o nessuno), anche
                        $x = x'0$ ne conterrebbe meno di due e quindi $x \notin \mathcal{L}$:
                        assurdo. Quindi $x'$ contiene almeno due $1$, cio\`e
                        $x' \in \mathcal{L}$ e, per ipotesi induttiva,
                        $\hat{\delta}(q_0, x') = q_2$. Allora
                        \begin{equation*}
                            \hat{\delta}(q_0, x) = \hat{\delta}(q_0, x'0)
                            \overset{\text{def}}{=} \delta(\hat{\delta}(q_0, x'), 0)
                            \overset{\text{Hp}}{=} \delta(q_2, 0) = q_2 \in \mathcal{F}.
                        \end{equation*}
                    \item $x = x'1$. Se $x'$ non contenesse $1$, $x = x'1$ ne conterrebbe
                        uno solo e quindi $x \notin \mathcal{L}$: assurdo. Quindi $x'$ contiene
                        almeno due $1$, e quindi $x' \in \mathcal{L}$, oppure ne contiene
                        esattamente uno, e quindi $x' \notin \mathcal{L}$.

                        Se $x' \in \mathcal{L}$, per ipotesi induttiva
                        $\hat{\delta}(q_0, x') = q_2$, quindi
                        \begin{equation*}
                            \hat{\delta}(q_0, x) = \hat{\delta}(q_0, x'1)
                            \overset{\text{def}}{=} \delta(\hat{\delta}(q_0, x'), 1)
                            \overset{\text{Hp}}{=} \delta(q_2, 1) = q_2 \in \mathcal{F}.
                        \end{equation*}

                        Se $x' \notin \mathcal{L}$, $x'$ contiene un solo $1$ e, per ipotesi
                        induttiva, $\hat{\delta}(q_0, x') = q_1$: \`e qui che serve lo
                        scorporo, perch\'e altrimenti si saprebbe solo che
                        $\hat{\delta}(q_0, x') \neq q_2$. Quindi
                        \begin{equation*}
                            \hat{\delta}(q_0, x) = \hat{\delta}(q_0, x'1)
                            \overset{\text{def}}{=} \delta(\hat{\delta}(q_0, x'), 1)
                            \overset{\text{Hp}}{=} \delta(q_1, 1) = q_2 \in \mathcal{F}.
                        \end{equation*}
                \end{enumerate}
            \item $x \notin \mathcal{L}$ e $x$ non contiene $1$. Il caso $x = x'1$ \`e
                assurdo, perch\'e $x$ conterrebbe un $1$, quindi $x = x'0$. Se $x'$
                contenesse almeno un $1$, lo conterrebbe anche $x$: assurdo. Perci\`o $x'$ non
                contiene $1$: $x' \notin \mathcal{L}$ e, per ipotesi induttiva,
                $\hat{\delta}(q_0, x') = q_0$. Allora
                \begin{equation*}
                    \hat{\delta}(q_0, x) = \hat{\delta}(q_0, x'0)
                    \overset{\text{def}}{=} \delta(\hat{\delta}(q_0, x'), 0)
                    \overset{\text{Hp}}{=} \delta(q_0, 0) = q_0.
                \end{equation*}
            \item $x \notin \mathcal{L}$ e $x$ contiene un solo $1$.
                \begin{enumerate}[label=\arabic{enumi}\alph*.]
                    \item $x = x'1$. Se $x'$ contenesse almeno un $1$, $x = x'1$ ne
                        conterrebbe almeno due e quindi $x \in \mathcal{L}$: assurdo. Perci\`o
                        $x'$ non contiene $1$: $x' \notin \mathcal{L}$ e, per ipotesi
                        induttiva, $\hat{\delta}(q_0, x') = q_0$. Allora
                        \begin{equation*}
                            \hat{\delta}(q_0, x) = \hat{\delta}(q_0, x'1)
                            \overset{\text{def}}{=} \delta(\hat{\delta}(q_0, x'), 1)
                            \overset{\text{Hp}}{=} \delta(q_0, 1) = q_1.
                        \end{equation*}
                    \item $x = x'0$. Se $x'$ non contenesse $1$, nemmeno $x = x'0$ ne
                        conterrebbe; se ne contenesse almeno due, $x \in \mathcal{L}$. In
                        entrambi i casi \`e assurdo, perci\`o $x'$ contiene un solo $1$:
                        $x' \notin \mathcal{L}$ e, per ipotesi induttiva,
                        $\hat{\delta}(q_0, x') = q_1$. Allora
                        \begin{equation*}
                            \hat{\delta}(q_0, x) = \hat{\delta}(q_0, x'0)
                            \overset{\text{def}}{=} \delta(\hat{\delta}(q_0, x'), 0)
                            \overset{\text{Hp}}{=} \delta(q_1, 0) = q_1.
                        \end{equation*}
                \end{enumerate}
        \end{enumerate}
        In tutti i casi lo stato raggiunto da $x$ \`e quello previsto da $\Psi(x)$, quindi
        $\Psi(x)$ vale anche per $|x| = n + 1$. Per induzione vale per ogni $x \in \Sigma^*$:
        in particolare $\hat{\delta}(q_0, x) \in \mathcal{F}$ se e solo se $x \in \mathcal{L}$,
        ovvero $\mathcal{L}(\mathcal{M}) = \mathcal{L}$.
    \end{esempio}

    \section{Automi a stati finiti non deterministici}
    In un \emph{automa a stati finiti non deterministico} (NFA) l'unico elemento che
    cambia rispetto ai DFA \`e la funzione di transizione. L'obiettivo \`e stabilire se il
    non determinismo aumenti la potenza di calcolo.

    Anche in un NFA l'unica forma di memoria \`e lo stato in cui l'automa si trova, ma la
    funzione di transizione \`e non deterministica: da uno stato, leggendo un simbolo,
    l'automa pu\`o eseguire zero, una o pi\`u transizioni, cio\`e raggiungere un
    qualsiasi insieme di stati, anche vuoto.

    \begin{definition}[Automa non deterministico]
        Un \emph{automa a stati finiti non deterministico} (NFA) \`e una quintupla
        $\langle \mathcal{Q}, \Sigma, \delta, q_0, \mathcal{F} \rangle$, dove:
        \begin{itemize}
            \item $\mathcal{Q}$ \`e un insieme finito di stati;
            \item $\Sigma$ \`e un alfabeto;
            \item $\delta \colon \mathcal{Q} \times \Sigma \to \mathcal{P}(\mathcal{Q})$ \`e
                la \emph{funzione di transizione}, che associa a ogni stato e a ogni
                simbolo l'insieme degli stati raggiungibili;
            \item $q_0 \in \mathcal{Q}$ \`e lo stato \emph{iniziale};
            \item $\mathcal{F} \subseteq \mathcal{Q}$ \`e l'insieme degli stati
                \emph{finali}.
        \end{itemize}
    \end{definition}

    Come per i DFA, la funzione $\delta$ va estesa alle stringhe. La funzione
    $\hat{\delta}$, che si ottiene in modo univoco da $\delta$, \`e definita come segue:
    \begin{equation*}
        \hat{\delta} \colon \mathcal{Q} \times \Sigma^* \to \mathcal{P}(\mathcal{Q}),
        \qquad
        \begin{cases}
            \hat{\delta}(q, \varepsilon) = \{q\}, \\
            \hat{\delta}(q, wa) = \bigcup_{p \in \hat{\delta}(q, w)} \delta(p, a).
        \end{cases}
    \end{equation*}
    Ad esempio, si supponga che, leggendo $w$ a partire da $q$, l'automa possa arrivare
    negli stati $p_1$, $p_2$ e $p_3$, ovvero $\hat{\delta}(q, w) = \{p_1, p_2, p_3\}$.
    Leggendo poi il simbolo $a$, da ciascuno di essi si raggiunge l'insieme
    $\delta(p_i, a)$, e $\hat{\delta}(q, wa)$ \`e l'unione di questi insiemi, come mostrato
    nella Figura~\ref{fig:nfa-delta}:
    \begin{equation*}
        \hat{\delta}(q, wa) = \delta(p_1, a) \cup \delta(p_2, a) \cup \delta(p_3, a).
    \end{equation*}
    Nella figura $\delta(p_1, a) = \{r_1, r_2, r_3\}$, $\delta(p_2, a) = \{r_4, r_5, r_6\}$ e
    $\delta(p_3, a) = \{r_7, r_8\}$, quindi
    \begin{equation*}
        \hat{\delta}(q, wa) = \{r_1, r_2, r_3\} \cup \{r_4, r_5, r_6\} \cup \{r_7, r_8\}
        = \{r_1, \dots, r_8\}.
    \end{equation*}
    In generale gli insiemi $\delta(p_i, a)$ possono avere stati in comune, che
    nell'unione compaiono una volta sola. Leggendo $wa$ un simbolo alla volta, l'automa
    attraversa quindi gli insiemi di stati
    \begin{equation*}
        \hat{\delta}(q, \varepsilon) = \{q\}
        \;\xrightarrow{\;w\;}\;
        \hat{\delta}(q, w) = \{p_1, p_2, p_3\}
        \;\xrightarrow{\;a\;}\;
        \hat{\delta}(q, wa) = \{r_1, \dots, r_8\}.
    \end{equation*}

    \begin{figure}[H]
        \centering
        \begin{tikzpicture}[
            freccia/.style={-{Stealth}, thick},
            gruppo/.style={draw, dashed, rounded corners=6pt, inner sep=5pt}
        ]
            % Stati raggiunti da q leggendo w
            \node[punto, label=left:$q$] (q) at (0, 0) {};
            \foreach \i/\y in {1/2.0, 2/0, 3/-1.75}
                \node[punto, label=above:$p_\i$] (p\i) at (3, \y) {};
            \node[gruppo, fit=(p1)(p2)(p3), inner xsep=12pt, inner ysep=16pt,
                label=below:{$\hat{\delta}(q, w)$}] {};
            \draw[freccia] (q) to[bend left=15] node[above left] {$w$} (p1);
            \draw[freccia] (q) to (p2);
            \draw[freccia] (q) to[bend right=15] (p3);

            % Da ogni p_i, leggendo a, si raggiunge l'insieme delta(p_i, a)
            \foreach \n/\r/\y in {t11/1/2.55, t12/2/2.0, t13/3/1.45,
                    t21/4/0.55, t22/5/0, t23/6/-0.55, t31/7/-1.475, t32/8/-2.025}
                \node[inner sep=1.5pt] (\n) at (6.2, \y) {$r_\r$};
            \foreach \i/\k in {1/3, 2/3, 3/2} {
                \foreach \j in {1,...,\k}
                    \draw[freccia] (p\i) -- (t\i\j);
                \node[gruppo, fit=(t\i1)(t\i2)(t\i\k),
                    label={[label distance=16pt]right:{$\delta(p_\i, a)$}}] (D\i) {};
            }
            \node at (4.4, 2.5) {$a$};
            \node at (4.4, 0.5) {$a$};
            \node at (4.4, -1.3) {$a$};

            % L'unione degli insiemi delta(p_i, a)
            \node[draw, thick, evidenza, rounded corners=10pt, inner sep=9pt,
                fit=(D1)(D2)(D3), label=above:{$\hat{\delta}(q, wa)$}] {};
        \end{tikzpicture}
        \caption{Calcolo di $\hat{\delta}(q, wa)$ in un NFA: l'unione degli insiemi
            $\delta(p_i, a)$, per ogni $p_i \in \hat{\delta}(q, w)$.}
        \label{fig:nfa-delta}
    \end{figure}

    Il linguaggio accettato da un NFA
    $\mathcal{N} = \langle \mathcal{Q}, \Sigma, \delta, q_0, \mathcal{F} \rangle$ \`e
    l'insieme delle stringhe per cui almeno uno degli stati raggiunti \`e finale:
    \begin{equation*}
        \mathcal{L}(\mathcal{N})
        = \{x \in \Sigma^* \mid \hat{\delta}(q_0, x) \cap \mathcal{F} \neq \emptyset\}.
    \end{equation*}

    \begin{theorem}[di Rabin e Scott, 1959]
        \label{thm:rabin-scott}
        \leavevmode
        \begin{enumerate}[label=\textup{(\alph*)}]
            \item Sia $\mathcal{M} = \langle \mathcal{Q}, \Sigma, \delta, q_0, \mathcal{F}
                \rangle$ un DFA. Allora esiste un NFA $\mathcal{N} = \langle \mathcal{Q}',
                \Sigma, \delta', q_0, \mathcal{F} \rangle$ che riconosce lo stesso
                linguaggio, ovvero $\mathcal{L}(\mathcal{M}) = \mathcal{L}(\mathcal{N})$.
            \item Sia $\mathcal{N} = \langle \mathcal{Q}, \Sigma, \delta, q_0, \mathcal{F}
                \rangle$ un NFA. Allora esiste un DFA $\mathcal{M} = \langle \mathcal{Q}',
                \Sigma, \delta', q_0', \mathcal{F}' \rangle$ che riconosce lo stesso
                linguaggio, ovvero $\mathcal{L}(\mathcal{N}) = \mathcal{L}(\mathcal{M})$.
        \end{enumerate}
    \end{theorem}

    \begin{proof}[Dimostrazione del Teorema~\ref{thm:rabin-scott}]
        L'obiettivo \`e dimostrare che NFA e DFA sono equivalenti.

        \textsc{Punto (a)}. Un DFA \`e un caso particolare di NFA, in cui ogni transizione
        porta in un solo stato. Si definisce quindi $\mathcal{N}$ con
        $\mathcal{Q}' = \mathcal{Q}$, gli stessi $q_0$ e $\mathcal{F}$, e
        \begin{equation*}
            \delta'(q, a) = \{\delta(q, a)\}.
        \end{equation*}
        Per induzione su $|x|$ si verifica che $\hat{\delta}'(q_0, x) =
        \{\hat{\delta}(q_0, x)\}$. Per $x = \varepsilon$ entrambi i membri valgono $\{q_0\}$;
        per $x = wa$, con $|w| = n$, l'ipotesi induttiva dice che
        $\hat{\delta}'(q_0, w) = \{\hat{\delta}(q_0, w)\}$, quindi
        \begin{align*}
            \hat{\delta}'(q_0, wa)
                &= \bigcup_{p \in \hat{\delta}'(q_0, w)} \delta'(p, a)
                    && \text{(definizione di $\hat{\delta}'$ per gli NFA)} \\
                &= \bigcup_{p \in \{\hat{\delta}(q_0, w)\}} \delta'(p, a)
                    && \text{(ipotesi induttiva)} \\
                &= \delta'(\hat{\delta}(q_0, w), a)
                    && \text{(unione di un solo stato)} \\
                &= \{\delta(\hat{\delta}(q_0, w), a)\}
                    && \text{(definizione di $\delta'$)} \\
                &= \{\hat{\delta}(q_0, wa)\}
                    && \text{(definizione di $\hat{\delta}$)}.
        \end{align*}
        Quindi $\hat{\delta}'(q_0, x) \cap \mathcal{F} \neq \emptyset$ se e solo se
        $\hat{\delta}(q_0, x) \in \mathcal{F}$, ovvero $\mathcal{L}(\mathcal{N}) =
        \mathcal{L}(\mathcal{M})$.

        \textsc{Punto (b)}. L'idea \`e che ogni stato di $\mathcal{M}$ rappresenti
        l'insieme degli stati in cui $\mathcal{N}$ pu\`o trovarsi dopo aver letto una
        stringa. Si definisce quindi $\mathcal{M}$ come segue:
        \begin{itemize}
            \item $\mathcal{Q}' = \mathcal{P}(\mathcal{Q})$: gli stati di $\mathcal{M}$ sono
                i sottoinsiemi di $\mathcal{Q}$, che sono $2^{|\mathcal{Q}|}$ e quindi in
                numero finito\footnote{Sarebbe pi\`u preciso definire
                $\mathcal{Q}' = \{q_1, \dots, q_{2^{|\mathcal{Q}|}}\}$ e stabilire una
                corrispondenza biunivoca tra questi stati e gli elementi di
                $\mathcal{P}(\mathcal{Q})$; usare direttamente i sottoinsiemi come stati
                rende per\`o la dimostrazione molto pi\`u snella.};
            \item $q_0' = \{q_0\}$: all'inizio $\mathcal{N}$ si trova solo in $q_0$;
            \item $\delta' \colon \mathcal{Q}' \times \Sigma \to \mathcal{Q}'$, con
                $\delta'(S, a) = \bigcup_{p \in S} \delta(p, a)$ per ogni
                $S \in \mathcal{Q}'$ e $a \in \Sigma$: da $S$, leggendo $a$, si
                raggiungono tutti gli stati raggiungibili da almeno uno degli stati di $S$.
                Il risultato \`e un sottoinsieme di $\mathcal{Q}$, cio\`e un solo stato di
                $\mathcal{Q}'$, quindi $\mathcal{M}$ \`e deterministico;
            \item $\mathcal{F}' = \{S \in \mathcal{Q}' \mid S \cap \mathcal{F} \neq
                \emptyset\}$: un insieme di stati \`e finale se contiene almeno uno stato
                finale di $\mathcal{N}$.
        \end{itemize}

        Si dimostra per induzione su $|x|$ che, per ogni $x \in \Sigma^*$,
        \begin{equation*}
            \hat{\delta}'(q_0', x) = \hat{\delta}(q_0, x),
        \end{equation*}
        cio\`e che dopo aver letto $x$ l'automa $\mathcal{M}$ si trova nello stato che
        corrisponde esattamente all'insieme degli stati raggiunti da $\mathcal{N}$.

        \textsc{Base}. Per $|x| = 0$ si ha $x = \varepsilon$ e
        \begin{equation*}
            \hat{\delta}'(q_0', \varepsilon) = q_0' = \{q_0\}
            = \hat{\delta}(q_0, \varepsilon).
        \end{equation*}

        \textsc{Passo induttivo}. Si suppone che l'uguaglianza valga per tutte le stringhe
        di lunghezza $n$ e sia $x = wa$, con $|w| = n$ e $a \in \Sigma$. Allora
        \begin{align*}
            \hat{\delta}'(q_0', wa)
                &= \delta'(\hat{\delta}'(q_0', w), a)
                    && \text{(definizione di $\hat{\delta}'$)} \\
                &= \delta'(\hat{\delta}(q_0, w), a)
                    && \text{(ipotesi induttiva)} \\
                &= \bigcup_{p \in \hat{\delta}(q_0, w)} \delta(p, a)
                    && \text{(definizione di $\delta'$)} \\
                &= \hat{\delta}(q_0, wa)
                    && \text{(definizione di $\hat{\delta}$ per gli NFA)}.
        \end{align*}

        Infine, per ogni $x \in \Sigma^*$,
        \begin{align*}
            x \in \mathcal{L}(\mathcal{M})
                &\iff \hat{\delta}'(q_0', x) \in \mathcal{F}' \\
                &\iff \hat{\delta}(q_0, x) \in \mathcal{F}'
                    && \text{(uguaglianza appena dimostrata)} \\
                &\iff \hat{\delta}(q_0, x) \cap \mathcal{F} \neq \emptyset
                    && \text{(definizione di $\mathcal{F}'$)} \\
                &\iff x \in \mathcal{L}(\mathcal{N}),
        \end{align*}
        quindi $\mathcal{L}(\mathcal{M}) = \mathcal{L}(\mathcal{N})$.
    \end{proof}

    DFA e NFA riconoscono quindi la stessa classe di linguaggi, i linguaggi regolari: il
    non determinismo non aumenta la potenza di calcolo. In altre parole, la possibilit\`a
    di raggiungere pi\`u stati diversi leggendo lo stesso simbolo \emph{non} cambia la
    classe dei linguaggi riconoscibili.

    \begin{esempio}
        Si consideri l'NFA della Figura~\ref{fig:nfa-esempio}, con stato iniziale $q_0$ e
        unico stato finale $q_2$.

        \begin{figure}[H]
            \centering
            \begin{tikzpicture}[automa]
                \node[state, initial] (q0) at (0, 0) {$q_0$};
                \node[state] (q1) at (4, 0) {$q_1$};
                \node[state, accepting] (q2) at (2, 2.6) {$q_2$};
                \path (q0) edge[loop below] node {$0, 1$} ()
                    (q1) edge[loop below] node {$0$} ()
                    (q2) edge[loop above] node {$0, 1$} ()
                    (q0) edge[bend left=15] node {$1$} (q1)
                    (q1) edge[bend left=15] node {$1$} (q0)
                    (q1) edge[bend left=15] node[pos=0.65] {$1$} (q2)
                    (q2) edge[bend left=15] node {$0, 1$} (q1)
                    (q2) edge node[swap] {$1$} (q0);

                % Matrice di transizione, a destra del diagramma
                \node at (8, 1.2) {%
                    \begin{tabular}{ccc}
                        \toprule
                        $\delta$ & $0$ & $1$ \\
                        \midrule
                        $q_0$ & $\{q_0\}$ & $\{q_0, q_1\}$ \\
                        $q_1$ & $\{q_1\}$ & $\{q_0, q_2\}$ \\
                        $q_2$ & $\{q_1, q_2\}$ & $\{q_0, q_1, q_2\}$ \\
                        \bottomrule
                    \end{tabular}};
            \end{tikzpicture}
            \caption{Un NFA e la sua matrice di transizione.}
            \label{fig:nfa-esempio}
        \end{figure}

        Si applica la costruzione del punto (b) della dimostrazione del
        Teorema~\ref{thm:rabin-scott}, che \`e del tutto meccanica:
        ogni sottoinsieme di $\mathcal{Q} = \{q_0, q_1, q_2\}$ \`e uno stato del DFA, e la
        sua riga nella matrice di transizione si ottiene unendo le righe dell'NFA dei suoi
        elementi. Ad esempio,
        \begin{equation*}
            \delta'(\{q_0, q_1\}, 1) = \delta(q_0, 1) \cup \delta(q_1, 1)
            = \{q_0, q_1\} \cup \{q_0, q_2\} = \{q_0, q_1, q_2\}.
        \end{equation*}
        Gli stati finali sono i sottoinsiemi che contengono $q_2$, l'unico stato finale
        dell'NFA. Si ottiene la Tabella~\ref{tab:sottoinsiemi}, in cui a ogni sottoinsieme
        non vuoto si d\`a un nome $q_i'$.

        \begin{table}[H]
            \centering
            \begin{tabular}{ccccc}
                \toprule
                \textsc{Stato} & \textsc{Sottoinsieme} & $0$ & $1$ & \textsc{Finale} \\
                \midrule
                & $\emptyset$ & $\emptyset$ & $\emptyset$ & \\
                $q_0'$ & $\{q_0\}$ & $\{q_0\}$ & $\{q_0, q_1\}$ & \\
                $q_1'$ & $\{q_1\}$ & $\{q_1\}$ & $\{q_0, q_2\}$ & \\
                $q_2'$ & $\{q_2\}$ & $\{q_1, q_2\}$ & $\{q_0, q_1, q_2\}$ & $\checkmark$ \\
                $q_3'$ & $\{q_0, q_1\}$ & $\{q_0, q_1\}$ & $\{q_0, q_1, q_2\}$ & \\
                $q_4'$ & $\{q_0, q_2\}$ & $\{q_0, q_1, q_2\}$ & $\{q_0, q_1, q_2\}$
                    & $\checkmark$ \\
                $q_5'$ & $\{q_1, q_2\}$ & $\{q_1, q_2\}$ & $\{q_0, q_1, q_2\}$
                    & $\checkmark$ \\
                $q_6'$ & $\{q_0, q_1, q_2\}$ & $\{q_0, q_1, q_2\}$ & $\{q_0, q_1, q_2\}$
                    & $\checkmark$ \\
                \bottomrule
            \end{tabular}
            \caption{Costruzione per sottoinsiemi del DFA equivalente all'NFA della
                Figura~\ref{fig:nfa-esempio}.}
            \label{tab:sottoinsiemi}
        \end{table}

        Sostituendo a ogni sottoinsieme il suo nome, e chiamando per semplicit\`a $q_i$
        lo stato $q_i'$, si ottiene il DFA della Figura~\ref{fig:dfa-sottoinsiemi}, in cui
        lo stato $\emptyset$ \`e omesso.

        \begin{figure}[H]
            \centering
            \begin{tikzpicture}[
                automa,
                gruppo/.style={draw, dashed, thin, -, rounded corners=12pt, inner sep=7pt}
            ]
                \node[state, initial] (q0) at (2.2, 3.4) {$q_0$};
                \node[state] (q3) at (2.2, 1.7) {$q_3$};
                \node[state, accepting] (q6) at (2.2, 0) {$q_6$};
                \node[state] (q1) at (5.0, 3.4) {$q_1$};
                \node[state, accepting] (q4) at (5.0, 1.7) {$q_4$};
                \node[state, accepting] (q2) at (-0.4, 2.4) {$q_2$};
                \node[state, accepting] (q5) at (-0.4, 0.6) {$q_5$};
                \path (q0) edge[loop above] node {$0$} ()
                    (q0) edge node {$1$} (q3)
                    (q3) edge[loop right] node {$0$} ()
                    (q3) edge node {$1$} (q6)
                    (q6) edge[loop below] node {$0, 1$} ()
                    (q1) edge[loop above] node {$0$} ()
                    (q1) edge node {$1$} (q4)
                    (q4) edge node {$0, 1$} (q6)
                    (q2) edge node[swap] {$0$} (q5)
                    (q2) edge node {$1$} (q6)
                    (q5) edge[loop below] node {$0$} ()
                    (q5) edge node[swap] {$1$} (q6);

                % Stati non raggiungibili da q0
                \node[gruppo, fit={(q1) (q4) (5.0, 5.0)}] {};
                \node[gruppo, fit={(q2) (q5) (-0.4, -1.0)}] {};

                % Matrice di transizione, a destra del diagramma
                \node at (8.6, 1.7) {%
                    \begin{tabular}{ccc}
                        \toprule
                        $\delta'$ & $0$ & $1$ \\
                        \midrule
                        $q_0$ & $q_0$ & $q_3$ \\
                        $q_1$ & $q_1$ & $q_4$ \\
                        $q_2$ & $q_5$ & $q_6$ \\
                        $q_3$ & $q_3$ & $q_6$ \\
                        $q_4$ & $q_6$ & $q_6$ \\
                        $q_5$ & $q_5$ & $q_6$ \\
                        $q_6$ & $q_6$ & $q_6$ \\
                        \bottomrule
                    \end{tabular}};
            \end{tikzpicture}
            \caption{Il DFA ottenuto dalla Tabella~\ref{tab:sottoinsiemi}; i riquadri
                tratteggiati racchiudono gli stati non raggiungibili da $q_0$.}
            \label{fig:dfa-sottoinsiemi}
        \end{figure}

        \`E interessante notare che alcuni stati non sono raggiungibili da $q_0$. Gli
        stati $q_1$ e $q_2$ non hanno archi entranti da altri stati, mentre $q_4$ e $q_5$
        ne hanno, ma solo da $q_1$ e da $q_2$. Quello che conta, infatti, non \`e avere
        archi entranti, ma l'esistenza di un cammino da $q_0$: gli unici stati
        raggiungibili sono $q_0$, $q_3$ e $q_6$.

        Gli stati non raggiungibili non cambiano il linguaggio riconosciuto e si possono
        eliminare, ottenendo il DFA minimo della Figura~\ref{fig:dfa-minimo}.\footnote{In
        questo caso basta eliminare gli stati non raggiungibili; in generale, per ottenere
        il DFA minimo pu\`o servire anche unire stati equivalenti.} \`E lo stesso automa
        dell'esempio sulle stringhe con almeno due $1$ (Figura~\ref{fig:dfa-due-uno}), con
        $q_3$ e $q_6$ al posto di $q_1$ e $q_2$: l'NFA di partenza riconosce quindi proprio
        quel linguaggio.

        \begin{figure}[H]
            \centering
            \begin{tikzpicture}[automa]
                \node[state, initial] (q0) {$q_0$};
                \node[state, right=of q0] (q3) {$q_3$};
                \node[state, accepting, right=of q3] (q6) {$q_6$};
                \path (q0) edge node {$1$} (q3)
                    (q3) edge node {$1$} (q6)
                    (q0) edge[loop above] node {$0$} ()
                    (q3) edge[loop above] node {$0$} ()
                    (q6) edge[loop above] node {$0, 1$} ();
            \end{tikzpicture}
            \caption{Il DFA minimo, ottenuto eliminando gli stati non raggiungibili.}
            \label{fig:dfa-minimo}
        \end{figure}
    \end{esempio}

    \section{Automi a stati finiti non deterministici con
        \texorpdfstring{$\varepsilon$}{ε}-transizioni}
    Ci si chiede ora se la possibilit\`a di cambiare stato \emph{senza} leggere simboli
    cambi la classe dei linguaggi riconoscibili.

    \begin{definition}[Automa non deterministico con $\varepsilon$-transizioni]
        Un \emph{automa a stati finiti non deterministico con $\varepsilon$-transizioni}
        ($\varepsilon$-NFA) \`e una quintupla
        $\langle \mathcal{Q}, \Sigma, \delta, q_0, \mathcal{F} \rangle$ che differisce da un
        NFA solo per la funzione di transizione:
        \begin{itemize}
            \item $\mathcal{Q}$ \`e un insieme finito di stati;
            \item $\Sigma$ \`e un alfabeto;
            \item $\delta \colon \mathcal{Q} \times (\Sigma \cup \{\varepsilon\}) \to
                \mathcal{P}(\mathcal{Q})$ \`e la \emph{funzione di transizione}: come negli
                NFA, da uno stato si raggiunge un insieme di stati, ma lo si pu\`o fare anche
                senza leggere alcun simbolo, seguendo un arco etichettato con $\varepsilon$
                (\emph{$\varepsilon$-transizione});
            \item $q_0 \in \mathcal{Q}$ \`e lo stato \emph{iniziale};
            \item $\mathcal{F} \subseteq \mathcal{Q}$ \`e l'insieme degli stati
                \emph{finali}.
        \end{itemize}
    \end{definition}

    Prima di estendere $\delta$ alle stringhe \`e necessario definire la funzione
    $\eclosure \colon \mathcal{Q} \to \mathcal{P}(\mathcal{Q})$, che associa a ogni stato
    $q$ l'insieme degli stati raggiungibili da $q$ senza leggere simboli, ovvero seguendo
    zero o pi\`u archi etichettati con $\varepsilon$:
    \begin{equation*}
        \eclosure(q) = \{p \in \mathcal{Q} \mid p \text{ \`e raggiungibile da } q
        \text{ con zero o pi\`u } \varepsilon\text{-transizioni}\}.
    \end{equation*}
    Equivalentemente, $\eclosure(q)$ \`e il pi\`u piccolo insieme che contiene $q$ e che,
    per ogni suo stato $p$, contiene anche $\delta(p, \varepsilon)$. In particolare
    $q \in \eclosure(q)$, perch\'e $q$ si raggiunge con zero $\varepsilon$-transizioni, e
    $\delta(q, \varepsilon) \subseteq \eclosure(q)$. L'inclusione pu\`o essere stretta: se
    $q \xrightarrow{\varepsilon} p \xrightarrow{\varepsilon} r$, allora
    $r \in \eclosure(q)$, anche se $r$ non appartiene necessariamente a
    $\delta(q, \varepsilon)$. Per un insieme di stati $S \subseteq \mathcal{Q}$ si pone
    \begin{equation*}
        \eclosure(S) = \bigcup_{s \in S} \eclosure(s).
    \end{equation*}

    Si pu\`o ora definire la funzione estesa $\hat{\delta}$:
    \begin{equation*}
        \hat{\delta} \colon \mathcal{Q} \times \Sigma^* \to \mathcal{P}(\mathcal{Q}),
        \qquad
        \begin{cases}
            \hat{\delta}(q, \varepsilon) = \eclosure(q), \\
            \hat{\delta}(q, xa) = \bigcup_{p \in \hat{\delta}(q, x)}
                \eclosure(\delta(p, a)).
        \end{cases}
    \end{equation*}
    La ricorsione avviene sulla stringa e non sugli stati: $\hat{\delta}(q, xa)$ richiama
    $\hat{\delta}(q, x)$ con lo stesso stato di partenza $q$ e la stringa privata
    dell'ultimo simbolo, fino al caso base $\hat{\delta}(q, \varepsilon)$.

    \begin{esempio}
        Si consideri l'automa della Figura~\ref{fig:enfa-esempio} e si calcoli
        $\hat{\delta}(q, a)$. Poich\'e $a = \varepsilon a$, si applica il secondo caso
        della definizione con $x = \varepsilon$.

        \begin{figure}[H]
            \centering
            \begin{tikzpicture}[automa, node distance=2.3cm]
                \node[state] (q) {$q$};
                \node[state, right=of q] (q1) {$q'$};
                \node[state, right=of q1] (q2) {$q''$};
                \node[state, right=of q2] (q3) {$q'''$};
                \path (q) edge node {$a$} (q1)
                    (q1) edge node {$\varepsilon$} (q2)
                    (q2) edge node {$\varepsilon$} (q3);
            \end{tikzpicture}
            \caption{Un automa con due $\varepsilon$-transizioni consecutive.}
            \label{fig:enfa-esempio}
        \end{figure}

        Da $q$ non escono archi etichettati con $\varepsilon$, quindi
        $\hat{\delta}(q, \varepsilon) = \eclosure(q) = \{q\}$ e l'unione \`e su un solo
        stato. Leggendo $a$ si ottiene $\delta(q, a) = \{q'\}$, e infine si seguono tutte
        le catene di archi $\varepsilon$ a partire da $q'$:
        \begin{align*}
            \hat{\delta}(q, a)
                &= \bigcup_{p \in \hat{\delta}(q, \varepsilon)} \eclosure(\delta(p, a))
                    && \text{(definizione di $\hat{\delta}$)} \\
                &= \eclosure(\{q'\})
                    && \text{(estensione di $\eclosure$ agli insiemi)} \\
                &= \{q', q'', q'''\}.
        \end{align*}
        Per calcolare la chiusura si parte da $\{q'\}$ e si aggiungono gli stati
        raggiungibili con un arco $\varepsilon$ da quelli gi\`a presenti, finch\'e non se
        ne trovano di nuovi: prima $q''$, poi $q'''$. Da $q'''$ non escono archi
        $\varepsilon$, quindi ci si ferma. A differenza di DFA e NFA, $\hat{\delta}(q, a)$
        pu\`o quindi essere diverso da $\delta(q, a) = \{q'\}$.

        Le $\varepsilon$-transizioni contano anche prima di leggere un simbolo: in un
        automa $q \xrightarrow{\varepsilon} r \xrightarrow{a} s$ si ha
        $\delta(q, a) = \emptyset$, ma $\hat{\delta}(q, \varepsilon) = \{q, r\}$ e quindi
        $s \in \hat{\delta}(q, a)$.
    \end{esempio}

    \begin{esempio}
        Si calcoli $\hat{\delta}(q, ab)$ per l'automa della
        Figura~\ref{fig:enfa-stack}, svolgendo per intero la ricorsione. Le transizioni
        dell'automa sono $\delta(q, a) = \{q'\}$, $\delta(q', \varepsilon) = \{q''\}$,
        $\delta(q'', b) = \{s\}$ e $\delta(s, \varepsilon) = \{t\}$; tutte le altre danno
        $\emptyset$.

        \begin{figure}[H]
            \centering
            \begin{tikzpicture}[automa, node distance=1.8cm]
                \node[state] (q) {$q$};
                \node[state, right=of q] (q1) {$q'$};
                \node[state, right=of q1] (q2) {$q''$};
                \node[state, right=of q2] (s) {$s$};
                \node[state, right=of s] (t) {$t$};
                \path (q) edge node {$a$} (q1)
                    (q1) edge node {$\varepsilon$} (q2)
                    (q2) edge node {$b$} (s)
                    (s) edge node {$\varepsilon$} (t);
            \end{tikzpicture}
            \caption{Un automa in cui il simbolo $b$ si pu\`o leggere solo dopo una
                $\varepsilon$-transizione.}
            \label{fig:enfa-stack}
        \end{figure}

        \textsc{Discesa}. Ogni chiamata toglie l'ultimo simbolo della stringa, fino al caso
        base:
        \begin{alignat*}{2}
            &\hat{\delta}(q, ab) & \qquad & ab = a \cdot b \text{: si richiama }
                \hat{\delta}(q, a), \\
            &\qquad \hat{\delta}(q, a) & & a = \varepsilon \cdot a \text{: si richiama }
                \hat{\delta}(q, \varepsilon), \\
            &\qquad\qquad \hat{\delta}(q, \varepsilon) & & \text{caso base.}
        \end{alignat*}

        \textsc{Risalita}. Ogni livello legge il proprio simbolo e chiude con le
        $\varepsilon$-transizioni. Il caso base restituisce
        \begin{equation*}
            \hat{\delta}(q, \varepsilon) = \eclosure(q) = \{q\}.
        \end{equation*}
        Il livello intermedio legge $a$:
        \begin{align*}
            \hat{\delta}(q, a)
                &= \bigcup_{p \in \{q\}} \eclosure(\delta(p, a))
                    && \text{(unione su un solo stato)} \\
                &= \eclosure(\delta(q, a)) = \eclosure(\{q'\}) \\
                &= \{q', q''\}
                    && \text{($\delta(q', \varepsilon) = \{q''\}$, $\delta(q'',
                    \varepsilon) = \emptyset$)}.
        \end{align*}
        Il livello pi\`u esterno legge $b$, partendo dall'insieme appena restituito:
        \begin{align*}
            \hat{\delta}(q, ab)
                &= \bigcup_{p \in \{q', q''\}} \eclosure(\delta(p, b)) \\
                &= \eclosure(\delta(q', b)) \cup \eclosure(\delta(q'', b))
                    && \text{(unione su due stati)} \\
                &= \eclosure(\emptyset) \cup \eclosure(\{s\}) \\
                &= \emptyset \cup \{s, t\} = \{s, t\}
                    && \text{($\delta(s, \varepsilon) = \{t\}$, $\delta(t,
                    \varepsilon) = \emptyset$)}.
        \end{align*}

        A ogni livello l'insieme su cui si fa l'unione \`e il valore restituito dal livello
        sottostante: $\{q\}$ per la lettura di $a$, $\{q', q''\}$ per la lettura di $b$.
        Il ramo $p = q'$ non contribuisce, perch\'e da $q'$ non si legge $b$; \`e il ramo
        $p = q''$, raggiunto solo grazie alla $\varepsilon$-closure applicata dopo $a$, a
        portare in $s$.
    \end{esempio}

    \begin{theorem}[Equivalenza tra NFA ed $\varepsilon$-NFA]
        \label{thm:equiv-nfa-enfa}
        \leavevmode
        \begin{enumerate}[label=\textup{(\alph*)}]
            \item Sia $\mathcal{N}_\varepsilon = \langle \mathcal{Q}, \Sigma, \delta, q_0,
                \mathcal{F} \rangle$ un $\varepsilon$-NFA. Allora esiste un NFA
                $\mathcal{N} = \langle \mathcal{Q}', \Sigma, \delta', q_0', \mathcal{F}'
                \rangle$ tale che $\mathcal{L}(\mathcal{N}_\varepsilon) =
                \mathcal{L}(\mathcal{N})$.
            \item Sia $\mathcal{N} = \langle \mathcal{Q}, \Sigma, \delta, q_0, \mathcal{F}
                \rangle$ un NFA. Allora esiste un $\varepsilon$-NFA $\mathcal{N}_\varepsilon
                = \langle \mathcal{Q}', \Sigma, \delta', q_0', \mathcal{F}' \rangle$ tale
                che $\mathcal{L}(\mathcal{N}) = \mathcal{L}(\mathcal{N}_\varepsilon)$.
        \end{enumerate}
    \end{theorem}

    \begin{proof}[Dimostrazione del Teorema~\ref{thm:equiv-nfa-enfa}]
        Anche in questo caso si deve dimostrare un'equivalenza, quindi servono entrambi i
        versi.

        \textsc{Punto (a)}. La dimostrazione si articola in tre passi:
        \begin{enumerate}
            \item costruire un NFA $\mathcal{N}$ a partire da $\mathcal{N}_\varepsilon$;
            \item dimostrare che, dopo aver letto una stringa non vuota, $\mathcal{N}$ e
                $\mathcal{N}_\varepsilon$ si trovano nello stesso insieme di stati;
            \item dedurne che $\mathcal{N}$ e $\mathcal{N}_\varepsilon$ accettano le
                stesse stringhe.
        \end{enumerate}

        \emph{Passo 1: costruzione di $\mathcal{N}$}. Gli stati, l'alfabeto e lo stato
        iniziale restano quelli di $\mathcal{N}_\varepsilon$, cio\`e
        $\mathcal{Q}' = \mathcal{Q}$ e $q_0' = q_0$; cambiano solo la funzione di
        transizione e gli stati finali.
        \begin{itemize}
            \item La nuova funzione di transizione ingloba le $\varepsilon$-transizioni
                dell'automa originale, sia quelle che precedono sia quelle che seguono la
                lettura del simbolo:
                \begin{equation*}
                    \delta'(q, a) \coloneqq \hat{\delta}(q, a)
                    = \bigcup_{p \in \eclosure(q)} \eclosure(\delta(p, a)).
                \end{equation*}
            \item Lo stato iniziale diventa finale se da esso si pu\`o raggiungere uno
                stato finale con sole $\varepsilon$-transizioni, perch\'e in questo caso
                $\mathcal{N}_\varepsilon$ accetta $\varepsilon$:
                \begin{equation*}
                    \mathcal{F}' =
                    \begin{cases}
                        \mathcal{F} \cup \{q_0\} & \text{se } \eclosure(q_0) \cap
                            \mathcal{F} \neq \emptyset, \\
                        \mathcal{F} & \text{altrimenti.}
                    \end{cases}
                \end{equation*}
        \end{itemize}

        \emph{Passo 2: stessi stati raggiunti}. Si dimostra per induzione su $|x|$
        che\footnote{La dimostrazione per induzione di questo passo non \`e stata svolta a
        lezione.}
        \begin{equation*}
            \hat{\delta}'(q_0, x) = \hat{\delta}(q_0, x)
            \qquad \text{per ogni } x \in \Sigma^* \text{ con } x \neq \varepsilon.
        \end{equation*}
        L'ipotesi $x \neq \varepsilon$ \`e necessaria, perch\'e per la stringa vuota
        l'uguaglianza non vale: $\hat{\delta}'(q_0, \varepsilon) = \{q_0\}$, mentre
        $\hat{\delta}(q_0, \varepsilon) = \eclosure(q_0)$.

        \emph{Passo 3: stesse stringhe accettate}. Bisogna dimostrare che, per ogni
        $x \in \Sigma^*$,\footnote{La dimostrazione di questo passo non \`e stata svolta
        a lezione.}
        \begin{equation*}
            \hat{\delta}'(q_0, x) \cap \mathcal{F}' \neq \emptyset \iff
            \hat{\delta}(q_0, x) \cap \mathcal{F} \neq \emptyset.
        \end{equation*}
        Questo passo usa il passo~2, che va quindi dimostrato prima. Poich\'e il passo~2
        vale solo per le stringhe non vuote, si distinguono due casi.

        \textit{Caso $x = \varepsilon$}. $\mathcal{N}$ accetta $\varepsilon$ se e solo
        se $q_0 \in \mathcal{F}'$, mentre $\mathcal{N}_\varepsilon$ accetta
        $\varepsilon$ se e solo se $\eclosure(q_0) \cap \mathcal{F} \neq \emptyset$.
        Basta quindi dimostrare che
        \begin{equation*}
            q_0 \in \mathcal{F}' \iff \eclosure(q_0) \cap \mathcal{F} \neq \emptyset,
        \end{equation*}
        che segue dalla definizione di $\mathcal{F}'$.
        \begin{itemize}[leftmargin=3em]
            \item[($\Leftarrow$)] Se $\eclosure(q_0) \cap \mathcal{F} \neq \emptyset$, si
                applica il primo caso della definizione: $\mathcal{F}' = \mathcal{F} \cup
                \{q_0\}$, quindi $q_0 \in \mathcal{F}'$.
            \item[($\Rightarrow$)] Se $q_0 \in \mathcal{F}'$, ci sono due possibilit\`a:
                \begin{itemize}
                    \item \`e stato applicato il primo caso della definizione, perch\'e
                        $\eclosure(q_0) \cap \mathcal{F} \neq \emptyset$, e la condizione
                        vale gi\`a;
                    \item \`e stato applicato il secondo caso, $\mathcal{F}' =
                        \mathcal{F}$: allora $q_0 \in \mathcal{F}$, cio\`e $q_0$ era gi\`a
                        finale. Poich\'e $q_0 \in \eclosure(q_0)$ sempre, dato che con zero
                        $\varepsilon$-transizioni si resta in $q_0$, si ha
                        $q_0 \in \eclosure(q_0) \cap \mathcal{F}$, che quindi non \`e
                        vuoto.
                \end{itemize}
        \end{itemize}

        \textit{Caso $x \neq \varepsilon$}. Per il passo~2 si pu\`o porre
        $S = \hat{\delta}'(q_0, x) = \hat{\delta}(q_0, x)$, e bisogna dimostrare che
        \begin{equation*}
            S \cap \mathcal{F}' \neq \emptyset \iff S \cap \mathcal{F} \neq \emptyset.
        \end{equation*}
        \begin{itemize}[leftmargin=3em]
            \item[($\Leftarrow$)] $\mathcal{F}'$ \`e un sovrainsieme di $\mathcal{F}$,
                quindi se esiste $f \in S \cap \mathcal{F}$ allora
                $f \in S \cap \mathcal{F}'$.
            \item[($\Rightarrow$)] Se $S \cap \mathcal{F}'$ contiene uno stato diverso da
                $q_0$, questo appartiene anche a $\mathcal{F}$ e non c'\`e altro da
                dimostrare. Resta il caso in cui l'unico stato in comune sia $q_0$, con
                $q_0 \notin \mathcal{F}$. Allora valgono sia $q_0 \in \mathcal{F}'$ sia
                $q_0 \in S$, e da ciascuna delle due condizioni segue un fatto:
                \begin{itemize}
                    \item da $q_0 \in \mathcal{F}'$, per la definizione di
                        $\mathcal{F}'$, esiste uno stato finale
                        $f \in \eclosure(q_0) \cap \mathcal{F}$, cio\`e uno stato finale
                        raggiungibile da $q_0$ con sole $\varepsilon$-transizioni;
                    \item da $q_0 \in S$ segue $\eclosure(q_0) \subseteq S$. Infatti
                        $q_0 \in S$ significa che, dopo aver letto $x$, l'automa pu\`o
                        trovarsi in $q_0$. Poich\'e $x \neq \varepsilon$, l'ultimo passo
                        di $\hat{\delta}$ legge un simbolo e termina con una
                        $\varepsilon$-closure, che da $q_0$ segue tutte le
                        $\varepsilon$-transizioni possibili e non si ferma finch\'e trova
                        stati nuovi: tutti gli stati raggiungibili da $q_0$ con sole
                        $\varepsilon$-transizioni entrano quindi in $S$.
                \end{itemize}
                Messi insieme, i due fatti danno la tesi:
                \begin{equation*}
                    q_0 \in S \cap \mathcal{F}'
                    \implies
                    \bigl(\exists f \in \eclosure(q_0) \cap \mathcal{F}\bigr) \land
                    \bigl(\eclosure(q_0) \subseteq S\bigr)
                    \implies
                    f \in S \cap \mathcal{F}.
                \end{equation*}
        \end{itemize}
        In entrambi i casi $\mathcal{N}$ e $\mathcal{N}_\varepsilon$ accettano le stesse
        stringhe, quindi $\mathcal{L}(\mathcal{N}) = \mathcal{L}(\mathcal{N}_\varepsilon)$.

        \textsc{Punto (b)}. Questo verso \`e banale: un NFA \`e un $\varepsilon$-NFA privo
        di archi etichettati con $\varepsilon$, ovvero con $\delta(q, \varepsilon) =
        \emptyset$ per ogni $q \in \mathcal{Q}$.
    \end{proof}

    \section{Operazioni tra linguaggi}
    \label{sec:operazioni-linguaggi}
    Siano $\mathcal{L}, \mathcal{L}_1, \mathcal{L}_2$ linguaggi sull'alfabeto $\Sigma$.
    Oltre alle operazioni insiemistiche della Sezione~\ref{sec:operazioni-insiemi}, come
    l'unione
    \begin{equation*}
        \mathcal{L}_1 \cup \mathcal{L}_2
        = \{x \mid x \in \mathcal{L}_1 \lor x \in \mathcal{L}_2\},
    \end{equation*}
    sui linguaggi sono definite due operazioni specifiche: la concatenazione e la chiusura
    di Kleene.

    \begin{definition}[Concatenazione]
        La \emph{concatenazione} di $\mathcal{L}_1$ e $\mathcal{L}_2$ \`e il linguaggio
        delle stringhe ottenute facendo seguire a una stringa di $\mathcal{L}_1$ una
        stringa di $\mathcal{L}_2$:
        \begin{equation*}
            \mathcal{L}_1 \cdot \mathcal{L}_2
            = \{x \cdot y \mid x \in \mathcal{L}_1 \land y \in \mathcal{L}_2\}.
        \end{equation*}
    \end{definition}
    Se uno dei due linguaggi contiene la stringa vuota, la concatenazione contiene l'altro.
    Ad esempio, se $\varepsilon \in \mathcal{L}_1$, ogni $y \in \mathcal{L}_2$ si scrive
    come $\varepsilon \cdot y$, quindi
    \begin{equation*}
        \varepsilon \in \mathcal{L}_1 \implies
            \mathcal{L}_2 \subseteq \mathcal{L}_1 \cdot \mathcal{L}_2,
        \qquad\qquad
        \varepsilon \in \mathcal{L}_2 \implies
            \mathcal{L}_1 \subseteq \mathcal{L}_1 \cdot \mathcal{L}_2.
    \end{equation*}

    \begin{definition}[Chiusura di Kleene]
        Le \emph{potenze} di $\mathcal{L}$ si definiscono per concatenazione iterata:
        \begin{equation*}
            \mathcal{L}^0 = \{\varepsilon\},
            \qquad\qquad
            \mathcal{L}^{i+1} = \mathcal{L} \cdot \mathcal{L}^i.
        \end{equation*}
        La \emph{chiusura di Kleene} $\mathcal{L}^*$ \`e l'unione di tutte le potenze di
        $\mathcal{L}$, mentre $\mathcal{L}^+$ esclude la potenza $\mathcal{L}^0$:
        \begin{equation*}
            \mathcal{L}^* = \bigcup_{i \geq 0} \mathcal{L}^i,
            \qquad\qquad
            \mathcal{L}^+ = \bigcup_{i \geq 1} \mathcal{L}^i.
        \end{equation*}
    \end{definition}

    \begin{esempio}
        Sia $\mathcal{L} = \{ab\}$. Allora
        \begin{align*}
            \mathcal{L}^0 &= \{\varepsilon\}, \\
            \mathcal{L}^1 &= \{ab\} \cdot \{\varepsilon\} = \{ab\}, \\
            \mathcal{L}^2 &= \{ab\} \cdot \mathcal{L}^1 = \{ab\} \cdot \{ab\} = \{abab\},
        \end{align*}
        e cos\`i via, quindi $\mathcal{L}^* = \{\varepsilon, ab, abab, ababab, \dots\}$.
        Con $\mathcal{L} = \{ab, cd\}$, invece,
        \begin{align*}
            \mathcal{L}^0 &= \{\varepsilon\}, \\
            \mathcal{L}^1 &= \{ab, cd\} \cdot \{\varepsilon\} = \{ab, cd\}, \\
            \mathcal{L}^2 &= \{ab, cd\} \cdot \mathcal{L}^1
                = \{abab, abcd, cdab, cdcd\},
        \end{align*}
        e cos\`i via.
    \end{esempio}

    \section{Espressioni regolari}
    \begin{definition}[Espressione regolare]
        Un'\emph{espressione regolare} (RE) \`e uno strumento formale utilizzato per
        descrivere e generare in modo compatto le stringhe che appartengono a un
        linguaggio regolare. Le espressioni regolari su un alfabeto $\Sigma$ sono definite
        per induzione:
        \begin{itemize}
            \item $\emptyset$ \`e un'espressione regolare, che denota l'insieme vuoto;
            \item $\varepsilon$ \`e un'espressione regolare, che denota l'insieme
                $\{\varepsilon\}$;
            \item per ogni simbolo $a \in \Sigma$, $a$ \`e un'espressione regolare, che
                denota l'insieme $\{a\}$;
            \item se $r$ e $s$ sono espressioni regolari che denotano rispettivamente gli
                insiemi $R$ e $S$, allora $(r + s)$, $(r \cdot s)$ e $(r^*)$ sono
                espressioni regolari, che denotano rispettivamente gli insiemi $R \cup S$,
                $R \cdot S$ e $R^*$.
        \end{itemize}
    \end{definition}

    Se $r$ \`e un'espressione regolare, si indica con $\mathcal{L}(r)$ il linguaggio
    denotato da $r$. Due espressioni regolari sono uguali se denotano lo stesso
    linguaggio,
    \begin{equation*}
        r = s \iff \mathcal{L}(r) = \mathcal{L}(s),
    \end{equation*}
    e tra espressioni regolari valgono uguaglianze che ne permettono la manipolazione
    algebrica.

    \begin{esempio}
        Sia $\Sigma = \{0, 1\}$ e si consideri l'espressione regolare $r = 1(0 + 1)^*00$.
        Il linguaggio associato \`e
        \begin{equation*}
            \mathcal{L}(r)
            = \mathcal{L}(1) \bigl(\mathcal{L}(0) \cup \mathcal{L}(1)\bigr)^*
                \mathcal{L}(0) \mathcal{L}(0)
            = \{1x00 \mid x \in \Sigma^*\},
        \end{equation*}
        ovvero l'insieme delle stringhe che rappresentano in binario, senza zeri
        iniziali, i multipli di $4$ maggiori di zero.
    \end{esempio}

    Il teorema seguente mostra che ogni linguaggio descritto da un'espressione regolare
    \`e riconosciuto da un $\varepsilon$-NFA.

    \begin{theorem}[di McNaughton e Yamada, 1960]
        \label{thm:mcnaughton-yamada}
        Sia $r$ un'espressione regolare. Allora esiste un $\varepsilon$-NFA $\mathcal{M}$
        tale che $\mathcal{L}(\mathcal{M}) = \mathcal{L}(r)$.
    \end{theorem}

    \begin{proof}[Dimostrazione del Teorema~\ref{thm:mcnaughton-yamada}]
        Si dimostra, pi\`u precisamente, che esiste un $\varepsilon$-NFA con un unico stato
        finale che riconosce $\mathcal{L}(r)$. Si procede per induzione sulla struttura di
        $r$, costruendo un automa per ciascun caso base e per ciascuna regola induttiva.

        \textsc{Base}. I tre casi base sono riconosciuti dagli automi della
        Figura~\ref{fig:re-base}:
        \begin{itemize}
            \item se $r = \varepsilon$, l'automa riconosce $\mathcal{L}(r) =
                \{\varepsilon\}$, perch\'e lo stato iniziale \`e anche finale e non ci sono
                transizioni;
            \item se $r = \emptyset$, l'automa riconosce $\mathcal{L}(r) = \emptyset$,
                perch\'e lo stato finale non \`e raggiungibile;
            \item se $r = a$, con $a \in \Sigma$, l'automa riconosce
                $\mathcal{L}(r) = \{a\}$.
        \end{itemize}

        \begin{figure}[H]
            \centering
            \begin{tikzpicture}[automa]
                % r = epsilon
                \node[state, initial, accepting] (e0) at (0, 0) {$q_0$};
                \node at (0, -1.1) {$r = \varepsilon$};

                % r = insieme vuoto
                \node[state, initial] (v0) at (3.3, 0) {$q_0$};
                \node[state, accepting] (vf) at (5.6, 0) {$q_f$};
                \node at (4.45, -1.1) {$r = \emptyset$};

                % r = a
                \node[state, initial] (a0) at (8.9, 0) {$q_0$};
                \node[state, accepting] (af) at (11.2, 0) {$q_f$};
                \path (a0) edge node {$a$} (af);
                \node at (10.05, -1.1) {$r = a$};
            \end{tikzpicture}
            \caption{Gli $\varepsilon$-NFA dei casi base.}
            \label{fig:re-base}
        \end{figure}

        \textsc{Passo induttivo}. Per ipotesi induttiva, per ogni sottoespressione $r_i$
        di $r$ esiste un $\varepsilon$-NFA $\mathcal{M}_i$, con stato iniziale $q_0^i$ e
        unico stato finale $q_f^i$, che riconosce $\mathcal{L}(r_i)$. Nelle figure
        seguenti ogni $\mathcal{M}_i$ \`e racchiuso in un riquadro tratteggiato, di cui si
        mostrano solo lo stato iniziale e quello finale. Il nuovo automa ha stato iniziale
        $q_0$ e unico stato finale $q_f$: gli stati $q_f^i$ non sono pi\`u finali.

        \emph{Unione}. Se $r = r_1 + r_2$, l'automa della Figura~\ref{fig:re-unione}
        riconosce $\mathcal{L}(r) = \mathcal{L}(r_1) \cup \mathcal{L}(r_2)$: da $q_0$ si
        sceglie, con una $\varepsilon$-transizione, quale dei due automi percorrere.

        \begin{figure}[H]
            \centering
            \begin{tikzpicture}[automa]
                \node[state, initial] (q0) at (0, 0) {$q_0$};
                \node[state] (a0) at (2.6, 1.2) {$q_0^1$};
                \node[state] (af) at (5.6, 1.2) {$q_f^1$};
                \node[state] (b0) at (2.6, -1.2) {$q_0^2$};
                \node[state] (bf) at (5.6, -1.2) {$q_f^2$};
                \node[state, accepting] (qf) at (8.2, 0) {$q_f$};

                \draw[interno] (a0) -- node[above] {$\mathcal{M}_1$} (af);
                \draw[interno] (b0) -- node[above] {$\mathcal{M}_2$} (bf);
                \node[sottoautoma, fit=(a0)(af)] {};
                \node[sottoautoma, fit=(b0)(bf)] {};

                \path (q0) edge node {$\varepsilon$} (a0)
                    (q0) edge node[swap] {$\varepsilon$} (b0)
                    (af) edge node {$\varepsilon$} (qf)
                    (bf) edge node[swap] {$\varepsilon$} (qf);
            \end{tikzpicture}
            \caption{Costruzione per l'unione $r_1 + r_2$.}
            \label{fig:re-unione}
        \end{figure}

        \emph{Concatenazione}. Se $r = r_1 r_2$, l'automa della
        Figura~\ref{fig:re-concatenazione} riconosce $\mathcal{L}(r) = \mathcal{L}(r_1)
        \cdot \mathcal{L}(r_2)$: dopo aver letto una stringa di $\mathcal{L}(r_1)$ si
        passa, senza leggere simboli, all'inizio di $\mathcal{M}_2$.

        \begin{figure}[H]
            \centering
            \begin{tikzpicture}[automa]
                \node[state, initial] (q0) at (0, 0) {$q_0$};
                \node[state] (a0) at (2.3, 0) {$q_0^1$};
                \node[state] (af) at (4.6, 0) {$q_f^1$};
                \node[state] (b0) at (7.0, 0) {$q_0^2$};
                \node[state] (bf) at (9.3, 0) {$q_f^2$};
                \node[state, accepting] (qf) at (11.6, 0) {$q_f$};

                \draw[interno] (a0) -- node[above] {$\mathcal{M}_1$} (af);
                \draw[interno] (b0) -- node[above] {$\mathcal{M}_2$} (bf);
                \node[sottoautoma, fit=(a0)(af)] {};
                \node[sottoautoma, fit=(b0)(bf)] {};

                \path (q0) edge node {$\varepsilon$} (a0)
                    (af) edge node {$\varepsilon$} (b0)
                    (bf) edge node {$\varepsilon$} (qf);
            \end{tikzpicture}
            \caption{Costruzione per la concatenazione $r_1 r_2$.}
            \label{fig:re-concatenazione}
        \end{figure}

        \emph{Chiusura di Kleene}. Se $r = r_1^*$, l'automa della
        Figura~\ref{fig:re-stella} riconosce $\mathcal{L}(r) = \mathcal{L}(r_1)^*$: l'arco
        da $q_f^1$ a $q_0^1$ permette di ripetere $\mathcal{M}_1$ quante volte si vuole,
        mentre quello da $q_0$ a $q_f$ accetta la stringa vuota.

        \begin{figure}[H]
            \centering
            \begin{tikzpicture}[automa]
                \node[state, initial] (q0) at (0, 0) {$q_0$};
                \node[state] (a0) at (2.6, 0) {$q_0^1$};
                \node[state] (af) at (5.6, 0) {$q_f^1$};
                \node[state, accepting] (qf) at (8.2, 0) {$q_f$};

                \draw[interno] (a0) -- node[above] {$\mathcal{M}_1$} (af);
                \node[sottoautoma, fit=(a0)(af)] {};

                \path (q0) edge node {$\varepsilon$} (a0)
                    (af) edge node {$\varepsilon$} (qf)
                    (af) edge[bend right=55] node[swap] {$\varepsilon$} (a0)
                    (q0) edge[bend right=35] node[swap] {$\varepsilon$} (qf);
            \end{tikzpicture}
            \caption{Costruzione per la chiusura di Kleene $r_1^*$.}
            \label{fig:re-stella}
        \end{figure}
    \end{proof}

    \begin{esempio}
        Si consideri l'espressione regolare $r = 0^* + 1^* + (01)^*$, che denota il
        linguaggio
        \begin{equation*}
            \mathcal{L}(r) = \{x \in \{0, 1\}^* \mid \exists m \in \mathbb{N}.\;
            x = 0^m \lor x = 1^m \lor x = (01)^m\}.
        \end{equation*}
        L'automa della Figura~\ref{fig:re-adhoc} riconosce correttamente il linguaggio, ma
        \`e costruito ad hoc e non segue la costruzione vista nella dimostrazione del
        Teorema~\ref{thm:mcnaughton-yamada}.

        \begin{figure}[H]
            \centering
            \begin{tikzpicture}[automa]
                \node[state, initial] (q0) at (0, 0) {$q_0$};
                \node[state, accepting] (q1) at (2.8, 1.8) {$q_1$};
                \node[state, accepting] (q2) at (2.8, 0) {$q_2$};
                \node[state, accepting] (q3) at (2.8, -1.8) {$q_3$};
                \node[state] (q4) at (5.4, -1.8) {$q_4$};
                \path (q0) edge node {$\varepsilon$} (q1)
                    (q0) edge node {$\varepsilon$} (q2)
                    (q0) edge node[swap] {$\varepsilon$} (q3)
                    (q1) edge[loop right] node {$0$} ()
                    (q2) edge[loop right] node {$1$} ()
                    (q3) edge[bend left=20] node {$0$} (q4)
                    (q4) edge[bend left=20] node {$1$} (q3);
            \end{tikzpicture}
            \caption{Un $\varepsilon$-NFA costruito ad hoc per $0^* + 1^* + (01)^*$.}
            \label{fig:re-adhoc}
        \end{figure}

        L'automa della Figura~\ref{fig:re-costruzione}, invece, ricalca la costruzione
        dell'unione: un nuovo stato iniziale $q_0$ e un unico stato finale $q_f$, collegati
        con $\varepsilon$-transizioni agli automi di $0^*$, $1^*$ e $(01)^*$. Per
        leggibilit\`a, questi ultimi sono semplificati rispetto alla costruzione completa,
        che userebbe pi\`u stati e pi\`u $\varepsilon$-transizioni.

        \begin{figure}[H]
            \centering
            \begin{tikzpicture}[automa]
                % Stati
                \node[state, initial] (q0) at (0, 0) {$q_0$};

                \node[state] (p0) at (3, 2) {$P_0$};
                \node[state] (p1) at (3, 0) {$P_1$};
                \node[state] (p2) at (3, -2) {$P_2$};
                \node[state] (p3) at (5.5, -2) {$P_3$};

                % Nodo qf posizionato tra P0 e P1 spostato a destra
                \node[state, accepting] (qf) at (7, 1.25) {$q_f$};

                % Transizioni epsilon da q0
                \path
                    (q0) edge[bend left=20] node[above left] {$\varepsilon$} (p0)
                    (q0) edge[bend right=10] node[below] {$\varepsilon$} (p1)
                    (q0) edge[bend right=20] node[below left] {$\varepsilon$} (p2);

                % Transizioni da P0 e P1
                \path
                    (p0) edge[loop above, looseness=4] node {$0$} (p0)
                    (p0) edge[bend left=10] node[above] {$\varepsilon$} (qf)

                    (p1) edge[loop above, looseness=4] node {$1$} (p1)
                    (p1) edge[bend right=10] node[below right] {$\varepsilon$} (qf);

                % Transizioni tra P2, P3 e verso qf
                \path
                    (p2) edge[bend left=20] node[above] {$0$} (p3)
                    (p3) edge[bend left=20] node[below] {$1$} (p2)
                    (p2) edge[out=60, in=250, looseness=1.2]
                    node[below right] {$\varepsilon$} (qf);
            \end{tikzpicture}
            \caption{Un $\varepsilon$-NFA per $0^* + 1^* + (01)^*$ con la struttura della
                costruzione dell'unione.}
            \label{fig:re-costruzione}
        \end{figure}
    \end{esempio}

    Ricapitolando, tutti i modelli visti descrivono la stessa classe di linguaggi:
    \begin{equation*}
        \text{linguaggi regolari} \iff \text{DFA} \iff \text{NFA} \iff
        \varepsilon\text{-NFA} \iff \text{RE}.
    \end{equation*}

    \section{Propriet\`a dei linguaggi regolari}
    Si indica con $\mathrm{LR}$ l'insieme dei linguaggi regolari:
    \begin{equation*}
        \mathrm{LR} = \{\mathcal{L} \subseteq \Sigma^* \mid
            \text{``$\mathcal{L}$ \`e un linguaggio regolare''}\}.
    \end{equation*}
    I linguaggi regolari godono di alcune propriet\`a di chiusura.

    \begin{theorem}[Propriet\`a di chiusura]
        \label{thm:chiusura}
        Siano $\mathcal{L}, \mathcal{L}_1, \mathcal{L}_2 \in \mathrm{LR}$. Allora:
        \begin{enumerate}[label=\textup{(\alph*)}]
            \item $\mathrm{LR}$ \`e chiuso rispetto all'unione, alla concatenazione e alla
                chiusura di Kleene:
                \begin{equation*}
                    \mathcal{L}_1 \cup \mathcal{L}_2 \in \mathrm{LR},
                    \qquad
                    \mathcal{L}_1 \cdot \mathcal{L}_2 \in \mathrm{LR},
                    \qquad
                    \mathcal{L}^* \in \mathrm{LR};
                \end{equation*}
            \item $\mathrm{LR}$ \`e chiuso rispetto al complemento:
                $\overline{\mathcal{L}} = \Sigma^* \setminus \mathcal{L} \in \mathrm{LR}$;
            \item $\mathrm{LR}$ \`e chiuso rispetto all'intersezione:
                $\mathcal{L}_1 \cap \mathcal{L}_2 \in \mathrm{LR}$.
        \end{enumerate}
    \end{theorem}

    \begin{proof}[Dimostrazione del Teorema~\ref{thm:chiusura}]
        \textsc{Punto (a)}. Deriva direttamente dall'equivalenza con le espressioni
        regolari: se $\mathcal{L}_1 = \mathcal{L}(r_1)$, $\mathcal{L}_2 = \mathcal{L}(r_2)$
        e $\mathcal{L} = \mathcal{L}(r)$, i tre linguaggi sono denotati rispettivamente
        dalle espressioni regolari $r_1 + r_2$, $r_1 r_2$ e $r^*$.

        \textsc{Punto (b)}. Sia $\mathcal{M} = \langle \mathcal{Q}, \Sigma, \delta, q_0,
        \mathcal{F} \rangle$ un DFA che riconosce $\mathcal{L}$. Basta ridefinire l'insieme
        degli stati finali come $\mathcal{F}' = \mathcal{Q} \setminus \mathcal{F}$: il DFA
        ottenuto accetta esattamente le stringhe rifiutate da $\mathcal{M}$, cio\`e
        $\overline{\mathcal{L}}$.

        \textsc{Punto (c)}. Per le leggi di De Morgan,
        \begin{equation*}
            \mathcal{L}_1 \cap \mathcal{L}_2
            = \overline{\overline{\mathcal{L}_1} \cup \overline{\mathcal{L}_2}},
        \end{equation*}
        e il membro destro \`e regolare per i punti (a) e (b).
    \end{proof}

    La chiusura rispetto all'unione e all'intersezione vale solo per un numero
    \emph{finito} di linguaggi: anche se ogni $\mathcal{L}_m$ \`e regolare, l'unione o
    l'intersezione di infiniti linguaggi regolari pu\`o non esserlo. In simboli,
    \begin{gather*}
        \bigl(\forall m \in \mathbb{N}.\; \mathcal{L}_m \in \mathrm{LR}\bigr)
        \centernot\implies \bigcup_{m \in \mathbb{N}} \mathcal{L}_m \in \mathrm{LR}, \\
        \bigl(\forall m \in \mathbb{N}.\; \mathcal{L}_m \in \mathrm{LR}\bigr)
        \centernot\implies \bigcap_{m \in \mathbb{N}} \mathcal{L}_m \in \mathrm{LR}.
    \end{gather*}

    \section{Problemi di decidibilit\`a}
    \begin{definition}[Problema decidibile]
        Un problema \`e \emph{decidibile} se esiste un algoritmo che, ricevuta in input
        un'istanza qualsiasi del problema, termina sempre in un numero finito di passi e
        fornisce la risposta corretta (s\`i o no).
    \end{definition}

    \begin{theorem}[Problemi decidibili sui linguaggi regolari]
        \label{thm:decidibilita}
        Sia $\mathcal{L}$ un linguaggio regolare e sia
        $\mathcal{M} = \langle \mathcal{Q}, \Sigma, \delta, q_0, \mathcal{F} \rangle$ un
        DFA tale che $\mathcal{L} = \mathcal{L}(\mathcal{M})$. I seguenti problemi sono
        decidibili:
        \begin{enumerate}[label=\textup{(\alph*)}]
            \item \textup{\textsc{Vacuit\`a}}: stabilire se $\mathcal{L} \neq \emptyset$.
                Vale
                \begin{equation*}
                    \mathcal{L} \neq \emptyset \iff
                    \exists x \in \mathcal{L}(\mathcal{M}).\; |x| < |\mathcal{Q}|;
                \end{equation*}
            \item \textup{\textsc{Infinit\`a}}: stabilire se $|\mathcal{L}| = \omega$,
                ovvero se $\mathcal{L}$ \`e infinito. Detto $n = |\mathcal{Q}|$ il numero
                di stati di $\mathcal{M}$, vale
                \begin{equation*}
                    |\mathcal{L}| = \omega \iff
                    \exists x \in \mathcal{L}(\mathcal{M}).\; n \leq |x| < 2n;
                \end{equation*}
            \item \textup{\textsc{Equivalenza}}: dati due DFA $\mathcal{M}_1$ e
                $\mathcal{M}_2$, stabilire se $\mathcal{L}(\mathcal{M}_1) =
                \mathcal{L}(\mathcal{M}_2)$.
        \end{enumerate}
    \end{theorem}
    Per i primi due problemi basta quindi controllare se $\mathcal{M}$ accetta almeno una
    delle stringhe di lunghezza limitata, che sono in numero finito.

    \section{Relazioni di equivalenza e linguaggi}
    \begin{definition}[Relazione indotta da un linguaggio]
        Sia $\mathcal{L} \subseteq \Sigma^*$. La relazione
        $R_{\mathcal{L}} \subseteq \Sigma^* \times \Sigma^*$ \`e definita da
        \begin{equation*}
            x \mathrel{R_{\mathcal{L}}} y \iff
            \forall z \in \Sigma^*.\; (xz \in \mathcal{L} \iff yz \in \mathcal{L}).
        \end{equation*}
    \end{definition}
    Due stringhe sono quindi in relazione se, prolungandole con uno stesso suffisso $z$
    qualsiasi, si ottengono stringhe che stanno entrambe in $\mathcal{L}$ o entrambe
    fuori.

    \begin{definition}[Relazione indotta da un DFA]
        Sia $\mathcal{M} = \langle \mathcal{Q}, \Sigma, \delta, q_0, \mathcal{F} \rangle$
        un DFA. La relazione $R_{\mathcal{M}} \subseteq \Sigma^* \times \Sigma^*$ \`e
        definita da
        \begin{equation*}
            x \mathrel{R_{\mathcal{M}}} y \iff \hat{\delta}(q_0, x) = \hat{\delta}(q_0, y).
        \end{equation*}
    \end{definition}
    Due stringhe sono quindi in relazione se portano l'automa nello stesso stato.
    $R_{\mathcal{L}}$ e $R_{\mathcal{M}}$ sono entrambe relazioni di equivalenza.

    \begin{definition}[Relazione invariante a destra]
        Una relazione di equivalenza $R \subseteq \Sigma^* \times \Sigma^*$ \`e
        \emph{invariante a destra} se
        \begin{equation*}
            x \mathrel{R} y \implies \forall z \in \Sigma^*.\; xz \mathrel{R} yz.
        \end{equation*}
    \end{definition}
    L'invarianza a destra \`e una propriet\`a che una relazione pu\`o avere o non avere, e
    la condizione \`e una sola implicazione, non un'equivalenza: da $xz \mathrel{R} yz$ non
    segue, in generale, $x \mathrel{R} y$.

    $R_{\mathcal{M}}$ \`e invariante a destra: se due stringhe portano l'automa nello
    stesso stato, leggendo poi uno stesso suffisso $z$ si arriva ancora nello stesso
    stato. Poich\'e $\hat{\delta}(q_0, xz) = \hat{\delta}(\hat{\delta}(q_0, x), z)$, si ha
    \begin{equation*}
        \hat{\delta}(q_0, x) = \hat{\delta}(q_0, y)
        \implies
        \forall z \in \Sigma^*.\; \hat{\delta}(q_0, xz) = \hat{\delta}(q_0, yz)
        \implies
        \forall z \in \Sigma^*.\; xz \mathrel{R_{\mathcal{M}}} yz.
    \end{equation*}
    Anche $R_{\mathcal{L}}$ \`e invariante a destra. Se $x \mathrel{R_{\mathcal{L}}} y$ e
    $w \in \Sigma^*$, per ogni $z \in \Sigma^*$ si applica la definizione di
    $R_{\mathcal{L}}$ al suffisso $wz$, ottenendo
    \begin{equation*}
        (xw)z \in \mathcal{L} \iff (yw)z \in \mathcal{L},
    \end{equation*}
    cio\`e $xw \mathrel{R_{\mathcal{L}}} yw$.

    \begin{definition}[Linguaggio indotto da uno stato]
        Sia $\mathcal{M}$ un DFA e sia $q$ un suo stato. Il linguaggio \emph{indotto}
        da $q$ \`e
        \begin{equation*}
            \mathcal{L}_q = \{x \in \Sigma^* \mid \hat{\delta}(q_0, x) = q\},
        \end{equation*}
        ovvero il linguaggio ``accettato'' dallo stato $q$.
    \end{definition}
    Le stringhe accettate da $\mathcal{M}$ sono quelle che portano in uno stato finale,
    quindi
    \begin{equation*}
        \mathcal{L}(\mathcal{M}) = \bigcup_{q \in \mathcal{F}} \mathcal{L}_q.
    \end{equation*}

    \begin{definition}[Indice]
        L'\emph{indice} di una relazione di equivalenza \`e il numero delle sue classi di
        equivalenza. La relazione \`e \emph{di indice finito} se le sue classi di
        equivalenza sono in numero finito.
    \end{definition}

    \begin{definition}[Raffinamento]
        Siano $R$ e $R'$ relazioni di equivalenza sullo stesso insieme. $R$ \`e un
        \emph{raffinamento} di $R'$ se
        \begin{equation*}
            x \mathrel{R} y \implies x \mathrel{R'} y.
        \end{equation*}
    \end{definition}
    Le seguenti affermazioni sono equivalenti alla definizione:
    \begin{itemize}
        \item ogni classe di equivalenza di $R$ \`e contenuta in una classe di $R'$, cio\`e
            $[x]_R \subseteq [x]_{R'}$ per ogni $x$;
        \item ogni classe di equivalenza di $R'$ \`e unione di classi di equivalenza
            di $R$.
    \end{itemize}
    Ne segue che l'indice di $R$ \`e maggiore o uguale a quello di $R'$ e, in
    particolare, che se $R$ \`e di indice finito anche $R'$ lo \`e. Queste ultime sono
    solo conseguenze: un indice maggiore o uguale non basta perch\'e $R$ raffini $R'$.

    \begin{theorem}[di Myhill e Nerode]
        \label{thm:myhill-nerode}
        Sia $\mathcal{L} \subseteq \Sigma^*$. I seguenti enunciati sono equivalenti:
        \begin{enumerate}[label=\textup{(\arabic*)}]
            \item $\mathcal{L}$ \`e riconosciuto da un DFA, cio\`e \`e regolare;
            \item $\mathcal{L}$ \`e unione di classi di equivalenza di una relazione di
                equivalenza su $\Sigma^*$ invariante a destra e di indice finito;
            \item $R_{\mathcal{L}}$ \`e di indice finito.
        \end{enumerate}
    \end{theorem}

    \begin{proof}[Dimostrazione del Teorema~\ref{thm:myhill-nerode}]
        Si dimostrano le implicazioni in catena,
        \begin{equation*}
            (1) \implies (2) \implies (3) \implies (1),
        \end{equation*}
        da cui segue che i tre enunciati sono equivalenti.

        \textsc{(1)} $\Rightarrow$ \textsc{(2)}. Si suppone che valga (1), cio\`e che
        $\mathcal{L}$ sia regolare: esiste un DFA
        $\mathcal{M} = \langle \mathcal{Q}, \Sigma, \delta, q_0, \mathcal{F} \rangle$ tale
        che $\mathcal{L} = \mathcal{L}(\mathcal{M})$. Si dimostra che vale (2) con la
        relazione $R_{\mathcal{M}}$. Come visto,
        \begin{equation*}
            \mathcal{L} = \bigcup_{q \in \mathcal{F}} \mathcal{L}_q
            = \bigcup_{q \in \mathcal{F}} \{x \in \Sigma^* \mid \hat{\delta}(q_0, x) = q\},
        \end{equation*}
        quindi basta verificare che:
        \begin{itemize}
            \item $R_{\mathcal{M}}$ \`e invariante a destra, come mostrato sopra;
            \item i linguaggi indotti $\mathcal{L}_q$ sono le classi di equivalenza di
                $R_{\mathcal{M}}$. Infatti, se $x \in \mathcal{L}_q$, per ogni
                $y \in \Sigma^*$ si ha
                \begin{equation*}
                    x \mathrel{R_{\mathcal{M}}} y
                    \iff \hat{\delta}(q_0, y) = \hat{\delta}(q_0, x) = q
                    \iff y \in \mathcal{L}_q,
                \end{equation*}
                cio\`e $\mathcal{L}_q = [x]_{R_{\mathcal{M}}}$;
            \item $R_{\mathcal{M}}$ \`e di indice finito: $\mathcal{M}$ \`e a stati finiti,
                quindi $|\mathcal{Q}|$ \`e finito, e le classi di equivalenza sono una per
                ogni stato raggiungibile da $q_0$, quindi al pi\`u
                $|\mathcal{Q}|$.
        \end{itemize}
        Quindi $\mathcal{L}$ \`e unione delle classi di equivalenza $\mathcal{L}_q$, con
        $q \in \mathcal{F}$, di una relazione invariante a destra e di indice finito,
        ovvero vale (2).

        \textsc{(2)} $\Rightarrow$ \textsc{(3)}. Si suppone che valga (2): $\mathcal{L}$
        \`e unione di classi di equivalenza di una relazione $R$ invariante a destra e di
        indice finito (quindi $\mathcal{L}$ non taglia le classi: non \`e possibile che,
        di due elementi in relazione tra loro, uno stia in $\mathcal{L}$ e l'altro fuori). Si dimostra che $R$ \`e un raffinamento di $R_{\mathcal{L}}$, cio\`e
        che
        \begin{equation*}
            \forall x \in \Sigma^*.\; [x]_R \subseteq [x]_{R_{\mathcal{L}}}.
        \end{equation*}
        Siano $x \in \Sigma^*$ e $y \in [x]_R$, cio\`e $x \mathrel{R} y$. Allora
        \begin{align*}
            x \mathrel{R} y
                &\implies \forall z \in \Sigma^*.\; xz \mathrel{R} yz
                    && \text{($R$ \`e invariante a destra)} \\
                &\implies \forall z \in \Sigma^*.\;
                    (xz \in \mathcal{L} \iff yz \in \mathcal{L})
                    && \text{($\mathcal{L}$ non taglia le classi di $R$)} \\
                &\iff x \mathrel{R_{\mathcal{L}}} y
                    && \text{(definizione di $R_{\mathcal{L}}$)} \\
                &\iff y \in [x]_{R_{\mathcal{L}}}
                    && \text{(definizione di classe di equivalenza)}.
        \end{align*}
        Si \`e partiti da $y \in [x]_R$ e si \`e arrivati a $y \in [x]_{R_{\mathcal{L}}}$,
        quindi $[x]_R \subseteq [x]_{R_{\mathcal{L}}}$. Dunque $R$ \`e un raffinamento di $R_{\mathcal{L}}$ e, poich\'e $R$ \`e di indice
        finito, anche $R_{\mathcal{L}}$ \`e di indice finito: vale (3).

        \textsc{(3)} $\Rightarrow$ \textsc{(1)}. Si suppone che valga (3):
        $R_{\mathcal{L}}$ \`e di indice finito. La dimostrazione si articola in tre passi:
        \begin{enumerate}
            \item costruire un DFA $\mathcal{M}'$ i cui stati sono le classi di
                equivalenza di $R_{\mathcal{L}}$;
            \item dimostrare che, leggendo una stringa $y$ a partire dalla classe di $x$,
                $\mathcal{M}'$ arriva nella classe di $xy$;
            \item dedurne che $\mathcal{L} = \mathcal{L}(\mathcal{M}')$.
        \end{enumerate}

        \emph{Passo 1: costruzione di $\mathcal{M}'$}. Si pone
        $\mathcal{M}' = \langle \mathcal{Q}', \Sigma', \delta', q_0', \mathcal{F}' \rangle$
        con:
        \begin{itemize}
            \item $\mathcal{Q}' = \{[x]_{R_{\mathcal{L}}} \mid x \in \Sigma^*\}$, che \`e
                un insieme finito per ipotesi;
            \item $\Sigma' = \Sigma$, l'alfabeto di $\mathcal{L}$;
            \item $\delta' \colon \mathcal{Q}' \times \Sigma \to \mathcal{Q}'$, con
                $\delta'([x]_{R_{\mathcal{L}}}, a) \coloneqq [xa]_{R_{\mathcal{L}}}$;
            \item $q_0' = [\varepsilon]_{R_{\mathcal{L}}}$;
            \item $\mathcal{F}' = \{[x]_{R_{\mathcal{L}}} \mid x \in \mathcal{L}\}$.
        \end{itemize}
        Quella di $\delta'$ \`e una buona definizione: il risultato non dipende
        dall'elemento $x$ scelto per rappresentare la classe. Infatti, se
        $[x]_{R_{\mathcal{L}}} = [y]_{R_{\mathcal{L}}}$, cio\`e
        $x \mathrel{R_{\mathcal{L}}} y$, allora
        \begin{align*}
            x \mathrel{R_{\mathcal{L}}} y
                &\implies \forall z \in \Sigma^*.\; xz \mathrel{R_{\mathcal{L}}} yz
                    && \text{($R_{\mathcal{L}}$ \`e invariante a destra)} \\
                &\implies xa \mathrel{R_{\mathcal{L}}} ya
                    && \text{(caso $z = a$)} \\
                &\iff [xa]_{R_{\mathcal{L}}} = [ya]_{R_{\mathcal{L}}}.
        \end{align*}

        \emph{Passo 2: stato raggiunto leggendo $y$}. Si dimostra, per induzione sulla
        lunghezza $|y|$ della stringa letta, che
        \begin{equation*}
            \forall [x]_{R_{\mathcal{L}}} \in \mathcal{Q}'.\;
            \hat{\delta}'([x]_{R_{\mathcal{L}}}, y) = [xy]_{R_{\mathcal{L}}}.
        \end{equation*}
        \begin{itemize}
            \item \textit{Base} ($y = \varepsilon$). Leggendo $\varepsilon$ non si cambia
                stato e $x = x\varepsilon$ per ogni $x$, quindi
                \begin{equation*}
                    \hat{\delta}'([x]_{R_{\mathcal{L}}}, \varepsilon)
                    = [x]_{R_{\mathcal{L}}} = [x\varepsilon]_{R_{\mathcal{L}}}.
                \end{equation*}
            \item \textit{Passo induttivo}. Si suppone che l'uguaglianza valga per ogni
                $y$ con $|y| \leq m$ e sia $ya$ una stringa di lunghezza $m + 1$. Allora
                \begin{align*}
                    \hat{\delta}'([x]_{R_{\mathcal{L}}}, ya)
                        &= \delta'(\hat{\delta}'([x]_{R_{\mathcal{L}}}, y), a)
                            && \text{(definizione di $\hat{\delta}'$)} \\
                        &= \delta'([xy]_{R_{\mathcal{L}}}, a)
                            && \text{(ipotesi induttiva)} \\
                        &= [xya]_{R_{\mathcal{L}}}
                            && \text{(definizione di $\delta'$)}.
                \end{align*}
        \end{itemize}

        \emph{Passo 3: $\mathcal{L} = \mathcal{L}(\mathcal{M}')$}. Bisogna dimostrare che
        \begin{equation*}
            \forall y \in \Sigma^*.\;
            y \in \mathcal{L}(\mathcal{M}') \iff y \in \mathcal{L}.
        \end{equation*}
        Si usa il passo~2 con $x = \varepsilon$, cio\`e partendo dallo stato iniziale
        $q_0' = [\varepsilon]_{R_{\mathcal{L}}}$: la stringa $y$ resta qualsiasi, quindi
        serve il passo~2 per intero e non solo il suo caso base. Poich\'e $\varepsilon$
        concatenata a sinistra non cambia la stringa, $[\varepsilon y]_{R_{\mathcal{L}}} =
        [y]_{R_{\mathcal{L}}}$, e si ottiene
        \begin{align*}
            y \in \mathcal{L}(\mathcal{M}')
                &\iff \hat{\delta}'(q_0', y) \in \mathcal{F}'
                    && \text{(definizione di $\mathcal{L}(\mathcal{M}')$)} \\
                &\iff \hat{\delta}'([\varepsilon]_{R_{\mathcal{L}}}, y) \in \mathcal{F}'
                    && \text{(definizione di $q_0'$)} \\
                &\iff [\varepsilon y]_{R_{\mathcal{L}}} = [y]_{R_{\mathcal{L}}}
                    \in \mathcal{F}'
                    && \text{(passo~2 con $x = \varepsilon$)} \\
                &\iff y \in \mathcal{L}
                    && \text{(definizione di $\mathcal{F}'$)}.
        \end{align*}
        Nell'ultimo passaggio, $[y]_{R_{\mathcal{L}}} \in \mathcal{F}' \iff
        y \in \mathcal{L}$, il verso ($\Leftarrow$) \`e immediato, mentre il verso
        ($\Rightarrow$) richiede un argomento in pi\`u:
        \begin{itemize}[leftmargin=3em]
            \item[($\Leftarrow$)] se $y \in \mathcal{L}$, allora
                $[y]_{R_{\mathcal{L}}} \in \mathcal{F}'$ per definizione di
                $\mathcal{F}'$;
            \item[($\Rightarrow$)] se $[y]_{R_{\mathcal{L}}} \in \mathcal{F}'$, allora
                $[y]_{R_{\mathcal{L}}} = [x]_{R_{\mathcal{L}}}$ per qualche
                $x \in \mathcal{L}$, non necessariamente $y$ stesso, quindi
                $y \mathrel{R_{\mathcal{L}}} x$. Prendendo $z = \varepsilon$ nella
                definizione di $R_{\mathcal{L}}$ si ottiene
                $y \in \mathcal{L} \iff x \in \mathcal{L}$ e, poich\'e
                $x \in \mathcal{L}$, anche $y \in \mathcal{L}$.
        \end{itemize}
        Quindi $\mathcal{L} = \mathcal{L}(\mathcal{M}')$: $\mathcal{L}$ \`e regolare e
        vale (1).
    \end{proof}

    %TEORICA%
    \chapter{Esercizi}
    \section{Automi a stati finiti deterministici}
    Gli esercizi seguenti dimostrano che un linguaggio \`e regolare costruendo un DFA che
    lo riconosce, con il metodo della Sezione~\ref{sec:dimostrare-regolare}.

    \esercizio
    \textit{Si dimostri che il linguaggio $\mathcal{L} = \{0^{3n} \mid n \geq 0\}$
    sull'alfabeto $\Sigma = \{0, 1\}$ \`e regolare.}

    \textsc{Svolgimento}. Il linguaggio contiene le stringhe formate solo da zeri, in numero
    multiplo di $3$. Lo riconosce il DFA
    $\mathcal{M} = \langle \{q_0, q_1, q_2, q_3\}, \{0, 1\}, \delta, q_0, \{q_0\} \rangle$
    della Figura~\ref{fig:es-zeri}: gli stati $q_0$, $q_1$ e $q_2$ contano gli zeri letti
    modulo $3$, mentre $q_3$ \`e lo stato in cui si finisce, senza pi\`u uscirne, appena si
    legge un $1$.

    \begin{figure}[H]
        \centering
        \begin{tikzpicture}[automa]
            \node[state, initial, accepting] (q0) {$q_0$};
            \node[state, right=of q0] (q1) {$q_1$};
            \node[state, below=of q1] (q2) {$q_2$};
            \node[state, below=of q0] (q3) {$q_3$};
            \path (q0) edge node {$0$} (q1)
                (q1) edge node {$0$} (q2)
                (q2) edge node[pos=0.2] {$0$} (q0)
                (q0) edge node[swap] {$1$} (q3)
                (q1) edge node[pos=0.2, swap] {$1$} (q3)
                (q2) edge node {$1$} (q3)
                (q3) edge[loop below] node {$0, 1$} ();

            % Matrice di transizione, a destra del diagramma
            \node[right=1.8cm of q1, yshift=-1.25cm] {%
                \begin{tabular}{ccc}
                    \toprule
                    $\delta$ & $0$ & $1$ \\
                    \midrule
                    $q_0$ & $q_1$ & $q_3$ \\
                    $q_1$ & $q_2$ & $q_3$ \\
                    $q_2$ & $q_0$ & $q_3$ \\
                    $q_3$ & $q_3$ & $q_3$ \\
                    \bottomrule
                \end{tabular}};
        \end{tikzpicture}
        \caption{DFA che riconosce $\{0^{3n} \mid n \geq 0\}$ e la sua matrice di
            transizione.}
        \label{fig:es-zeri}
    \end{figure}

    Bisogna dimostrare che $\mathcal{L}(\mathcal{M}) = \mathcal{L}$, cio\`e
    \begin{equation*}
        \begin{cases}
            x \in \mathcal{L} \implies \hat{\delta}(q_0, x) \in \mathcal{F}, \\
            x \notin \mathcal{L} \implies \hat{\delta}(q_0, x) \notin \mathcal{F}.
        \end{cases}
    \end{equation*}
    Come nell'esempio della Sezione~\ref{sec:dimostrare-regolare}, la seconda condizione
    non basta per il passo induttivo, perch\'e non dice in quale dei tre stati non finali
    si trovi l'automa. La si scorpora quindi secondo il significato degli stati, ottenendo
    la propriet\`a pi\`u forte
    \begin{equation*}
        \Psi(x) =
        \begin{cases}
            x \in \mathcal{L} \implies \hat{\delta}(q_0, x) = q_0,
                & \text{(1)} \\
            x \notin \mathcal{L} \land x = 0^{3i+1} \implies \hat{\delta}(q_0, x) = q_1,
                & \text{(2)} \\
            x \notin \mathcal{L} \land x = 0^{3i+2} \implies \hat{\delta}(q_0, x) = q_2,
                & \text{(3)} \\
            x \notin \mathcal{L} \land \text{``$x$ contiene almeno un $1$''}
                \implies \hat{\delta}(q_0, x) = q_3,
                & \text{(4)}
        \end{cases}
    \end{equation*}
    con $i \geq 0$. Le quattro condizioni coprono tutte le stringhe e si escludono a
    vicenda: una stringa o contiene almeno un $1$, oppure \`e della forma $0^k$, e in tal
    caso $k$ diviso per $3$ d\`a resto $0$, $1$ o $2$.

    Per $|x| \leq 2$ esistono sia stringhe in $\mathcal{L}$ sia stringhe non in
    $\mathcal{L}$, quindi si sceglie questa come base:
    \begin{equation*}
        \left.
        \begin{alignedat}{3}
            |x| = 0 \quad & x = \varepsilon \in \mathcal{L} & \qquad
                & \hat{\delta}(q_0, \varepsilon) = q_0
                &\quad \text{(caso 1 di $\Psi$)} \\
            |x| = 1 \quad & x = 0 \notin \mathcal{L} &
                & \hat{\delta}(q_0, 0) = q_1
                &\quad \text{(caso 2 di $\Psi$)} \\
            & x = 1 \notin \mathcal{L} &
                & \hat{\delta}(q_0, 1) = q_3
                &\quad \text{(caso 4 di $\Psi$)} \\
            |x| = 2 \quad & x = 00 \notin \mathcal{L} &
                & \hat{\delta}(q_0, 00) = q_2
                &\quad \text{(caso 3 di $\Psi$)} \\
            & x = 01 \notin \mathcal{L} &
                & \hat{\delta}(q_0, 01) = q_3
                &\quad \text{(caso 4 di $\Psi$)} \\
            & x = 10 \notin \mathcal{L} &
                & \hat{\delta}(q_0, 10) = q_3
                &\quad \text{(caso 4 di $\Psi$)} \\
            & x = 11 \notin \mathcal{L} &
                & \hat{\delta}(q_0, 11) = q_3
                &\quad \text{(caso 4 di $\Psi$)}
        \end{alignedat}
        \right\}
        \quad \text{\textsc{Base}}
    \end{equation*}

    \textsc{Passo induttivo}. Si suppone che $\Psi(x)$ valga per ogni $x \in \Sigma^*$ con
    $|x| \leq n$, per un certo $n \geq 2$, e sia $x$ una stringa con $|x| = n + 1$. Allora
    $x = x'0$ oppure $x = x'1$, con $|x'| = n$: per ipotesi induttiva vale $\Psi(x')$, e
    $x'$ soddisfa esattamente una delle quattro condizioni. Si distingue un caso per
    ciascuna di esse.
    \begin{enumerate}
        \item $x' \in \mathcal{L}$, cio\`e $x' = 0^{3i}$: per ipotesi induttiva
            $\hat{\delta}(q_0, x') = q_0$.
            \begin{enumerate}[label=\arabic{enumi}\alph*.]
                \item $x = x'0 = 0^{3i+1}$, quindi $x \notin \mathcal{L}$ (caso 2 di
                    $\Psi$):
                    \begin{equation*}
                        \hat{\delta}(q_0, x'0)
                        \overset{\text{def}}{=} \delta(\hat{\delta}(q_0, x'), 0)
                        \overset{\text{Hp}}{=} \delta(q_0, 0) = q_1.
                    \end{equation*}
                \item $x = x'1$ contiene un $1$, quindi $x \notin \mathcal{L}$ (caso 4 di
                    $\Psi$):
                    \begin{equation*}
                        \hat{\delta}(q_0, x'1)
                        \overset{\text{def}}{=} \delta(\hat{\delta}(q_0, x'), 1)
                        \overset{\text{Hp}}{=} \delta(q_0, 1) = q_3.
                    \end{equation*}
            \end{enumerate}
        \item $x' \notin \mathcal{L}$ e $x' = 0^{3i+1}$: per ipotesi induttiva
            $\hat{\delta}(q_0, x') = q_1$.
            \begin{enumerate}[label=\arabic{enumi}\alph*.]
                \item $x = x'0 = 0^{3i+2}$, quindi $x \notin \mathcal{L}$ (caso 3 di
                    $\Psi$):
                    \begin{equation*}
                        \hat{\delta}(q_0, x'0)
                        \overset{\text{def}}{=} \delta(\hat{\delta}(q_0, x'), 0)
                        \overset{\text{Hp}}{=} \delta(q_1, 0) = q_2.
                    \end{equation*}
                \item $x = x'1$ contiene un $1$, quindi $x \notin \mathcal{L}$ (caso 4 di
                    $\Psi$):
                    \begin{equation*}
                        \hat{\delta}(q_0, x'1)
                        \overset{\text{def}}{=} \delta(\hat{\delta}(q_0, x'), 1)
                        \overset{\text{Hp}}{=} \delta(q_1, 1) = q_3.
                    \end{equation*}
            \end{enumerate}
        \item $x' \notin \mathcal{L}$ e $x' = 0^{3i+2}$: per ipotesi induttiva
            $\hat{\delta}(q_0, x') = q_2$.
            \begin{enumerate}[label=\arabic{enumi}\alph*.]
                \item $x = x'0 = 0^{3i+3} = 0^{3(i+1)}$, quindi $x \in \mathcal{L}$ (caso 1
                    di $\Psi$):
                    \begin{equation*}
                        \hat{\delta}(q_0, x'0)
                        \overset{\text{def}}{=} \delta(\hat{\delta}(q_0, x'), 0)
                        \overset{\text{Hp}}{=} \delta(q_2, 0) = q_0.
                    \end{equation*}
                \item $x = x'1$ contiene un $1$, quindi $x \notin \mathcal{L}$ (caso 4 di
                    $\Psi$):
                    \begin{equation*}
                        \hat{\delta}(q_0, x'1)
                        \overset{\text{def}}{=} \delta(\hat{\delta}(q_0, x'), 1)
                        \overset{\text{Hp}}{=} \delta(q_2, 1) = q_3.
                    \end{equation*}
            \end{enumerate}
        \item $x' \notin \mathcal{L}$ e $x'$ contiene almeno un $1$: per ipotesi induttiva
            $\hat{\delta}(q_0, x') = q_3$. Qualunque sia l'ultimo simbolo $a \in \{0, 1\}$,
            anche $x = x'a$ contiene almeno un $1$, quindi $x \notin \mathcal{L}$ (caso 4 di
            $\Psi$):
            \begin{equation*}
                \hat{\delta}(q_0, x'a)
                \overset{\text{def}}{=} \delta(\hat{\delta}(q_0, x'), a)
                \overset{\text{Hp}}{=} \delta(q_3, a) = q_3.
            \end{equation*}
    \end{enumerate}
    In tutti i casi lo stato raggiunto da $x$ \`e quello previsto da $\Psi(x)$, quindi
    $\Psi(x)$ vale anche per $|x| = n + 1$ e, per induzione, per ogni $x \in \Sigma^*$. In
    particolare, se $x \in \mathcal{L}$ allora $\hat{\delta}(q_0, x) = q_0 \in \mathcal{F}$;
    se $x \notin \mathcal{L}$ allora $\hat{\delta}(q_0, x) \in \{q_1, q_2, q_3\}$, e nessuno
    di questi stati \`e finale. Quindi $\mathcal{L}(\mathcal{M}) = \mathcal{L}$.

    \esercizio
    \textit{Si dimostri che il linguaggio
    $\mathcal{L} = \{x \in \Sigma^* \mid \exists v \in \Sigma^*,\, a \in \Sigma.\;
    x = v0a\}$ sull'alfabeto $\Sigma = \{0, 1\}$, formato dalle stringhe il cui penultimo
    simbolo \`e $0$, \`e regolare.}

    \textsc{Svolgimento}. Lo riconosce il DFA
    $\mathcal{M} = \langle \{q_0, q_1, q_2, q_3\}, \{0, 1\}, \delta, q_0, \{q_2, q_3\}
    \rangle$ della Figura~\ref{fig:es-penultimo}. Ogni stato ricorda gli ultimi due simboli
    letti: $q_0$ corrisponde a $11$, $q_1$ a $10$, $q_2$ a $00$ e $q_3$ a $01$. Gli stati
    finali sono quelli in cui il penultimo simbolo \`e $0$, cio\`e $q_2$ e $q_3$.

    \begin{figure}[H]
        \centering
        \begin{tikzpicture}[automa]
            \node[state, initial] (q0) {$q_0$};
            \node[state, right=of q0] (q1) {$q_1$};
            \node[state, accepting, below=of q1] (q2) {$q_2$};
            \node[state, accepting, below=of q0] (q3) {$q_3$};
            \path (q0) edge node {$0$} (q1)
                (q0) edge[loop above] node {$1$} ()
                (q1) edge node {$0$} (q2)
                (q1) edge[bend left=12] node[pos=0.3] {$1$} (q3)
                (q3) edge[bend left=12] node[pos=0.3] {$0$} (q1)
                (q2) edge[loop right] node {$0$} ()
                (q2) edge node {$1$} (q3)
                (q3) edge node {$1$} (q0);

            % Matrice di transizione, a destra del diagramma
            \node[right=2.6cm of q1, yshift=-1.25cm] {%
                \begin{tabular}{ccc}
                    \toprule
                    $\delta$ & $0$ & $1$ \\
                    \midrule
                    $q_0$ & $q_1$ & $q_0$ \\
                    $q_1$ & $q_2$ & $q_3$ \\
                    $q_2$ & $q_2$ & $q_3$ \\
                    $q_3$ & $q_1$ & $q_0$ \\
                    \bottomrule
                \end{tabular}};
        \end{tikzpicture}
        \caption{DFA che riconosce le stringhe con penultimo simbolo $0$ e la sua matrice di
            transizione.}
        \label{fig:es-penultimo}
    \end{figure}

    Le stringhe con meno di due simboli non hanno un penultimo simbolo, quindi non
    appartengono a $\mathcal{L}$. Per le stringhe con almeno due simboli si scorpora la
    tesi secondo gli ultimi due simboli letti, ottenendo
    \begin{equation*}
        \Psi(x) =
        \begin{cases}
            x \in \mathcal{L} \land x = v00 \implies \hat{\delta}(q_0, x) = q_2,
                & \text{(1)} \\
            x \in \mathcal{L} \land x = v01 \implies \hat{\delta}(q_0, x) = q_3,
                & \text{(2)} \\
            x \notin \mathcal{L} \land x = v10 \implies \hat{\delta}(q_0, x) = q_1,
                & \text{(3)} \\
            x \notin \mathcal{L} \land x = v11 \implies \hat{\delta}(q_0, x) = q_0,
                & \text{(4)}
        \end{cases}
    \end{equation*}
    con $v \in \Sigma^*$. Per $|x| \geq 2$ le quattro condizioni coprono tutte le stringhe
    e si escludono a vicenda, perch\'e dipendono solo dagli ultimi due simboli.

    La prima lunghezza per cui esistono sia stringhe in $\mathcal{L}$ sia stringhe non in
    $\mathcal{L}$ \`e $|x| = 2$, quindi la base comprende tutte le stringhe con
    $|x| \leq 2$. Quelle pi\`u corte non soddisfano nessuna condizione di $\Psi$ e si
    verificano direttamente sulla tesi:
    \begin{equation*}
        \left.
        \begin{alignedat}{3}
            |x| = 0 \quad & x = \varepsilon \notin \mathcal{L} & \qquad
                & \hat{\delta}(q_0, \varepsilon) = q_0 \notin \mathcal{F} & & \\
            |x| = 1 \quad & x = 0 \notin \mathcal{L} &
                & \hat{\delta}(q_0, 0) = q_1 \notin \mathcal{F} & & \\
            & x = 1 \notin \mathcal{L} &
                & \hat{\delta}(q_0, 1) = q_0 \notin \mathcal{F} & & \\
            |x| = 2 \quad & x = 00 \in \mathcal{L} &
                & \hat{\delta}(q_0, 00) = q_2
                &\quad & \text{(caso 1 di $\Psi$)} \\
            & x = 01 \in \mathcal{L} &
                & \hat{\delta}(q_0, 01) = q_3
                && \text{(caso 2 di $\Psi$)} \\
            & x = 10 \notin \mathcal{L} &
                & \hat{\delta}(q_0, 10) = q_1
                && \text{(caso 3 di $\Psi$)} \\
            & x = 11 \notin \mathcal{L} &
                & \hat{\delta}(q_0, 11) = q_0
                && \text{(caso 4 di $\Psi$)}
        \end{alignedat}
        \right\}
        \quad \text{\textsc{Base}}
    \end{equation*}

    \textsc{Passo induttivo}. Si suppone che $\Psi(x)$ valga per ogni $x \in \Sigma^*$ con
    $2 \leq |x| \leq n$, per un certo $n \geq 2$, e sia $x$ una stringa con
    $|x| = n + 1$. Allora $x = x'0$ oppure $x = x'1$, con $|x'| = n \geq 2$: per ipotesi
    induttiva vale $\Psi(x')$, e $x'$ soddisfa esattamente una delle quattro condizioni.
    Gli ultimi due simboli di $x$ sono l'ultimo simbolo di $x'$ seguito da quello appena
    letto.
    \begin{enumerate}
        \item $x' \in \mathcal{L}$ e $x' = v00$: per ipotesi induttiva
            $\hat{\delta}(q_0, x') = q_2$.
            \begin{enumerate}[label=\arabic{enumi}\alph*.]
                \item $x = x'0 = (v0)00$, quindi $x \in \mathcal{L}$ (caso 1 di $\Psi$):
                    \begin{equation*}
                        \hat{\delta}(q_0, x'0)
                        \overset{\text{def}}{=} \delta(\hat{\delta}(q_0, x'), 0)
                        \overset{\text{Hp}}{=} \delta(q_2, 0) = q_2.
                    \end{equation*}
                \item $x = x'1 = (v0)01$, quindi $x \in \mathcal{L}$ (caso 2 di $\Psi$):
                    \begin{equation*}
                        \hat{\delta}(q_0, x'1)
                        \overset{\text{def}}{=} \delta(\hat{\delta}(q_0, x'), 1)
                        \overset{\text{Hp}}{=} \delta(q_2, 1) = q_3.
                    \end{equation*}
            \end{enumerate}
        \item $x' \in \mathcal{L}$ e $x' = v01$: per ipotesi induttiva
            $\hat{\delta}(q_0, x') = q_3$.
            \begin{enumerate}[label=\arabic{enumi}\alph*.]
                \item $x = x'0 = (v0)10$, quindi $x \notin \mathcal{L}$ (caso 3 di $\Psi$):
                    \begin{equation*}
                        \hat{\delta}(q_0, x'0)
                        \overset{\text{def}}{=} \delta(\hat{\delta}(q_0, x'), 0)
                        \overset{\text{Hp}}{=} \delta(q_3, 0) = q_1.
                    \end{equation*}
                \item $x = x'1 = (v0)11$, quindi $x \notin \mathcal{L}$ (caso 4 di $\Psi$):
                    \begin{equation*}
                        \hat{\delta}(q_0, x'1)
                        \overset{\text{def}}{=} \delta(\hat{\delta}(q_0, x'), 1)
                        \overset{\text{Hp}}{=} \delta(q_3, 1) = q_0.
                    \end{equation*}
            \end{enumerate}
        \item $x' \notin \mathcal{L}$ e $x' = v10$: per ipotesi induttiva
            $\hat{\delta}(q_0, x') = q_1$.
            \begin{enumerate}[label=\arabic{enumi}\alph*.]
                \item $x = x'0 = (v1)00$, quindi $x \in \mathcal{L}$ (caso 1 di $\Psi$):
                    \begin{equation*}
                        \hat{\delta}(q_0, x'0)
                        \overset{\text{def}}{=} \delta(\hat{\delta}(q_0, x'), 0)
                        \overset{\text{Hp}}{=} \delta(q_1, 0) = q_2.
                    \end{equation*}
                \item $x = x'1 = (v1)01$, quindi $x \in \mathcal{L}$ (caso 2 di $\Psi$):
                    \begin{equation*}
                        \hat{\delta}(q_0, x'1)
                        \overset{\text{def}}{=} \delta(\hat{\delta}(q_0, x'), 1)
                        \overset{\text{Hp}}{=} \delta(q_1, 1) = q_3.
                    \end{equation*}
            \end{enumerate}
        \item $x' \notin \mathcal{L}$ e $x' = v11$: per ipotesi induttiva
            $\hat{\delta}(q_0, x') = q_0$.
            \begin{enumerate}[label=\arabic{enumi}\alph*.]
                \item $x = x'0 = (v1)10$, quindi $x \notin \mathcal{L}$ (caso 3 di $\Psi$):
                    \begin{equation*}
                        \hat{\delta}(q_0, x'0)
                        \overset{\text{def}}{=} \delta(\hat{\delta}(q_0, x'), 0)
                        \overset{\text{Hp}}{=} \delta(q_0, 0) = q_1.
                    \end{equation*}
                \item $x = x'1 = (v1)11$, quindi $x \notin \mathcal{L}$ (caso 4 di $\Psi$):
                    \begin{equation*}
                        \hat{\delta}(q_0, x'1)
                        \overset{\text{def}}{=} \delta(\hat{\delta}(q_0, x'), 1)
                        \overset{\text{Hp}}{=} \delta(q_0, 1) = q_0.
                    \end{equation*}
            \end{enumerate}
    \end{enumerate}
    In tutti i casi lo stato raggiunto da $x$ \`e quello previsto da $\Psi(x)$, quindi
    $\Psi(x)$ vale per ogni $x$ con $|x| \geq 2$. Ne segue la tesi: se $x \in \mathcal{L}$,
    allora $|x| \geq 2$ e $x$ termina con $00$ o $01$, quindi
    $\hat{\delta}(q_0, x) \in \{q_2, q_3\} = \mathcal{F}$; se $x \notin \mathcal{L}$, allora
    $x$ ha meno di due simboli, e la base mostra che $\hat{\delta}(q_0, x) \notin
    \mathcal{F}$, oppure termina con $10$ o $11$, e $\hat{\delta}(q_0, x) \in \{q_0, q_1\}$,
    stati non finali. Quindi $\mathcal{L}(\mathcal{M}) = \mathcal{L}$.

    \esercizio
    \textit{Si dimostri che il linguaggio
    $\mathcal{L} = \{x \in \Sigma^* \mid \text{``$x$ contiene almeno due $0$
    consecutivi''}\}$ sull'alfabeto $\Sigma = \{0, 1\}$ \`e regolare.}

    \textsc{Svolgimento}. Lo riconosce il DFA
    $\mathcal{M} = \langle \{q_0, q_1, q_2\}, \{0, 1\}, \delta, q_0, \{q_2\} \rangle$
    della Figura~\ref{fig:es-due-zeri}. Finch\'e non ha letto due $0$ consecutivi l'automa
    ricorda soltanto l'ultimo simbolo: \`e in $q_1$ se l'ultimo simbolo \`e $0$ e in $q_0$
    altrimenti. Appena legge il secondo $0$ consecutivo passa in $q_2$, da cui non esce
    pi\`u.

    \begin{figure}[H]
        \centering
        \begin{tikzpicture}[automa]
            \node[state, initial] (q0) {$q_0$};
            \node[state, right=of q0] (q1) {$q_1$};
            \node[state, accepting, right=of q1] (q2) {$q_2$};
            \path (q0) edge node {$0$} (q1)
                (q1) edge[bend left] node {$1$} (q0)
                (q1) edge node {$0$} (q2)
                (q0) edge[loop above] node {$1$} ()
                (q2) edge[loop above] node {$0, 1$} ();

            % Matrice di transizione, a destra del diagramma
            \node[right=1.8cm of q2] {%
                \begin{tabular}{ccc}
                    \toprule
                    $\delta$ & $0$ & $1$ \\
                    \midrule
                    $q_0$ & $q_1$ & $q_0$ \\
                    $q_1$ & $q_2$ & $q_0$ \\
                    $q_2$ & $q_2$ & $q_2$ \\
                    \bottomrule
                \end{tabular}};
        \end{tikzpicture}
        \caption{DFA che riconosce le stringhe con almeno due $0$ consecutivi e la sua
            matrice di transizione.}
        \label{fig:es-due-zeri}
    \end{figure}

    Si scorpora la tesi secondo il significato degli stati, ottenendo
    \begin{equation*}
        \Psi(x) =
        \begin{cases}
            x \notin \mathcal{L} \land x = w0 \implies \hat{\delta}(q_0, x) = q_1,
                & \text{(1)} \\
            x \in \mathcal{L} \implies \hat{\delta}(q_0, x) = q_2,
                & \text{(2)} \\
            x \notin \mathcal{L} \land x = w1 \implies \hat{\delta}(q_0, x) = q_0,
                & \text{(3)}
        \end{cases}
    \end{equation*}
    con $w \in \Sigma^*$. Le tre condizioni si escludono a vicenda e coprono tutte le
    stringhe tranne $\varepsilon$, che non ha un ultimo simbolo e si verifica a parte.

    La prima lunghezza per cui esistono sia stringhe in $\mathcal{L}$ sia stringhe non in
    $\mathcal{L}$ \`e $|x| = 2$, quindi la base comprende tutte le stringhe con
    $|x| \leq 2$:
    \begin{equation*}
        \left.
        \begin{alignedat}{3}
            |x| = 0 \quad & x = \varepsilon \notin \mathcal{L} & \qquad
                & \hat{\delta}(q_0, \varepsilon) = q_0 \notin \mathcal{F} & & \\
            |x| = 1 \quad & x = 0 \notin \mathcal{L} &
                & \hat{\delta}(q_0, 0) = q_1
                &\quad & \text{(caso 1 di $\Psi$)} \\
            & x = 1 \notin \mathcal{L} &
                & \hat{\delta}(q_0, 1) = q_0
                && \text{(caso 3 di $\Psi$)} \\
            |x| = 2 \quad & x = 00 \in \mathcal{L} &
                & \hat{\delta}(q_0, 00) = q_2
                && \text{(caso 2 di $\Psi$)} \\
            & x = 01 \notin \mathcal{L} &
                & \hat{\delta}(q_0, 01) = q_0
                && \text{(caso 3 di $\Psi$)} \\
            & x = 10 \notin \mathcal{L} &
                & \hat{\delta}(q_0, 10) = q_1
                && \text{(caso 1 di $\Psi$)} \\
            & x = 11 \notin \mathcal{L} &
                & \hat{\delta}(q_0, 11) = q_0
                && \text{(caso 3 di $\Psi$)}
        \end{alignedat}
        \right\}
        \quad \text{\textsc{Base}}
    \end{equation*}

    \textsc{Passo induttivo}. Si suppone che $\Psi(x)$ valga per ogni $x \in \Sigma^*$ con
    $|x| \leq n$, per un certo $n \geq 2$, e sia $x$ una stringa con $|x| = n + 1$. Allora
    $x = x'0$ oppure $x = x'1$, con $|x'| = n \geq 2$: per ipotesi induttiva vale
    $\Psi(x')$, e $x'$ soddisfa esattamente una delle tre condizioni.
    \begin{enumerate}
        \item $x' \notin \mathcal{L}$ e $x' = w0$: per ipotesi induttiva
            $\hat{\delta}(q_0, x') = q_1$.
            \begin{enumerate}[label=\arabic{enumi}\alph*.]
                \item $x = x'0 = w00$ contiene due $0$ consecutivi, quindi
                    $x \in \mathcal{L}$ (caso 2 di $\Psi$):
                    \begin{equation*}
                        \hat{\delta}(q_0, x'0)
                        \overset{\text{def}}{=} \delta(\hat{\delta}(q_0, x'), 0)
                        \overset{\text{Hp}}{=} \delta(q_1, 0) = q_2.
                    \end{equation*}
                \item $x = x'1 = w01$: aggiungere un $1$ non crea due $0$ consecutivi,
                    quindi $x \notin \mathcal{L}$ e $x$ termina con $1$ (caso 3 di $\Psi$):
                    \begin{equation*}
                        \hat{\delta}(q_0, x'1)
                        \overset{\text{def}}{=} \delta(\hat{\delta}(q_0, x'), 1)
                        \overset{\text{Hp}}{=} \delta(q_1, 1) = q_0.
                    \end{equation*}
            \end{enumerate}
        \item $x' \in \mathcal{L}$: per ipotesi induttiva $\hat{\delta}(q_0, x') = q_2$.
            Qualunque sia l'ultimo simbolo $a \in \{0, 1\}$, $x = x'a$ contiene ancora i due
            $0$ consecutivi di $x'$, quindi $x \in \mathcal{L}$ (caso 2 di $\Psi$):
            \begin{equation*}
                \hat{\delta}(q_0, x'a)
                \overset{\text{def}}{=} \delta(\hat{\delta}(q_0, x'), a)
                \overset{\text{Hp}}{=} \delta(q_2, a) = q_2.
            \end{equation*}
        \item $x' \notin \mathcal{L}$ e $x' = w1$: per ipotesi induttiva
            $\hat{\delta}(q_0, x') = q_0$.
            \begin{enumerate}[label=\arabic{enumi}\alph*.]
                \item $x = x'0 = w10$: poich\'e $x'$ termina con $1$, aggiungere uno $0$ non
                    crea due $0$ consecutivi, quindi $x \notin \mathcal{L}$ e $x$ termina con
                    $0$ (caso 1 di $\Psi$):
                    \begin{equation*}
                        \hat{\delta}(q_0, x'0)
                        \overset{\text{def}}{=} \delta(\hat{\delta}(q_0, x'), 0)
                        \overset{\text{Hp}}{=} \delta(q_0, 0) = q_1.
                    \end{equation*}
                \item $x = x'1 = w11$: aggiungere un $1$ non crea due $0$ consecutivi,
                    quindi $x \notin \mathcal{L}$ e $x$ termina con $1$ (caso 3 di $\Psi$):
                    \begin{equation*}
                        \hat{\delta}(q_0, x'1)
                        \overset{\text{def}}{=} \delta(\hat{\delta}(q_0, x'), 1)
                        \overset{\text{Hp}}{=} \delta(q_0, 1) = q_0.
                    \end{equation*}
            \end{enumerate}
    \end{enumerate}
    In tutti i casi lo stato raggiunto da $x$ \`e quello previsto da $\Psi(x)$. Quindi
    $\hat{\delta}(q_0, x) = q_2 \in \mathcal{F}$ se $x \in \mathcal{L}$, mentre se
    $x \notin \mathcal{L}$ lo stato raggiunto \`e $q_0$ o $q_1$, che non sono finali:
    $\mathcal{L}(\mathcal{M}) = \mathcal{L}$.

    \esercizio
    \textit{Si dimostri che il linguaggio
    $\mathcal{L} = \{x \in \Sigma^* \mid \text{``$x$ non contiene due $1$
    consecutivi''}\}$ sull'alfabeto $\Sigma = \{0, 1\}$ \`e regolare.}

    \textsc{Svolgimento}. Lo riconosce il DFA
    $\mathcal{M} = \langle \{q_0, q_1, q_2\}, \{0, 1\}, \delta, q_0, \{q_0, q_1\} \rangle$
    della Figura~\ref{fig:es-no-uno}. Finch\'e non ha letto due $1$ consecutivi l'automa
    ricorda soltanto l'ultimo simbolo: \`e in $q_1$ se l'ultimo simbolo \`e $1$ e in $q_0$
    altrimenti. Appena legge il secondo $1$ consecutivo passa in $q_2$, da cui non esce
    pi\`u e che non \`e finale.

    \begin{figure}[H]
        \centering
        \begin{tikzpicture}[automa]
            \node[state, initial, accepting] (q0) {$q_0$};
            \node[state, accepting, right=of q0] (q1) {$q_1$};
            \node[state, right=of q1] (q2) {$q_2$};
            \path (q0) edge node {$1$} (q1)
                (q1) edge[bend left] node {$0$} (q0)
                (q1) edge node {$1$} (q2)
                (q0) edge[loop above] node {$0$} ()
                (q2) edge[loop above] node {$0, 1$} ();

            % Matrice di transizione, a destra del diagramma
            \node[right=1.8cm of q2] {%
                \begin{tabular}{ccc}
                    \toprule
                    $\delta$ & $0$ & $1$ \\
                    \midrule
                    $q_0$ & $q_0$ & $q_1$ \\
                    $q_1$ & $q_0$ & $q_2$ \\
                    $q_2$ & $q_2$ & $q_2$ \\
                    \bottomrule
                \end{tabular}};
        \end{tikzpicture}
        \caption{DFA che riconosce le stringhe senza due $1$ consecutivi e la sua matrice
            di transizione.}
        \label{fig:es-no-uno}
    \end{figure}

    Si scorpora la tesi secondo il significato degli stati, ottenendo
    \begin{equation*}
        \Psi(x) =
        \begin{cases}
            x \in \mathcal{L} \land x = w0 \implies \hat{\delta}(q_0, x) = q_0,
                & \text{(1)} \\
            x \in \mathcal{L} \land x = w1 \implies \hat{\delta}(q_0, x) = q_1,
                & \text{(2)} \\
            x \notin \mathcal{L} \implies \hat{\delta}(q_0, x) = q_2,
                & \text{(3)}
        \end{cases}
    \end{equation*}
    con $w \in \Sigma^*$. Anche qui le tre condizioni si escludono a vicenda e coprono
    tutte le stringhe tranne $\varepsilon$, che si verifica a parte.

    La prima lunghezza per cui esistono sia stringhe in $\mathcal{L}$ sia stringhe non in
    $\mathcal{L}$ \`e $|x| = 2$, quindi la base comprende tutte le stringhe con
    $|x| \leq 2$:
    \begin{equation*}
        \left.
        \begin{alignedat}{3}
            |x| = 0 \quad & x = \varepsilon \in \mathcal{L} & \qquad
                & \hat{\delta}(q_0, \varepsilon) = q_0 \in \mathcal{F} & & \\
            |x| = 1 \quad & x = 0 \in \mathcal{L} &
                & \hat{\delta}(q_0, 0) = q_0
                &\quad & \text{(caso 1 di $\Psi$)} \\
            & x = 1 \in \mathcal{L} &
                & \hat{\delta}(q_0, 1) = q_1
                && \text{(caso 2 di $\Psi$)} \\
            |x| = 2 \quad & x = 00 \in \mathcal{L} &
                & \hat{\delta}(q_0, 00) = q_0
                && \text{(caso 1 di $\Psi$)} \\
            & x = 01 \in \mathcal{L} &
                & \hat{\delta}(q_0, 01) = q_1
                && \text{(caso 2 di $\Psi$)} \\
            & x = 10 \in \mathcal{L} &
                & \hat{\delta}(q_0, 10) = q_0
                && \text{(caso 1 di $\Psi$)} \\
            & x = 11 \notin \mathcal{L} &
                & \hat{\delta}(q_0, 11) = q_2
                && \text{(caso 3 di $\Psi$)}
        \end{alignedat}
        \right\}
        \quad \text{\textsc{Base}}
    \end{equation*}

    \textsc{Passo induttivo}. Si suppone che $\Psi(x)$ valga per ogni $x \in \Sigma^*$ con
    $|x| \leq n$, per un certo $n \geq 2$, e sia $x$ una stringa con $|x| = n + 1$. Allora
    $x = x'0$ oppure $x = x'1$, con $|x'| = n \geq 2$: per ipotesi induttiva vale
    $\Psi(x')$, e $x'$ soddisfa esattamente una delle tre condizioni.
    \begin{enumerate}
        \item $x' \in \mathcal{L}$ e $x' = w0$: per ipotesi induttiva
            $\hat{\delta}(q_0, x') = q_0$.
            \begin{enumerate}[label=\arabic{enumi}\alph*.]
                \item $x = x'0 = w00$: aggiungere uno $0$ non crea due $1$ consecutivi,
                    quindi $x \in \mathcal{L}$ (caso 1 di $\Psi$):
                    \begin{equation*}
                        \hat{\delta}(q_0, x'0)
                        \overset{\text{def}}{=} \delta(\hat{\delta}(q_0, x'), 0)
                        \overset{\text{Hp}}{=} \delta(q_0, 0) = q_0.
                    \end{equation*}
                \item $x = x'1 = w01$: poich\'e $x'$ termina con $0$, aggiungere un $1$ non
                    crea due $1$ consecutivi, quindi $x \in \mathcal{L}$ (caso 2 di $\Psi$):
                    \begin{equation*}
                        \hat{\delta}(q_0, x'1)
                        \overset{\text{def}}{=} \delta(\hat{\delta}(q_0, x'), 1)
                        \overset{\text{Hp}}{=} \delta(q_0, 1) = q_1.
                    \end{equation*}
            \end{enumerate}
        \item $x' \in \mathcal{L}$ e $x' = w1$: per ipotesi induttiva
            $\hat{\delta}(q_0, x') = q_1$.
            \begin{enumerate}[label=\arabic{enumi}\alph*.]
                \item $x = x'0 = w10$: aggiungere uno $0$ non crea due $1$ consecutivi,
                    quindi $x \in \mathcal{L}$ (caso 1 di $\Psi$):
                    \begin{equation*}
                        \hat{\delta}(q_0, x'0)
                        \overset{\text{def}}{=} \delta(\hat{\delta}(q_0, x'), 0)
                        \overset{\text{Hp}}{=} \delta(q_1, 0) = q_0.
                    \end{equation*}
                \item $x = x'1 = w11$ termina con due $1$ consecutivi, quindi
                    $x \notin \mathcal{L}$ (caso 3 di $\Psi$):
                    \begin{equation*}
                        \hat{\delta}(q_0, x'1)
                        \overset{\text{def}}{=} \delta(\hat{\delta}(q_0, x'), 1)
                        \overset{\text{Hp}}{=} \delta(q_1, 1) = q_2.
                    \end{equation*}
            \end{enumerate}
        \item $x' \notin \mathcal{L}$: per ipotesi induttiva $\hat{\delta}(q_0, x') = q_2$.
            Qualunque sia l'ultimo simbolo $a \in \{0, 1\}$, $x = x'a$ contiene ancora i due
            $1$ consecutivi di $x'$, quindi $x \notin \mathcal{L}$ (caso 3 di $\Psi$):
            \begin{equation*}
                \hat{\delta}(q_0, x'a)
                \overset{\text{def}}{=} \delta(\hat{\delta}(q_0, x'), a)
                \overset{\text{Hp}}{=} \delta(q_2, a) = q_2.
            \end{equation*}
    \end{enumerate}
    In tutti i casi lo stato raggiunto da $x$ \`e quello previsto da $\Psi(x)$. Quindi
    se $x \in \mathcal{L}$ lo stato raggiunto \`e $q_0$ o $q_1$, entrambi finali, mentre se
    $x \notin \mathcal{L}$ \`e $q_2 \notin \mathcal{F}$: $\mathcal{L}(\mathcal{M}) =
    \mathcal{L}$.
\end{document}
